% Le fonctionnement de la mondialisation — Géographie Tle A — Cours signé OnBuch+
% Programme officiel MINESEC. Compiler avec : tectonic N12-fonctionnement-mondialisation.tex
\newcommand{\DOCMATIERE}{Géographie}
\newcommand{\DOCNIVEAU}{Tle A}
\newcommand{\DOCMODULE}{Module 2 : La libéralisation des échanges (la mondialisation)}
\newcommand{\DOCLECON}{N12}
\newcommand{\DOCTITRE}{Le fonctionnement de la mondialisation}
% OnBuch+ — préambule « cours signés » ENRICHI (compilé avec tectonic / XeLaTeX)
% Système visuel « Pop » d'OnBuch+ : fond crème (jamais blanc), bordures encre
% épaisses, ombres « dures » décalées, titres Archivo Black en capitales.
% Version MATHÉMATIQUES : graphes (pgfplots/tikz), boîtes pédagogiques
% (définition, propriété, méthode, attention, le savais-tu, exercice résolu,
% exercice, TP, bilan), page de garde et signature OnBuch+ ancrée en bas.
%
% Macros attendues dans main.tex AVANT \input{preamble} :
%   \DOCMATIERE  \DOCNIVEAU  \DOCMODULE  \DOCLECON (numéro)  \DOCTITRE
\documentclass[11pt]{article}

\usepackage{fontspec}
\usepackage[french]{babel}
\usepackage[a4paper,margin=2cm,top=3.4cm,bottom=2.8cm,headheight=1.4cm,footskip=1.3cm]{geometry}
\usepackage{amsmath,amssymb,amsfonts}
\usepackage{mathtools} % \xmapsto, \coloneqq... souvent utilisés par les modèles
\usepackage{cancel} % simplifications barrées (bilans de forces)
% \abs{z} (module, valeur absolue) et \norm{u} (norme), utilisés par les modèles
\providecommand{\abs}{}\let\abs\relax\DeclarePairedDelimiter{\abs}{\lvert}{\rvert}
\providecommand{\norm}{}\let\norm\relax\DeclarePairedDelimiter{\norm}{\lVert}{\rVert}
\usepackage[table]{xcolor} % [table] : fournit aussi \rowcolors (alternance de lignes)
\usepackage{pagecolor}
\usepackage{tikz}
\usetikzlibrary{arrows.meta,positioning,calc,shapes.geometric,shapes.misc,shapes.arrows,decorations.pathmorphing,decorations.pathreplacing,decorations.markings,patterns,shadows,mindmap,trees,fit,backgrounds,matrix,intersections,angles,quotes}
\usepackage{pgfplots}
\usepackage[version=4]{mhchem} % équations nucléaires \ce{^{235}_{92}U}
\usepackage[european,siunitx]{circuitikz} % schémas électriques
\pgfplotsset{compat=1.18}
\usepgfplotslibrary{fillbetween,groupplots} % aires entre courbes (calcul intégral)
\usepackage{enumitem}
\usepackage{fancyhdr}
% [most] charge toutes les bibliothèques tcolorbox (listings, minted, raster,
% documentation...) et gonfle inutilement la mémoire TeX sur un document long
% (~300 boîtes) : on ne charge que ce qui est réellement utilisé.
\usepackage[skins,breakable]{tcolorbox}
\usepackage{graphicx}
\graphicspath{{/home/user/t-t/geo-onbuch/images/}}
\usepackage{colortbl}
\usepackage{booktabs}
\usepackage{tabularx}
\usepackage{multirow}
\usepackage{diagbox} % case d'en-tête diagonale des tableaux à double entrée
% Colonnes extensibles centrée / gauche / droite (C, L, R), souvent utilisées par les modèles
\newcolumntype{C}{>{\centering\arraybackslash}X}
\newcolumntype{L}{>{\raggedright\arraybackslash}X}
\newcolumntype{R}{>{\raggedleft\arraybackslash}X}
\usepackage{array}
\usepackage{multicol}
\usepackage{siunitx}
% parse-numbers=false : accepte aussi \SI{2,5 \times 10^{-3}}{...}
\sisetup{locale=FR,output-decimal-marker={,},group-minimum-digits=4,parse-numbers=false}
\DeclareSIUnit\ohm{\ensuremath{\symup{\Omega}}} % U+2126 absent de la police
\DeclareSIUnit\division{div} % réglages d'oscilloscope (ms/div, V/div)
\DeclareSIUnit\tr{tr} % tours (vitesse de rotation, ex. \SI{1500}{\tr\per\minute})
\DeclareSIUnit\molar{M} % concentration molaire (ex. \SI{3}{\molar}, \SI{1}{\micro\molar})
\DeclareSIUnit\an{an} % année (le modèle confond parfois avec \year, un compteur TeX) — ex. \SI{5}{\kilo\meter\per\an}
\DeclareSIUnit\FCFA{FCFA} % franc CFA (ex. \SI{500}{\FCFA\per\kilo\gram})
\DeclareSIUnit\degreeBrix{\textdegree Brix} % teneur en sucre d'un sirop/jus (réfractométrie)
\DeclareSIUnit\month{mois} % le modèle utilise parfois \month, un compteur TeX, comme unité de durée
\DeclareSIUnit\microliter{\micro\liter} % le modèle écrit parfois \microliter en un seul token
\DeclareSIUnit\million{millions} % dénombrement (ex. \SI{15}{\million\per\mL} spermatozoïdes)
\usepackage{qrcode}
% \degree : commande fréquemment utilisée par les modèles pour un angle ou une
% température en dehors de \SI{}{} (non fournie par les paquets chargés).
\providecommand{\degree}{^{\circ}}
% \qed (fin de démonstration), utilisé par les modèles sans amsthm
\providecommand{\qed}{\hfill$\square$}
% \barre : double barre des valeurs interdites dans un tableau de variations (tkz-tab)
\providecommand{\barre}{\ensuremath{\|}}
% environnement proof (amsthm non chargé) utilisé par les modèles
\ifdefined\proof\else\newenvironment{proof}[1][Démonstration]{\par\noindent\textit{#1.}\ }{\qed\par}\fi
\usepackage{lastpage}
\usepackage{hyperref}
\hypersetup{hidelinks}

% ============================================================
% Palette « Pop » OnBuch+ — mêmes hex que la classe P (plus_ui.dart)
% ============================================================
\definecolor{poporange}{HTML}{F59321}
\definecolor{popdark}{HTML}{E07A0C}
\definecolor{popcream}{HTML}{FFF6E8}
\definecolor{popink}{HTML}{1C1714}
\definecolor{poppurple}{HTML}{7A5AE0}
\definecolor{popgreen}{HTML}{2E9E6A}
\definecolor{poppink}{HTML}{E0457B}
\definecolor{popblue}{HTML}{2F7DE1}
\definecolor{popgold}{HTML}{C98A12}
\definecolor{popmuted}{HTML}{8A7F76}
\definecolor{popline}{HTML}{EADFD2}
\definecolor{popgray}{HTML}{8A7F76}
\colorlet{popgrey}{popgray}
% Alias tolérants (le modèle nomme parfois les couleurs différemment)
\colorlet{popwhite}{white}
\colorlet{popblack}{popink}
\colorlet{popred}{poppink}
\colorlet{popyellow}{popgold}
% Teintes claires (fonds de tableaux/encadrés) — le modèle nomme parfois
% les couleurs avec un suffixe L (ex. popblueL) sans les avoir définies.
\colorlet{poporangeL}{poporange!20}
\colorlet{popdarkL}{popdark!20}
\colorlet{poppurpleL}{poppurple!20}
\colorlet{popgreenL}{popgreen!20}
\colorlet{poppinkL}{poppink!20}
\colorlet{popblueL}{popblue!20}
\colorlet{popgoldL}{popgold!20}
\colorlet{popredL}{popred!20}
\colorlet{popgrayL}{popgray!20}
% Teintes claires pour fonds de tableaux / zones
\colorlet{popblueL}{popblue!10!white}
\colorlet{poporangeL}{poporange!14!white}
\colorlet{popgreenL}{popgreen!12!white}
\colorlet{poppurpleL}{poppurple!12!white}
\colorlet{poppinkL}{poppink!12!white}
\colorlet{popgoldL}{popgold!14!white}

\pagecolor{popcream}

% ============================================================
% Typographie
% ============================================================
\defaultfontfeatures{Ligatures=TeX}
\setmainfont{PlusJakartaSans-Medium.ttf}[Path=../../../fonts/,BoldFont=PlusJakartaSans-Bold.ttf]
% Maths en sans-serif (Fira Math) avec chiffres et lettres droites de Plus
% Jakarta, pour que les formules se fondent dans le texte.
\usepackage[math-style=ISO,bold-style=ISO]{unicode-math}
\setmathfont{FiraMath-Regular.otf}[Path=../../../fonts/]
\setmathfont{PlusJakartaSans-Medium.ttf}[Path=../../../fonts/,range={up/num,up/{Latin,latin},\mathup}]
\usepackage{icomma} % virgule décimale sans espace en mode math
% Caractères mathématiques Unicode absents des polices de texte (Plus Jakarta,
% Archivo Black) : rendus par leur équivalent mathématique, en texte comme en
% formule — sinon ils s'affichent en carré vide (ex. « Arithmétique dans ℤ »).
\usepackage{newunicodechar}
\newunicodechar{ℝ}{\ensuremath{\mathbb{R}}}
\newunicodechar{ℤ}{\ensuremath{\mathbb{Z}}}
\newunicodechar{ℂ}{\ensuremath{\mathbb{C}}}
\newunicodechar{ℕ}{\ensuremath{\mathbb{N}}}
\newunicodechar{ℚ}{\ensuremath{\mathbb{Q}}}
\newunicodechar{𝔻}{\ensuremath{\mathbb{D}}}
\newunicodechar{α}{\ensuremath{\alpha}}
\newunicodechar{β}{\ensuremath{\beta}}
\newunicodechar{γ}{\ensuremath{\gamma}}
\newunicodechar{δ}{\ensuremath{\delta}}
\newunicodechar{ε}{\ensuremath{\varepsilon}}
\newunicodechar{θ}{\ensuremath{\theta}}
\newunicodechar{λ}{\ensuremath{\lambda}}
\newunicodechar{μ}{\ensuremath{\mu}}
\newunicodechar{σ}{\ensuremath{\sigma}}
\newunicodechar{τ}{\ensuremath{\tau}}
\newunicodechar{φ}{\ensuremath{\varphi}}
\newunicodechar{ψ}{\ensuremath{\psi}}
\newunicodechar{ω}{\ensuremath{\omega}}
\newunicodechar{Σ}{\ensuremath{\Sigma}}
% majuscules produites par \MakeUppercase dans les titres (α-aminés -> Α)
\newunicodechar{Α}{\ensuremath{\alpha}}
\newunicodechar{Β}{\ensuremath{\beta}}
\newunicodechar{Γ}{\ensuremath{\Gamma}}
\newunicodechar{Δ}{\ensuremath{\Delta}}
\newunicodechar{Ε}{\ensuremath{\varepsilon}}
\newunicodechar{Θ}{\ensuremath{\Theta}}
\newunicodechar{Λ}{\ensuremath{\Lambda}}
\newunicodechar{Μ}{\ensuremath{\mu}}
\newunicodechar{Τ}{\ensuremath{\tau}}
\newunicodechar{Φ}{\ensuremath{\Phi}}
\newunicodechar{Ψ}{\ensuremath{\Psi}}
\newunicodechar{Ω}{\ensuremath{\Omega}}
\newunicodechar{↦}{\ensuremath{\mapsto}}
\newunicodechar{∈}{\ensuremath{\in}}
\newunicodechar{∉}{\ensuremath{\notin}}
\newunicodechar{∧}{\ensuremath{\wedge}}
\newunicodechar{⇔}{\ensuremath{\Leftrightarrow}}
\newunicodechar{⟺}{\ensuremath{\Longleftrightarrow}}
\newunicodechar{⇒}{\ensuremath{\Rightarrow}}
\newunicodechar{⟹}{\ensuremath{\Longrightarrow}}
\newunicodechar{∘}{\ensuremath{\circ}}
\newunicodechar{∪}{\ensuremath{\cup}}
\newunicodechar{∩}{\ensuremath{\cap}}
\newunicodechar{≡}{\ensuremath{\equiv}}
\newunicodechar{⊂}{\ensuremath{\subset}}
\newunicodechar{⊥}{\ensuremath{\perp}}
\newunicodechar{∃}{\ensuremath{\exists}}
\newunicodechar{∀}{\ensuremath{\forall}}
\newunicodechar{∛}{\ensuremath{\sqrt[3]{\,}}}
\newunicodechar{∜}{\ensuremath{\sqrt[4]{\,}}}
\newunicodechar{⃗}{\ensuremath{{}^{\to}}}
\newunicodechar{ⁿ}{\ifmmode^{n}\else\textsuperscript{\ensuremath{n}}\fi}
\newunicodechar{ˣ}{\ifmmode^{x}\else\textsuperscript{\ensuremath{x}}\fi}
\newunicodechar{ᵇ}{\ifmmode^{b}\else\textsuperscript{\ensuremath{b}}\fi}
\newunicodechar{ʳ}{\ifmmode^{r}\else\textsuperscript{\ensuremath{r}}\fi}
\newunicodechar{ᵘ}{\ifmmode^{u}\else\textsuperscript{\ensuremath{u}}\fi}
\newunicodechar{ʸ}{\ifmmode^{y}\else\textsuperscript{\ensuremath{y}}\fi}
\newunicodechar{ᵃ}{\ifmmode^{a}\else\textsuperscript{\ensuremath{a}}\fi}
\newunicodechar{ᵉ}{\ifmmode^{e}\else\textsuperscript{\ensuremath{e}}\fi}
\newunicodechar{ᵏ}{\ifmmode^{k}\else\textsuperscript{\ensuremath{k}}\fi}
\newunicodechar{ᵖ}{\ifmmode^{p}\else\textsuperscript{\ensuremath{p}}\fi}
\newunicodechar{ᴺ}{\ifmmode^{N}\else\textsuperscript{\ensuremath{N}}\fi}
\newunicodechar{ᶜ}{\ifmmode^{c}\else\textsuperscript{\ensuremath{c}}\fi}
\newunicodechar{ʰ}{\ifmmode^{h}\else\textsuperscript{\ensuremath{h}}\fi}
\newunicodechar{ᵠ}{\ifmmode^{q}\else\textsuperscript{\ensuremath{q}}\fi}
\newunicodechar{ₙ}{\ifmmode_{n}\else\textsubscript{\ensuremath{n}}\fi}
\newunicodechar{ₐ}{\ifmmode_{a}\else\textsubscript{\ensuremath{a}}\fi}
\newunicodechar{ᵢ}{\ifmmode_{i}\else\textsubscript{\ensuremath{i}}\fi}
\newunicodechar{ᵣ}{\ifmmode_{r}\else\textsubscript{\ensuremath{r}}\fi}
\newunicodechar{ₖ}{\ifmmode_{k}\else\textsubscript{\ensuremath{k}}\fi}
\newunicodechar{ᵦ}{\ifmmode_{\beta}\else\textsubscript{\ensuremath{\beta}}\fi}
\newunicodechar{ₓ}{\ifmmode_{x}\else\textsubscript{\ensuremath{x}}\fi}
\newfontfamily\popdisplay{ArchivoBlack-Regular.ttf}[Path=../../../fonts/]
\newfontfamily\poplogo{SpaceGrotesk-Bold.ttf}[Path=../../../fonts/]
\newfontfamily\popsemi{PlusJakartaSans-SemiBold.ttf}[Path=../../../fonts/]
\newfontfamily\popxbold{PlusJakartaSans-ExtraBold.ttf}[Path=../../../fonts/]
\newfontfamily\popxboldls{PlusJakartaSans-ExtraBold.ttf}[Path=../../../fonts/,LetterSpace=3.0]

\setlist[enumerate]{leftmargin=1.7em,itemsep=5pt,topsep=4pt}
\setlist[itemize]{leftmargin=1.7em,itemsep=4pt,topsep=4pt,label={\color{poporange}\textbullet}}
\setlength{\parindent}{0pt}
\setlength{\parskip}{5pt}
\renewcommand{\arraystretch}{1.25}

% pgfplots : style Pop par défaut
\pgfplotsset{
  every axis/.append style={axis line style={popink,line width=1pt},tick style={popink},
    grid style={popline,line width=0.5pt},label style={font=\small},tick label style={font=\footnotesize}},
  popcurve/.style={poporange,line width=1.8pt,no markers,smooth},
}
% Même style disponible dans un \draw TikZ simple (hors pgfplots)
\tikzset{popcurve/.style={poporange,line width=1.8pt}}
\tikzset{popgrid/.style={popline,very thin}}
\colorlet{popcurve}{poporange} % le nom du style est parfois utilisé comme couleur

% ============================================================
% Pill / squiggle
% ============================================================
\newcommand{\poppill}[3]{%
  \tikz[baseline=-0.55ex]\node[draw=popink,line width=1.2pt,rounded corners=2.6pt,%
    fill=#2,inner xsep=6.5pt,inner ysep=3pt]{%
    \popxboldls\fontsize{7}{7}\selectfont\color{#3}\MakeUppercase{#1}};%
}
\newcommand{\popsquiggle}[1]{%
  \tikz[baseline=-0.4ex]{
    \draw[#1,line width=2.6pt,line cap=round]
      plot [smooth] coordinates {(0,0.09) (0.34,0.01) (0.68,0.21) (1.02,0.01) (1.36,0.10)};
  }%
}
% Mot-clé surligné (marqueur orange sous le texte)
\newcommand{\cle}[1]{{\popxbold\color{popdark}#1}}

% ============================================================
% En-tête / pied
% ============================================================
\pagestyle{fancy}
\fancyhf{}
\renewcommand{\headrulewidth}{0pt}
\renewcommand{\footrulewidth}{0pt}
\lhead{\raisebox{-0.15ex}{\poplogo\fontsize{13}{13}\selectfont\color{popink}OnBuch\color{poporange}+}}
\rhead{\poppill{\DOCMATIERE\ \textbullet\ \DOCNIVEAU\ \textbullet\ Leçon \DOCLECON}{white}{popink}}
\lfoot{\popxbold\fontsize{7.6}{7.6}\selectfont\color{popmuted}\MakeUppercase{Cours signé OnBuch+}\ \textbullet\ {\popsemi\DOCTITRE}}
\rfoot{\tikz[baseline=-0.6ex]\node[rounded corners=8.5pt,fill=popink,minimum height=17pt,minimum width=17pt,inner xsep=5pt,inner sep=0]{\popxbold\fontsize{8}{8}\selectfont\color{popcream}\thepage\,/\,\pageref*{LastPage}};}

% ============================================================
% Titres
% ============================================================
\newcommand{\chaptitre}[2]{%
  \begin{tcolorbox}[enhanced,colback=poporange,colframe=popink,boxrule=2.6pt,arc=11pt,
      drop shadow=popink,left=15pt,right=15pt,top=12pt,bottom=11pt,before skip=3pt,after skip=18pt]
    \poppill{#2}{popink}{popcream}\par\vspace{9pt}
    {\raggedright\hyphenpenalty=10000\exhyphenpenalty=10000\popdisplay\fontsize{22}{24}\selectfont\color{popcream}\MakeUppercase{#1}\par}
  \end{tcolorbox}
}
% Section numérotée : pastille orange avec le numéro + titre Archivo
\newcounter{popsec}
\newcommand{\coursec}[1]{%
  \stepcounter{popsec}\setcounter{popsub}{0}%
  \par\vspace{16pt}\noindent
  \begin{minipage}{\linewidth}
  \tikz[baseline=-0.7ex]\node[draw=popink,line width=1.6pt,fill=poporange,rounded corners=4pt,minimum size=22pt,inner sep=0]{\popdisplay\fontsize{12}{12}\selectfont\color{popcream}\thepopsec};%
  \hspace{8pt}{\popdisplay\fontsize{15}{17}\selectfont\color{popink}\MakeUppercase{#1}}\par
  \vspace{4pt}\noindent{\color{poporange}\rule{36pt}{3.6pt}}
  \end{minipage}\par\nopagebreak\vspace{8pt}
}
\newcounter{popsub}
\newcommand{\courssub}[1]{%
  \stepcounter{popsub}%
  \par\vspace{9pt}\noindent{\popxbold\fontsize{11.5}{13}\selectfont\color{popink}{\color{poporange}\thepopsec.\thepopsub}\hspace{6pt}#1}\par\nopagebreak\vspace{4pt}
}

% ============================================================
% Boîtes pédagogiques — même recette : fond blanc, bordure épaisse colorée,
% ombre dure encre, badge plein. Titre optionnel [#1].
% ============================================================
\tcbset{popbase/.style={breakable,enhanced,colback=white,boxrule=2.3pt,arc=9pt,
  shadow={3pt}{-3pt}{0pt}{popink},left=14pt,right=14pt,top=11pt,bottom=11pt,before skip=11pt,after skip=15pt,
  fontupper=\popsemi\fontsize{10.6}{14.5}\selectfont\color{popink},
  fontlower=\popsemi\fontsize{10.6}{14.5}\selectfont\color{popink}}}
% Coche dessinée (le glyphe ✓ n'existe pas dans les polices du document)
\newcommand{\popcheck}[1]{\tikz[baseline=-0.5ex]\draw[#1,line width=1.6pt,line cap=round,line join=round](-0.13,0.0)--(-0.04,-0.09)--(0.13,0.1);}
\providecommand{\coche}{\popcheck{popgreen}}
\providecommand{\resultat}[1]{\par\smallskip\noindent\fcolorbox{popgreen}{popgreenL}{\parbox{\dimexpr\linewidth-2\fboxsep-2\fboxrule\relax}{\textbf{Résultat.} #1}}\par\smallskip}
\newcommand{\popboxhead}[3]{\poppill{#1}{#2}{white}\ifx\relax#3\relax\else\hspace{6pt}{\popxbold\fontsize{10.6}{12}\selectfont\color{popink}#3}\fi\par\vspace{7pt}}

\newtcolorbox{aretenir}[1][]{popbase,colframe=popgreen,before upper={\popboxhead{À retenir}{popgreen}{#1}}}
\newtcolorbox{exemplebox}[1][]{popbase,colframe=poporange,before upper={\popboxhead{Exemple}{poporange}{#1}}}
\newtcolorbox{definition}[1][]{popbase,colframe=popblue,before upper={\popboxhead{Définition}{popblue}{#1}}}
\newtcolorbox{propriete}[1][]{popbase,colframe=poppurple,before upper={\popboxhead{Propriété}{poppurple}{#1}}}
\newtcolorbox{methode}[1][]{popbase,colframe=popgold,before upper={\popboxhead{Méthode}{popgold}{#1}}}
\newtcolorbox{attention}[1][]{popbase,colframe=poppink,before upper={\popboxhead{Attention}{poppink}{#1}}}
\newtcolorbox{experience}[1][]{popbase,colframe=popblue,colback=popblueL,before upper={\popboxhead{Expérience}{popblue}{#1}}}
% « Le savais-tu ? » : carte violette pleine (contexte, culture, Cameroun)
\newtcolorbox{savaistu}[1][]{popbase,colback=poppurple,colframe=popink,
  fontupper=\popsemi\fontsize{10.4}{14.2}\selectfont\color{white},
  before upper={\poppill{Le savais-tu\,?}{white}{poppurple}\ifx\relax#1\relax\else\hspace{6pt}{\popxbold\color{white}#1}\fi\par\vspace{7pt}}}
% Exercice résolu : énoncé en haut, \tcblower puis la solution sur fond crème
\newcounter{popexo}\newcounter{popexr}
\newtcolorbox{exoresolu}[1][]{popbase,colframe=poporange,colbacklower=popcream,
  segmentation style={popink,line width=1.2pt,dashed},
  before upper={\stepcounter{popexr}\popboxhead{Exercice résolu \thepopexr}{poporange}{#1}},
  before lower={\poppill{Solution}{popink}{popcream}\par\vspace{6pt}}}
% Exercice (non résolu en place) — la correction est dans la partie Corrigés
\newtcolorbox{exercice}[1][]{popbase,colframe=popink,
  before upper={\stepcounter{popexo}\popboxhead{Exercice \thepopexo}{popink}{#1}}}
% Corrigé
\newtcolorbox{corrige}[1][]{popbase,colframe=popgreen,colback=popgreenL,
  before upper={\popboxhead{Corrigé}{popgreen}{#1}}}
% Objectifs / prérequis / bilan
\newtcolorbox{objectifs}{popbase,colframe=popink,colback=poporangeL,
  before upper={\popboxhead{Objectifs de la leçon}{popink}{}}}
\newtcolorbox{prerequis}{popbase,colframe=popink,colback=popblueL,
  before upper={\popboxhead{Prérequis}{popblue}{}}}
\newtcolorbox{bilan}{popbase,colframe=popink,colback=white,boxrule=2.8pt,
  before upper={\popboxhead{Fiche bilan}{poporange}{L'essentiel à connaître pour l'examen}}}
% Figure encadrée avec légende
\newtcolorbox{popfigure}[1][]{enhanced,breakable=false,colback=white,colframe=popline,boxrule=1.4pt,arc=8pt,
  left=10pt,right=10pt,top=10pt,bottom=8pt,before skip=10pt,after skip=12pt,halign=center,
  lower separated=false,halign lower=center,
  fontlower=\popsemi\fontsize{9}{11.5}\selectfont\color{popmuted},
  #1}
\newcommand{\legende}[1]{\par\vspace{6pt}{\popsemi\fontsize{9}{11.5}\selectfont\color{popmuted}#1}}

% ============================================================
% Signature OnBuch+
% ============================================================
\newcommand{\poplogoinline}[1]{{\poplogo\fontsize{#1}{#1}\selectfont\color{popink}OnBuch\color{poporange}+}}
\newcommand{\signatureonbuch}{%
  \begin{tcolorbox}[enhanced,colback=popink,colframe=popink,boxrule=2.2pt,arc=10pt,
      drop shadow=poporange,left=14pt,right=14pt,top=10pt,bottom=10pt,before skip=0pt,after skip=0pt]
    \tikz[baseline=-0.7ex]\node[circle,fill=popgreen,draw=popcream,line width=1.3pt,minimum size=22pt,inner sep=0]{\popcheck{white}};%
    \hspace{9pt}\begin{minipage}[c]{0.62\linewidth}
      {\popxbold\fontsize{12}{13}\selectfont\color{popcream}Signé\ \textbullet\ {\poplogo OnBuch\color{poporange}+}}\par\vspace{2pt}
      {\raggedright\popsemi\fontsize{8}{10.5}\selectfont\color{popline}\MakeUppercase{Équipe pédagogique OnBuch+}\\\MakeUppercase{Conforme au programme officiel MINESEC}\par}
    \end{minipage}\hfill
    {\poplogo\fontsize{22}{22}\selectfont\color{popcream}OnBuch\color{poporange}+}
  \end{tcolorbox}%
}

% ============================================================
% Page de garde
% #1 titre · #2 matière · #3 niveau · #4 module · #5 n° leçon · #6 durée ·
% #7 description / objectifs (texte ou itemize)
% ============================================================
\newcommand{\pagegarde}[7]{%
  \thispagestyle{empty}
  \newgeometry{a4paper,margin=1.7cm,top=1.6cm,bottom=1.4cm}%
  \begin{tikzpicture}[remember picture,overlay]
    \fill[poporange!18] ([xshift=-3.2cm,yshift=-3.6cm]current page.north east) circle (4.6cm);
    \fill[poppurple!14] ([xshift=2.4cm,yshift=7.6cm]current page.south west) circle (3.8cm);
  \end{tikzpicture}%
  {\poplogo\fontsize{30}{30}\selectfont\color{popink}OnBuch\color{poporange}+}\hfill
  \raisebox{0.6ex}{\poppill{Cours signé \textbullet\ Premium}{popink}{popcream}}\par
  \vspace{1pt}\popsquiggle{poppurple}\par
  \vspace{8pt}
  \begin{tcolorbox}[enhanced,colback=poporange,colframe=popink,boxrule=2.8pt,arc=13pt,
      drop shadow=popink,left=18pt,right=18pt,top=12pt,bottom=13pt,after skip=10pt]
    \poppill{#2 \textbullet\ #3}{popink}{popcream}\hspace{5pt}\poppill{#4}{white}{popink}\par\vspace{8pt}
    {\popdisplay\fontsize{14}{14}\selectfont\color{popink}LEÇON #5}\par\vspace{4pt}
    {\raggedright\hyphenpenalty=10000\exhyphenpenalty=10000\popdisplay\fontsize{25}{27}\selectfont\color{popcream}\MakeUppercase{#1}\par}
    \vspace{8pt}
    {\popxbold\fontsize{9.5}{9.5}\selectfont\color{popink}\MakeUppercase{Durée conseillée : #6}}
  \end{tcolorbox}
  \begin{tcolorbox}[enhanced,colback=white,colframe=popink,boxrule=2.2pt,arc=10pt,
      drop shadow=popink,left=16pt,right=16pt,top=10pt,bottom=10pt,after skip=8pt]
    \poppill{Dans cette leçon}{poppurple}{white}\par\vspace{5pt}
    \popsemi\fontsize{9.3}{12.6}\selectfont\color{popink}#7
  \end{tcolorbox}
  \noindent\begin{minipage}[t]{0.487\linewidth}
    \begin{tcolorbox}[enhanced,colback=popgreen,colframe=popink,boxrule=2.2pt,arc=10pt,
        drop shadow=popink,left=11pt,right=11pt,top=10pt,bottom=10pt,height=3.1cm,valign=center]
      \tcbox[colback=white,colframe=popink,boxrule=1.3pt,arc=5pt,boxsep=0pt,nobeforeafter,
        left=3pt,right=3pt,top=3pt,bottom=3pt,valign=center]{\qrcode[height=1.9cm]{https://play.google.com/store/apps/details?id=com.luvvix.onbuch&hl=fr}}%
      \hspace{8pt}\begin{minipage}[c]{0.5\linewidth}
        {\popdisplay\fontsize{11}{12.5}\selectfont\color{white}\MakeUppercase{Télécharge OnBuch}}\par\vspace{4pt}
        {\raggedright\popsemi\fontsize{8.4}{11}\selectfont\color{white}Gratuit \textbullet\ Google Play\\Tuteur IA, annales, concours, résultats\par}
      \end{minipage}
    \end{tcolorbox}
  \end{minipage}\hfill
  \begin{minipage}[t]{0.487\linewidth}
    \begin{tcolorbox}[enhanced,colback=poppurple,colframe=popink,boxrule=2.2pt,arc=10pt,
        drop shadow=popink,left=11pt,right=11pt,top=10pt,bottom=10pt,height=3.1cm,valign=center]
      \tikz[baseline=-0.7ex]\node[circle,fill=white,draw=popink,line width=1.3pt,minimum size=1.6cm,inner sep=0]{\popdisplay\fontsize{24}{24}\selectfont\color{poppurple}+};%
      \hspace{8pt}\begin{minipage}[c]{0.6\linewidth}
        {\popdisplay\fontsize{11}{12.5}\selectfont\color{white}\MakeUppercase{Passe à OnBuch+}}\par\vspace{4pt}
        {\raggedright\popsemi\fontsize{8.4}{11}\selectfont\color{white}Tous les cours signés, Tuteur IA illimité, fiches PDF — onglet OnBuch+ de l'app\par}
      \end{minipage}
    \end{tcolorbox}
  \end{minipage}
  \vspace{8pt}
  \popcontenu
  \vfill
  \signatureonbuch
  \restoregeometry
  \newpage
}

% Rangée « Ce cours contient » (page de garde)
\newcommand{\poptile}[3]{% #1 couleur #2 gros texte #3 légende
  \begin{tcolorbox}[enhanced,colback=#1,colframe=popink,boxrule=2pt,arc=9pt,shadow={2.5pt}{-2.5pt}{0pt}{popink},
      left=6pt,right=6pt,top=9pt,bottom=9pt,halign=center,valign=center,height=1.75cm,width=\linewidth,nobeforeafter]
    {\popdisplay\fontsize{12.5}{13}\selectfont\color{white}\MakeUppercase{#2}}\par\vspace{3pt}
    {\centering\popsemi\fontsize{7.8}{9.5}\selectfont\color{white}#3\par}
  \end{tcolorbox}}
\newcommand{\popcontenu}{%
  \par\noindent{\popxboldls\fontsize{7.5}{7.5}\selectfont\color{popmuted}CE COURS SIGNÉ CONTIENT}\par\vspace{7pt}
  \noindent\begin{minipage}[t]{0.235\linewidth}\poptile{popblue}{Cours}{complet, illustré, pas à pas}\end{minipage}\hfill
  \begin{minipage}[t]{0.235\linewidth}\poptile{poporange}{Méthodes}{et exercices résolus}\end{minipage}\hfill
  \begin{minipage}[t]{0.235\linewidth}\poptile{poppink}{Exercices}{type examen + corrigés}\end{minipage}\hfill
  \begin{minipage}[t]{0.235\linewidth}\poptile{popgreen}{Fiche bilan}{l'essentiel à retenir}\end{minipage}\par}

% Sommaire
\newtcolorbox{sommaire}{enhanced,colback=white,colframe=popink,boxrule=2.2pt,arc=10pt,
  drop shadow=popink,left=18pt,right=18pt,top=13pt,bottom=13pt,before skip=6pt,after skip=20pt,
  before upper={\poppill{Sommaire}{poppurple}{white}\par\vspace{9pt}\popsemi\fontsize{10.6}{14.5}\selectfont\color{popink}}}

% Signature de fin de cours — poussée tout en bas de la dernière page
\newcommand{\findecours}{%
  \par\vfill\nopagebreak
  \begin{center}\popsquiggle{poporange}\end{center}\vspace{4pt}
  \signatureonbuch
}
\AtBeginDocument{\def\llbracket{\mathopen{[\mkern-3.5mu[}}\def\rrbracket{\mathclose{]\mkern-3.5mu]}}} % intervalles d'entiers (glyphe ⟦ absent de la police)
% Glyphes absents de Fira Math : redéfinis à partir de glyphes disponibles
\AtBeginDocument{%
  \def\setminus{\mathbin{\backslash}}\let\smallsetminus\setminus
  \def\vdots{\mathord{\vbox{\baselineskip3.2pt\lineskiplimit0pt\kern4pt\hbox{.}\hbox{.}\hbox{.}}}}%
  \def\ddots{\mathinner{\mkern1mu\raise6.5pt\hbox{.}\mkern2mu\raise3.6pt\hbox{.}\mkern2mu\raise0.7pt\hbox{.}\mkern1mu}}%
  \def\triangle{\mathord{\scalebox{0.9}{$\symup{\Delta}$}}}%
  \def\nabla{\mathord{\raisebox{\depth}{\scalebox{1}[-1]{$\symup{\Delta}$}}}}%
  \def\checkmark{\popcheck{popdark}}%
  \def\star{\mathbin{\ast}}\def\bigstar{\mathord{\ast}}%
  \def\diamond{\mathbin{\scalebox{0.6}{$\Diamond$}}}%
  \def\bigcup{\mathop{\mathchoice{\raisebox{-0.3em}{\scalebox{1.7}{$\cup$}}}{\scalebox{1.25}{$\cup$}}{\cup}{\cup}}\displaylimits}%
  \def\bigcap{\mathop{\mathchoice{\raisebox{-0.3em}{\scalebox{1.7}{$\cap$}}}{\scalebox{1.25}{$\cap$}}{\cap}{\cap}}\displaylimits}%
  \def\lesseqqgtr{\mathrel{\lessgtr}}\def\bigoplus{\mathop{\oplus}\displaylimits}%
}
\providecommand{\ssi}{si et seulement si} % abréviation fréquente
\AtBeginDocument{\providecommand{\wideparen}{\overparen}} % arc de cercle

%% FIGFIX-BEGIN v5 — correctifs globaux des figures (docs/onbuch-generation/tools/figfix)
\usepackage{adjustbox}\usepackage{etoolbox}\tcbuselibrary{raster}
% toute figure TikZ/pgfplots plus large que la ligne est réduite à la largeur disponible
\BeforeBeginEnvironment{tikzpicture}{\begin{adjustbox}{max width=\linewidth}}
\AfterEndEnvironment{tikzpicture}{\end{adjustbox}}
% légendes pgfplots lisibles par-dessus les courbes
\pgfplotsset{every axis/.append style={legend style={fill=white,fill opacity=0.92,text opacity=1,draw=black!15,inner sep=3pt}}}
% étiquettes d'arêtes (syntaxe edge["..."]) avec fond blanc
\tikzset{every edge quotes/.append style={fill=white,inner sep=1.2pt}}
%% FIGFIX-END
\begin{document}

\pagegarde{Le fonctionnement de la mondialisation}{Géographie}{Tle A}{Module 2 : La libéralisation des échanges (la mondialisation)}{N12}{2 h + dossiers 3 et 4 (4 h)}{Cette leçon analyse les mécanismes concrets de la mondialisation : les flux physiques de marchandises, les circulations de personnes, d'informations et de capitaux, ainsi que le commerce en ligne et les IDE. Elle aborde ensuite, à travers deux dossiers guidés, le commerce équitable comme réponse aux déséquilibres mondiaux et les problèmes majeurs liés à la mondialisation.
\vspace{4pt}
\begin{multicols}{2}\small
\begin{itemize}
\item À la fin de cette leçon, je sais définir la mondialisation et identifier ses principaux flux (marchandises, personnes, informations, capitaux).
\item Je sais observer, localiser et décrire sur une carte les grands flux mondiaux de produits agricoles et miniers.
\item Je sais légender une carte et construire un croquis des flux commerciaux mondiaux.
\item Je sais construire et interpréter un diagramme (bâtons, circulaire) à partir de données chiffrées sur les IDE ou les transferts de la diaspora.
\item Je sais inventorier, classer et interpréter des documents (cartes, photographies, tableaux statistiques) relatifs à la mondialisation.
\item Je sais expliquer le concept de commerce équitable, ses normes et ses critiques.
\end{itemize}\end{multicols}}

\begin{sommaire}
\begin{multicols}{2}
\begin{enumerate}
\item Les flux physiques de marchandises
\item La circulation des personnes
\item La circulation de l'information, des services et le commerce en ligne
\item Les mouvements financiers et les IDE
\item Dossier 3 : Le commerce équitable (activité guidée)
\item Dossier 4 : Les problèmes liés à la mondialisation (activité guidée)
\item TP : Cartographier la mondialisation et agir compétent
\item Activité d'intégration
\item Méthodes et exercices résolus
\item Exercices
\item Corrigés des exercices
\item Fiche bilan
\end{enumerate}
\end{multicols}
\end{sommaire}

\begin{objectifs}
\begin{itemize}
\item À la fin de cette leçon, je sais définir la mondialisation et identifier ses principaux flux (marchandises, personnes, informations, capitaux).
\item Je sais observer, localiser et décrire sur une carte les grands flux mondiaux de produits agricoles et miniers.
\item Je sais légender une carte et construire un croquis des flux commerciaux mondiaux.
\item Je sais construire et interpréter un diagramme (bâtons, circulaire) à partir de données chiffrées sur les IDE ou les transferts de la diaspora.
\item Je sais inventorier, classer et interpréter des documents (cartes, photographies, tableaux statistiques) relatifs à la mondialisation.
\item Je sais expliquer le concept de commerce équitable, ses normes et ses critiques.
\item Je sais analyser les problèmes liés à la mondialisation : immigration clandestine, cybercriminalité, conflits armés, marginalisation des pays pauvres.
\item Je sais proposer des actions de labellisation des produits camerounais pour une meilleure participation à la mondialisation.
\end{itemize}
\end{objectifs}

\begin{prerequis}
\begin{itemize}
\item La notion de mondialisation et ses acteurs (FTN, États, organisations internationales) vue en début de module.
\item Les inégalités de développement Nord/Sud (classe de Première).
\item Les notions de balance commerciale, exportations et importations.
\item La lecture et la légendation d'une carte (savoir-faire de Seconde).
\item Les grandes zones d'échanges et la Triade (amorce de la leçon 13).
\end{itemize}
\end{prerequis}

\newpage

\coursec{Les flux physiques de marchandises}

\textbf{Situation d'accroche.} Samedi matin, marché Mokolo à Yaoundé. Amina remplit son panier : un sac de cacao venu des plantations de l'Ouest, du riz parfumé importé de Thaïlande, et un téléphone portable fabriqué en Chine, payé avec l'argent que son frère, installé à Lyon, lui a envoyé la veille par transfert. En une seule matinée, sans quitter son quartier, Amina a participé à plusieurs flux mondiaux : des marchandises (cacao, riz, téléphone), des capitaux (le transfert de son frère), des personnes (la migration de ce dernier) et des informations (le message téléphonique qui a déclenché l'envoi). Pourtant, un détail la frappe : le kilo de cacao du planteur camerounais se vend quelques centaines de francs CFA, alors que la tablette de chocolat fabriquée avec ce même cacao se vend dix fois plus cher en Europe. Comment la mondialisation fonctionne-t-elle concrètement ? Qui en profite, qui en est exclu, et comment le Cameroun peut-il mieux s'y insérer ? C'est à ces questions que cette leçon va répondre, en commençant par le plus visible des flux : celui des marchandises.

\courssub{La mondialisation : une intensification des échanges}

\begin{definition}[La mondialisation]
La \cle{mondialisation} est le processus d'intensification des échanges (biens, personnes, capitaux, informations) qui met les espaces du monde en \cle{interdépendance} : ce qui se produit, se consomme ou se décide à un endroit du globe a des conséquences sur d'autres espaces, parfois très éloignés.
\end{definition}

L'intuition est simple : aucun pays ne produit tout ce qu'il consomme. Le Cameroun a du pétrole et du cacao, mais peu d'usines de téléphones ; la Chine a des usines, mais manque de pétrole ; l'Europe a des capitaux et des technologies, mais peu de cacao. Chaque espace échange donc ce qu'il possède en abondance contre ce qui lui manque. La mondialisation n'est pas un phénomène totalement nouveau (les caravanes transsahariennes et les comptoirs coloniaux en étaient des formes anciennes), mais elle s'est \textbf{accélérée} depuis les années 1980 grâce à trois moteurs : la baisse des coûts de transport (conteneurisation, grands porte-conteneurs), la révolution numérique (Internet, téléphonie mobile) et la libéralisation des échanges (baisse des droits de douane, accords commerciaux).

\begin{aretenir}[Les quatre grands flux de la mondialisation]
\begin{itemize}
\item les flux de \cle{marchandises} (produits agricoles, minerais, produits manufacturés) ;
\item la circulation des \cle{personnes} (migrations, tourisme) ;
\item la circulation de l'\cle{information} et des \cle{services} (Internet, banques, commerce en ligne) ;
\item les mouvements \cle{financiers} (IDE, transferts de la diaspora, fonds multilatéraux).
\end{itemize}
Cette section traite le premier flux ; les trois autres feront l'objet des sections suivantes.
\end{aretenir}

\courssub{Les grands produits agricoles échangés}

Les produits agricoles constituent une part essentielle du commerce mondial et la principale richesse exportée par de nombreux pays du Sud. On distingue deux grandes familles.

\textbf{Les céréales.} Le \cle{blé}, le \cle{riz} et le \cle{maïs} nourrissent l'humanité et circulent massivement. Les grands exportateurs de blé sont les États-Unis, le Canada, la Russie, l'Ukraine et la France ; le riz part surtout de l'Asie (Thaïlande, Vietnam, Inde) vers l'Afrique et le Moyen-Orient. Le riz thaïlandais acheté par Amina à Mokolo illustre ce flux : le Cameroun, qui consomme beaucoup plus de riz qu'il n'en produit, importe chaque année des centaines de milliers de tonnes (ordre de grandeur : plus de \num{600000} tonnes, voire \num{900000} tonnes certaines années).

\textbf{Les produits tropicaux.} Le \cle{cacao}, le \cle{café}, la \cle{banane}, le \cle{coton} et l'\cle{huile de palme} ne poussent que sous les tropiques, mais sont consommés surtout dans les pays riches du Nord. C'est le cas typique du cacao : produit en Côte d'Ivoire, au Ghana, au Cameroun, il est transformé et consommé en Europe et en Amérique du Nord. Le Cameroun exporte aussi du café arabica et robusta, des bananes (via la PHP, Plantations du Haut Penja), du coton (SODECOTON dans le Nord) et de l'huile de palme (SOCAPALM, PHC).

\courssub{Les minerais et les produits énergétiques}

La deuxième grande catégorie de marchandises échangées concerne les ressources du sous-sol.

\begin{itemize}
\item Le \cle{pétrole} et le \cle{gaz} : produits au Moyen-Orient (Arabie saoudite, Irak, Qatar), en Russie, en Afrique (Nigeria, Angola, Algérie, Cameroun de façon modeste) et en Amérique (États-Unis, Venezuela), ils sont consommés partout, surtout dans la Triade (Amérique du Nord, Europe, Asie de l'Est).
\item Le \cle{fer} : extrait en Australie, au Brésil, en Russie, il alimente les aciéries de Chine, premier sidérurgiste mondial.
\item La \cle{bauxite} : l'Australie, la Guinée et le Brésil dominent ; le Cameroun dispose d'importants gisements encore peu exploités (Minim-Martap, Ngaoundal).
\item Le \cle{cobalt} : la République démocratique du Congo fournit plus de la moitié de la production mondiale, indispensable aux batteries des téléphones et des voitures électriques.
\end{itemize}

On observe une régularité fondamentale : les \textbf{zones de production} (souvent au Sud) ne coïncident pas avec les \textbf{zones de consommation} (surtout au Nord et en Asie de l'Est). Cette dissociation spatiale oblige les matières premières à voyager, d'où l'importance des routes maritimes.

\courssub{Les routes maritimes et les points de passage stratégiques}

\begin{exemplebox}[Le règne du maritime]
Plus de \cle{80 \%} du commerce mondial en volume transite par voie \cle{maritime}, car le navire est le mode de transport le moins cher à la tonne. Le détroit de \cle{Malacca}, entre la Malaisie et l'Indonésie, voit passer environ un \textbf{quart} du commerce mondial : c'est le passage obligé entre l'océan Indien et la Chine.
\end{exemplebox}

Les flux maritimes empruntent des itinéraires contraints par la géographie : les \cle{détroits} et les \cle{canaux} sont des « goulots d'étranglement » où le trafic se concentre et qui ont une valeur stratégique immense.

\begin{itemize}
\item Le détroit d'\cle{Ormuz} (entre l'Iran et Oman) : passage de l'essentiel du pétrole du golfe Persique ;
\item le détroit de \cle{Malacca} : liaison Asie du Sud-Est -- Chine -- Japon ;
\item le détroit de \cle{Gibraltar} : porte occidentale de la Méditerranée ;
\item le canal de \cle{Suez} (Égypte, 1869) : raccourci Europe -- Asie sans contourner l'Afrique ;
\item le canal de \cle{Panama} (1914) : passage entre Atlantique et Pacifique sans contourner l'Amérique du Sud.
\end{itemize}

\begin{popfigure}
\includegraphics[width=\linewidth,height=9.5cm,keepaspectratio]{N12/shipping_routes.jpg}
\legende{Le trafic maritime commercial mondial (plus le bleu est foncé, plus le trafic est dense ; carte sans noms de lieux). Les trois grands axes se lisent d'un coup d'œil : l'Atlantique Nord, le Pacifique Nord et l'axe Méditerranée--Suez--océan Indien--Asie de l'Est. Les trafics se resserrent aux passages obligés : détroit de Gibraltar, canal de Suez, détroit de Malacca, canal de Panama. La carte ne distingue pas les produits transportés (pétrole, conteneurs) : elle montre le trafic de l'ensemble des navires de commerce. \\ Source : Wikimedia Commons, « Shipping routes » (densité du trafic maritime commercial), T. Hengl, CC BY-SA 3.0.}
\end{popfigure}

\begin{methode}[Légender une carte de flux]
\begin{enumerate}
\item Choisir des \textbf{flèches proportionnelles} aux volumes transportés (épaisseur ou largeur variable) ;
\item indiquer une \textbf{échelle} claire (ex. : 1 mm de largeur = 10 millions de tonnes) ;
\item rédiger un \textbf{titre précis} annonçant la nature du flux, sa direction et la date (ex. : « Les flux mondiaux de pétrole en 2020 ») ;
\item placer une \textbf{légende ordonnée} (du flux le plus important au plus faible) et l'orientation (flèche du nord).
\end{enumerate}
\end{methode}

\courssub{Les ports et les hubs : les portes du commerce mondial}

Les flux maritimes s'organisent autour de grands ports-hubs qui concentrent, trient et redistribuent les conteneurs. \cle{Shanghai} (Chine) est le premier port à conteneurs du monde, devant \cle{Singapour}, véritable plaque tournante de l'Asie du Sud-Est, et \cle{Rotterdam} (Pays-Bas), première porte d'entrée de l'Europe. Ces ports sont des interfaces : ils connectent le maritime aux réseaux ferroviaires, routiers et fluviaux de leur arrière-pays (le \emph{hinterland}).

Pour le Cameroun, le \cle{port de Douala} joue ce rôle de porte d'entrée : il traite l'écrasante majorité du commerce extérieur du pays (ordre de grandeur : plus de 90 \%) et dessert aussi le Tchad et la Centrafrique, pays enclavés. Le port en eaux profondes de \cle{Kribi}, dont la première phase a été mise en service en 2014, complète ce dispositif et accueille les plus grands navires (jusqu'à 300 m de long).

\courssub{La direction des flux : Sud-Nord, Nord-Sud et Sud-Sud}

Les flux de marchandises suivent des directions lisibles sur une carte :

\begin{itemize}
\item \textbf{Sud $\rightarrow$ Nord} : les pays en développement exportent des matières premières (cacao, pétrole, minerais) vers les pays riches ;
\item \textbf{Nord $\rightarrow$ Sud} : les pays riches exportent des produits manufacturés (machines, voitures, médicaments) à forte valeur ajoutée ;
\item \textbf{Sud $\rightarrow$ Sud} : flux croissants entre pays émergents et pays en développement, symbolisés par les échanges \cle{Chine -- Afrique} (la Chine est devenue le premier fournisseur et un partenaire commercial majeur du Cameroun et de nombreux pays africains).
\end{itemize}

\begin{attention}[Erreur fréquente]
Ne crois pas que les pays du Sud n'exportent \emph{que} des matières premières. La \cle{Chine} et plusieurs pays émergents (Corée du Sud, Vietnam, Bangladesh) exportent massivement des \textbf{produits manufacturés} : le téléphone d'Amina, fabriqué en Chine, en est la preuve. La fracture Nord/Sud est utile comme grille de lecture, mais elle est de plus en plus brouillée par l'émergence de nouvelles puissances industrielles.
\end{attention}

\courssub{Étude de cas : la filière cacao camerounaise}

Suivons une fève de cacao camerounais, de la plantation à la tablette de chocolat. Le planteur de la région du Centre ou du Littoral récolte les cabosses, les écabosse, fait fermenter et sécher les fèves, puis les vend à un acheteur agréé ou à une coopérative. Les fèves sont transportées par camion jusqu'au \cle{port de Douala}, d'où elles partent par navire vers l'Europe (Pays-Bas, Allemagne, Belgique). Là-bas, les chocolateries les torréfient, les broient et fabriquent le chocolat vendu dans le monde entier.

\begin{popfigure}
\begin{tikzpicture}[font=\small, node distance=0.5cm]
\node[draw, rounded corners, fill=popgreenL, text width=2.6cm, align=center] (prod) {Planteur camerounais\\ (Centre, Littoral, Sud-Ouest)};
\node[draw, rounded corners, fill=popgoldL, text width=2.4cm, align=center, right=1.1cm of prod] (douala) {Port de Douala\\ (exportation des fèves)};
\node[draw, rounded corners, fill=popblueL, text width=2.6cm, align=center, right=1.1cm of douala] (choco) {Chocolaterie européenne\\ (transformation, vente)};
\draw[-{Stealth}, very thick, popdark] (prod) -- (douala) node[midway, above, font=\scriptsize] {fèves brutes};
\draw[-{Stealth}, very thick, popdark] (douala) -- (choco) node[midway, above, font=\scriptsize] {transport maritime};
\end{tikzpicture}
\legende{Schéma de la filière cacao : du producteur camerounais au chocolatier européen via le port de Douala.}
\end{popfigure}

\begin{center}
\begin{tabularx}{\linewidth}{|l|X|X|X|}
\hline
\rowcolor{popblueL} \textbf{Étape de la filière} & \textbf{Planteur (production)} & \textbf{Transport, négoce, logistique} & \textbf{Transformation et distribution} \\
\hline
Part indicative du prix final & environ 5 à 10 \% & environ 20 \% & environ 70 à 75 \% \\
\hline
Localisation principale & Cameroun (Sud) & International (navires, courtiers) & Europe, Amérique du Nord (Nord) \\
\hline
Valeur ajoutée & Faible (matière première brute) & Moyenne & Élevée (technologie, marques, marketing) \\
\hline
\end{tabularx}
\end{center}

Le tableau ci-dessus montre le cœur du problème : le planteur, qui fournit le travail le plus pénible, ne touche qu'une faible part du prix final de la tablette, car la \cle{valeur ajoutée} se concentre dans la transformation et la distribution, réalisées au Nord. C'est ce qu'on appelle la \textbf{dégradation des termes de l'échange} : les matières premières brutes rapportent peu et leurs prix fluctuent fortement sur les marchés mondiaux (bourse de Londres, New York).

\begin{savaistu}[Le cacao camerounais]
Le Cameroun est environ le \cle{5\textsuperscript{e} producteur mondial} de cacao (ordre de grandeur : 250 000 à 300 000 tonnes par an), derrière la Côte d'Ivoire, le Ghana, l'Indonésie et l'Équateur. Pourtant, il transforme \textbf{localement environ 30 \%} de sa production (broyaige, beurre, poudre, chocolat) : l'essentiel part encore en fèves brutes. Des unités de transformation (comme les usines de la SIC Cacaos à Douala, Neo Industry à Kribi) cherchent à retenir davantage de valeur ajoutée au pays.
\end{savaistu}

\begin{exoresolu}[Commenter la répartition de la valeur dans la filière cacao]
Énoncé. À partir du schéma et du tableau de la filière cacao, expliquez en une dizaine de lignes pourquoi le producteur camerounais profite peu de la mondialisation, et proposez une solution.
\tcblower
Solution détaillée. Le producteur camerounais vend des fèves brutes, produit à faible valeur ajoutée dont le prix est fixé sur les marchés mondiaux (Londres, New York) et subit de fortes variations. Il ne touche qu'environ 5 à 10 \% du prix final de la tablette de chocolat. L'essentiel de la valeur (70 à 75 \%) est créé lors de la transformation (torréfaction, broyage, fabrication) et de la distribution, activités localisées en Europe, où se trouvent les capitaux, les technologies et les marques. Le planteur subit donc la mondialisation plus qu'il n'en profite. Pour mieux s'insérer, le Cameroun peut développer la \textbf{transformation locale} (usines de broyage, chocolateries), ce qui créerait des emplois et augmenterait la valeur des exportations ; il peut aussi labéliser sa production (commerce équitable, labels de qualité) pour obtenir de meilleurs prix, comme le montre le dossier 3 de cette leçon.
\end{exoresolu}

\begin{exercice}[À toi de jouer]
Sur une esquisse du planisphère, localise et nomme cinq points de passage stratégiques du commerce mondial, puis trace d'une flèche le trajet d'une cargaison de pétrole du golfe Persique vers le Japon, et d'une autre flèche le trajet du cacao camerounais vers Rotterdam. N'oublie ni le titre, ni la légende, ni l'orientation.
\end{exercice}

\begin{corrige}[Exercice : cartographier les flux]
La cargaison de pétrole part du détroit d'Ormuz, traverse l'océan Indien, passe par le détroit de Malacca, puis remonte la mer de Chine jusqu'au Japon. Le cacao part du port de Douala, longe les côtes ouest-africaines, passe au large du détroit de Gibraltar (ou directement vers le nord) et rejoint Rotterdam. Le titre doit préciser la nature des flux et la date (ex. « Flux de pétrole et de cacao, 2023 ») ; la légende distingue deux couleurs de flèches (ex. orange pour le pétrole, marron pour le cacao) ; la flèche du nord est orientée vers le haut.
\end{corrige}

\begin{aretenir}[L'essentiel de la section]
La mondialisation repose d'abord sur des flux physiques de marchandises : produits agricoles (céréales, produits tropicaux) et minerais (pétrole, fer, bauxite, cobalt) voyagent surtout par voie maritime (plus de 80 \% du commerce en volume), le long de routes contrôlées par des détroits et canaux stratégiques (Ormuz, Malacca, Suez, Panama, Gibraltar) et organisées autour de grands ports-hubs (Shanghai, Singapour, Rotterdam, Douala pour le Cameroun). Les flux vont des matières premières du Sud vers le Nord, des produits manufacturés du Nord vers le Sud, et de plus en plus entre pays du Sud. L'étude de la filière cacao montre que la valeur ajoutée se concentre là où se fait la transformation : c'est tout l'enjeu de l'insertion du Cameroun dans la mondialisation.
\end{aretenir}

\coursec{La circulation des personnes}

Dans la section précédente, nous avons suivi les porte-conteneurs qui transportent le pétrole du golfe Persique, le cacao du port de Douala ou le café d'Oku vers les usines du Nord. Mais la mondialisation ne fait pas voyager que des marchandises : elle fait aussi voyager les \emph{hommes}. À l'aéroport international de Douala, chaque jour, des centaines de Camerounais s'envolent vers Paris, Bruxelles, Montréal ou Dubaï ; chaque semaine, des touristes européens débarquent à Kribi pour découvrir les chutes de la Lobé. Qui sont ces voyageurs ? Où vont-ils ? Pourquoi partent-ils ? C'est toute la question de la \cle{circulation des personnes}, deuxième grand flux de la mondialisation.

\courssub{Les migrations internationales : volume et grandes directions}

\begin{definition}[Migration internationale]
Une \cle{migration internationale} est un déplacement durable d'une personne qui franchit une frontière d'État pour s'installer, de manière temporaire ou définitive, dans un pays autre que son pays de résidence habituelle. On parle d'\emph{émigration} du point de vue du pays de départ et d'\emph{immigration} du point de vue du pays d'accueil.
\end{definition}

Commençons par une intuition simple. Imagine la population du Cameroun tout entière — environ 28 millions d'habitants — multipliée par dix : tu obtiens à peu près le nombre de personnes qui vivent aujourd'hui hors de leur pays de naissance. Selon l'Organisation internationale pour les migrations (OIM) et la Division de la population de l'ONU, on compte environ \cle{281 millions de migrants internationaux} dans le monde en 2020, soit près de \num{3,6} \% de la population mondiale. Ce chiffre peut sembler modeste en proportion, mais il a presque doublé en trente ans : les migrations internationales sont un phénomène en forte croissance, comme le montre le diagramme ci-dessous.

\begin{popfigure}
\begin{center}
\begin{tikzpicture}
\begin{axis}[
    ybar,
    width=0.9\linewidth,
    height=7cm,
    bar width=0.6cm,
    ymin=0, ymax=320,
    ylabel={Migrants internationaux (millions)},
    symbolic x coords={1990,1995,2000,2005,2010,2015,2020},
    xtick=data,
    ytick={0,50,100,150,200,250,300},
    axis lines=left,
    every axis plot/.append style={fill=popblue, draw=popdark},
    nodes near coords,
    nodes near coords style={font=\small, color=popdark, anchor=south, yshift=2pt},
    enlarge x limits=0.08,
    tick label style={font=\small}
]
\addplot coordinates {(1990,153) (1995,161) (2000,173) (2005,191) (2010,221) (2015,248) (2020,281)};
\end{axis}
\end{tikzpicture}
\end{center}
\legende{Évolution du nombre de migrants internationaux dans le monde de 1990 à 2020 (en millions, valeurs arrondies). Source : ONU, Département des affaires économiques et sociales ; OIM, \emph{Rapport sur les migrations dans le monde} 2022.}
\end{popfigure}

\begin{attention}[Lire correctement ce graphique]
Ne conclus pas trop vite que « tout le monde migre ». Environ 96 \% de l'humanité vit toujours dans son pays de naissance. La mondialisation des personnes est réelle et croissante, mais elle reste minoritaire : ce sont surtout les \emph{marchandises}, les \emph{capitaux} et l'\emph{information} qui circulent massivement, beaucoup plus librement que les hommes, dont les déplacements sont contrôlés par les visas et les frontières.
\end{attention}

Dans quelles directions ces migrants se déplacent-ils ? L'analyse des flux fait apparaître deux grandes orientations :

\begin{itemize}
\item Les flux \cle{Sud-Nord} : des pays en développement (Afrique, Amérique latine, Asie du Sud) vers les pays riches (Amérique du Nord, Europe de l'Ouest, golfe Persique). C'est l'image la plus connue : le jeune Malien qui tente de rejoindre l'Espagne, le Camerounais qui part étudier puis travailler en France ou au Canada.
\item Les flux \cle{Sud-Sud} : entre pays en développement eux-mêmes. Ils sont aussi importants en volume que les flux Sud-Nord, mais beaucoup moins médiatisés. Exemples : les Burkinabè dans les plantations de Côte d'Ivoire, les Tchadiens et Centrafricains au Cameroun, les Zimbabwéens en Afrique du Sud, les Népalais dans les chantiers du Qatar.
\end{itemize}

\begin{popfigure}
\includegraphics[width=\linewidth,height=9.5cm,keepaspectratio]{N12/migrant_stock.png}
\legende{Nombre total d'immigrés par pays d'accueil en 2020 (échelle du jaune clair au bleu foncé, de moins de \num{10000} à plus de 10 millions de personnes). Les pôles d'accueil les plus importants sont en bleu foncé : États-Unis, Allemagne, Royaume-Uni, Arabie saoudite ; viennent ensuite, en bleu plus clair, la France, le Canada, la Russie, l'Inde ou l'Australie. Le croquis attendu utilise cette répartition pour situer les foyers d'accueil (Amérique du Nord, Europe de l'Ouest, golfe Persique), puis trace les flèches de départ depuis les foyers d'émigration (Mexique et Amérique centrale, Afrique, Asie du Sud, Asie de l'Est). \\ Source : Wikimedia Commons, « Migrant stock total, World, 2020 », Our World in Data (données ONU DAES), CC BY 4.0.}
\end{popfigure}

\begin{methode}[Construire un croquis de flux migratoires]
Pour réaliser un croquis de flux migratoires au brouillon puis au propre, procède par étapes :
\begin{enumerate}
\item Trace le fond de carte : contours simplifiés des continents, titre en haut, flèche du nord.
\item Représente les \textbf{foyers de départ} et les \textbf{foyers d'accueil} par des taches ou des points nommés (ex. : Afrique de l'Ouest, Europe de l'Ouest, golfe Persique).
\item Dessine les \textbf{flux} par des flèches dont l'\emph{épaisseur est proportionnelle} à l'importance du flux : une flèche épaisse pour Mexique $\rightarrow$ États-Unis, une flèche fine pour un flux secondaire.
\item Rédige une \textbf{légende numérotée} qui nomme chaque flèche et chaque tache.
\item Vérifie la cohérence : le titre doit annoncer le sujet et la date ; aucune flèche ne doit rester sans explication dans la légende.
\end{enumerate}
\end{methode}

\courssub{Les migrations de travail : de la main-d'œuvre non qualifiée à la fuite des cerveaux}

La grande majorité des migrants internationaux sont des \cle{migrants de travail} : ils partent chercher un emploi et de meilleurs revenus. Mais tous ne jouent pas le même rôle dans l'économie mondiale.

D'un côté, la \textbf{main-d'œuvre non qualifiée} alimente les secteurs pénibles et mal payés des pays riches : le BTP et le ménage en Europe, l'agriculture en Californie, les chantiers des stades et des gratte-ciel dans les émirats du golfe Persique (Qatar, Émirats arabes unis, Arabie saoudite), qui emploient des millions d'ouvriers indiens, pakistanais, bangladais, népalais et philippins.

De l'autre côté, les pays riches organisent une véritable \emph{chasse aux talents} : médecins, ingénieurs, informaticiens, chercheurs. C'est le phénomène de la \cle{fuite des cerveaux}.

\begin{definition}[Fuite des cerveaux]
La \cle{fuite des cerveaux} (en anglais \emph{brain drain}) est l'émigration des personnes qualifiées (diplômés du supérieur, cadres, professionnels de santé) des pays pauvres ou émergents vers les pays riches, attirés par de meilleurs salaires et de meilleures conditions de travail.
\end{definition}

Ce phénomène pose un problème majeur aux pays de départ : ils financent la formation de leurs élites (écoles, universités publiques) puis les voient partir. L'Afrique subsaharienne perd ainsi chaque année une partie importante de ses médecins et infirmiers formés, au profit de la Grande-Bretagne, du Canada, de la France ou des États-Unis. Le Cameroun n'échappe pas à cette logique : de nombreux médecins, infirmiers et ingénieurs camerounais exercent aujourd'hui en Allemagne, au Canada ou aux États-Unis, alors que les hôpitaux de Yaoundé ou de Garoua manquent de personnel spécialisé. Certains pays tentent de transformer cette perte en atout en encourageant le \emph{retour des compétences} ou l'investissement de la diaspora.

\courssub{Les migrations contraintes : réfugiés et déplacés}

Toutes les migrations ne sont pas choisies. Des dizaines de millions de personnes fuient la guerre, les persécutions ou les catastrophes naturelles.

\begin{attention}[Réfugié ou migrant économique ?]
Ne confonds jamais le \cle{réfugié} et le \cle{migrant économique}. Le réfugié fuit une \emph{contrainte} (guerre, persécution, danger de mort) et bénéficie d'un \emph{statut juridique international} défini par la Convention de Genève de 1951 relative au statut des réfugiés et protégé par le Haut-Commissariat des Nations unies pour les réfugiés (HCR). Le migrant économique part \emph{volontairement} pour améliorer ses conditions de vie : il ne relève pas de ce statut. Confondre les deux dans une copie est une erreur grave de définition.
\end{attention}

Selon le HCR, le monde comptait en 2022 plus de \num{100} millions de personnes déracinées de force, dont environ 35 millions de réfugiés ayant franchi une frontière (Syriens en Turquie et au Liban, Afghans au Pakistan et en Iran, Vénézuéliens en Colombie, Ukrainiens en Pologne après 2022) et une majorité de \emph{déplacés internes}, restés à l'intérieur de leur propre pays. Le Cameroun est lui-même concerné des deux côtés : il accueille des réfugiés centrafricains à l'Est et nigérians dans l'Extrême-Nord (camp de Minawao, près de Mokolo), tout en comptant des déplacés internes fuyant les attaques de Boko Haram dans l'Extrême-Nord et la crise anglophone dans le Nord-Ouest et le Sud-Ouest. Les catastrophes (inondations, sécheresses, cyclones) provoquent aussi des déplacements, appelés parfois « migrations environnementales », même si ce statut n'existe pas encore juridiquement.

\courssub{Le tourisme international : une circulation de personnes de loisir}

La circulation des personnes ne concerne pas que ceux qui partent s'installer : elle comprend aussi les déplacements temporaires, et d'abord le \cle{tourisme international}, c'est-à-dire tout voyage de loisir ou d'affaires d'au moins une nuit hors de son pays de résidence.

Avant la crise sanitaire de 2020, le monde comptait environ \num{1,4} milliard de touristes internationaux par an (source : Organisation mondiale du tourisme, OMT). L'observation des flux révèle une loi essentielle :

\begin{itemize}
\item Les flux \cle{Nord-Nord} dominent largement : les Européens voyagent surtout en Europe (France, Espagne et Italie sont les trois premières destinations mondiales), les Américains au Canada et en Europe. Le tourisme international reste d'abord une affaire de pays riches entre eux.
\item Mais le tourisme vers le \textbf{Sud} connaît un essor rapide : Mexique, Thaïlande, Maroc, Égypte, République dominicaine attirent des millions de visiteurs du Nord, séduits par le soleil, les prix bas et les paysages. C'est une source de devises précieuse pour ces pays.
\end{itemize}

L'Afrique ne capte encore qu'une faible part du tourisme mondial (de l'ordre de 5 \% des arrivées, ordre de grandeur à manier avec prudence). Le Cameroun, pourtant surnommé « l'Afrique en miniature » pour la diversité de ses paysages (plages de Kribi et Limbé, monts Mandara, parcs de Waza et de la Bénoué, mont Cameroun), reçoit moins d'un million de visiteurs par an : un potentiel considérable mais encore sous-exploité, freiné par le coût des transports, l'état des infrastructures et l'insécurité dans certaines régions.

\courssub{Les diasporas : des communautés actives entre deux rives}

\begin{definition}[Diaspora]
Une \cle{diaspora} est l'ensemble des membres d'une même communauté nationale ou culturelle dispersés dans plusieurs pays étrangers, qui maintiennent des liens économiques, culturels et affectifs avec leur pays d'origine.
\end{definition}

Les diasporas sont des acteurs majeurs de la mondialisation, car elles créent des \emph{ponts permanents} entre les pays. Leur rôle est double :

\begin{itemize}
\item \textbf{Rôle économique} : les migrants envoient régulièrement de l'argent à leurs familles restées au pays. Ces \cle{transferts de fonds} financent la scolarité des enfants, la santé, la construction de maisons, les petits commerces.
\item \textbf{Rôle culturel} : les diasporas diffusent la langue, la cuisine, la musique de leur pays (pense au rayonnement mondial du makossa ou des artistes camerounais installés à Paris) et servent de relais pour les entreprises et les investisseurs.
\end{itemize}

\begin{exemplebox}[Les transferts des migrants, un fleuve financier]
Selon la Banque mondiale, les transferts d'argent des migrants vers les pays en développement dépassent \cle{600 milliards de dollars par an} (environ 647 milliards en 2022, ordre de grandeur). C'est \emph{plus du triple} de l'aide publique au développement (APD) versée par les États riches, et presque autant que les investissements directs étrangers. Pour des pays comme le Nigeria, le Sénégal ou le Cameroun, ces envois représentent plusieurs points de PIB : une ressource vitale, et souvent plus stable que les capitaux spéculatifs.
\end{exemplebox}

La \textbf{diaspora camerounaise} est estimée à plusieurs centaines de milliers de personnes (certains avancent plus d'un million, mais les chiffres sont incertains : retiens un ordre de grandeur prudent). Elle est concentrée en \emph{Europe} (France en tête, puis Allemagne, Belgique, Grande-Bretagne, Italie) et en \emph{Amérique du Nord} (États-Unis, Canada, notamment le Québec francophone). Très qualifiée — étudiants, médecins, ingénieurs, cadres —, elle envoie chaque année des centaines de milliards de francs CFA au pays (ordre de grandeur : plusieurs centaines de millions de dollars par an selon la Banque mondiale), finance la construction de nombreuses villas à Douala, Yaoundé ou Bafoussam, et investit dans les écoles privées, les cliniques et l'agriculture. Des associations de ressortissants (par village ou par région) collectent aussi des fonds pour des projets collectifs : forages, salles de classe, centres de santé.

\courssub{Les espaces d'accueil privilégiés}

La carte mondiale des migrations montre que les migrants se concentrent dans quelques \cle{espaces d'accueil privilégiés}, qui correspondent aux pôles riches ou dynamiques de l'économie mondiale :

\begin{itemize}
\item \textbf{L'Amérique du Nord} : les États-Unis accueillent la première population immigrée du monde (plus de 45 millions de personnes nées à l'étranger), en provenance d'abord du Mexique, d'Amérique latine et d'Asie ; le Canada mène une politique d'immigration sélective très ouverte aux diplômés.
\item \textbf{L'Europe de l'Ouest} : Allemagne, France, Royaume-Uni, Italie et Espagne comptent des millions d'immigrés venus d'Afrique, de Turquie, d'Europe de l'Est et d'Asie. L'Allemagne est aujourd'hui le deuxième pays d'immigration du monde.
\item \textbf{Le golfe Persique} : cas spectaculaire, les monarchies pétrolières (Émirats arabes unis, Qatar, Koweït, Arabie saoudite) comptent jusqu'à 80 à 90 \% d'étrangers dans leur population active, employés dans le bâtiment, les services et le pétrole — mais sans perspective de naturalisation.
\end{itemize}

\begin{aretenir}[L'essentiel de la section]
\begin{itemize}
\item Environ \cle{281 millions de migrants internationaux} (2020), soit 3,6 \% de l'humanité, un chiffre en forte hausse depuis 1990.
\item Deux grandes directions : flux \cle{Sud-Nord} et flux \cle{Sud-Sud}, tout aussi volumineux.
\item Migrations de travail : main-d'œuvre non qualifiée (chantiers, agriculture, services) et \cle{fuite des cerveaux} des pays pauvres vers les pays riches.
\item Migrations contraintes : \cle{réfugiés} (statut juridique, Convention de Genève 1951, HCR) et déplacés internes — à ne pas confondre avec les migrants économiques.
\item Tourisme international : flux \cle{Nord-Nord} dominants, essor des destinations du Sud.
\item Les \cle{diasporas} transfèrent plus de 600 milliards de dollars par an vers le Sud, davantage que l'aide publique au développement ; la diaspora camerounaise est active en Europe et en Amérique du Nord.
\item Trois espaces d'accueil privilégiés : Amérique du Nord, Europe de l'Ouest, golfe Persique.
\end{itemize}
\end{aretenir}

\begin{exoresolu}[Commentaire de document : lire un tableau de flux migratoires]
\textbf{Énoncé.} Le tableau suivant indique les cinq premiers stocks de migrants internationaux par pays d'accueil en 2020 (en millions, ordres de grandeur, source ONU) : États-Unis : 51 ; Allemagne : 16 ; Arabie saoudite : 13 ; Russie : 12 ; Royaume-Uni : 9. Par ailleurs, les transferts de fonds des migrants vers les pays en développement atteignent environ 650 milliards de dollars en 2022, contre environ 200 milliards d'aide publique au développement. \textbf{Questions :} 1) Quels types d'espaces dominent l'accueil des migrants ? 2) Montrez que les migrations profitent aussi aux pays de départ. 3) Pourquoi la présence de l'Arabie saoudite dans ce classement illustre-t-elle une direction particulière des flux ?
\tcblower
\textbf{Solution détaillée.}
\textbf{1)} Les cinq premiers pays d'accueil sont tous des espaces riches ou à forte rente : deux pôles de la Triade (Amérique du Nord avec les États-Unis, Europe de l'Ouest avec l'Allemagne et le Royaume-Uni) et une monarchie pétrolière du golfe Persique (Arabie saoudite). Les migrants se concentrent donc là où se trouvent les emplois et les revenus élevés : les espaces d'accueil privilégiés sont les cœurs de l'économie mondiale.
\textbf{2)} Les pays de départ bénéficient des transferts de fonds de leurs diasporas : avec environ 650 milliards de dollars par an, ces envois dépassent plus de trois fois l'aide publique au développement ($650 / 200 = \num{3,25}$). Ils financent la consommation des familles, l'éducation, la santé et l'investissement (maisons, commerces). Les migrations alimentent donc un flux financier Sud-destiné considérable, qui compense partiellement la fuite des cerveaux.
\textbf{3)} L'Arabie saoudite illustre la direction \emph{Sud-Sud} et l'attractivité du golfe Persique : ses 13 millions de migrants sont pour l'essentiel des travailleurs venus d'Asie du Sud et du Sud-Est (Inde, Pakistan, Bangladesh, Philippines) et du monde arabe, employés dans le pétrole, le bâtiment et les services. Ce flux, peu médiatisé en Europe, est pourtant l'un des plus importants du monde.
\end{exoresolu}

\begin{savaistu}[Le Cameroun, carrefour migratoire de l'Afrique centrale]
Le Cameroun n'est pas seulement un pays de départ : c'est aussi un pays d'\emph{accueil} et de \emph{transit}. Il accueille des centaines de milliers de réfugiés (Centrafricains à l'Est, Nigérians dans l'Extrême-Nord), des migrants de travail tchadiens, nigérians et centrafricains dans le commerce et l'artisanat, et il voit transiter des candidats à l'émigration vers la Libye puis l'Europe. La mondialisation des personnes se joue donc aussi sur notre territoire.
\end{savaistu}

\begin{exercice}[À toi de jouer]
1) Définis en une phrase : migration internationale, fuite des cerveaux, diaspora. 2) Calcule la part des migrants internationaux dans la population mondiale en 2020 (281 millions de migrants pour environ \num{7,8} milliards d'habitants) et commente le résultat en deux lignes. 3) Sur un fond de planisphère, construis le croquis des principaux flux migratoires mondiaux en appliquant la méthode étudiée (taches, flèches proportionnelles, légende numérotée, titre, nord). 4) Explique en cinq lignes pourquoi le golfe Persique est devenu un espace d'accueil majeur alors qu'il ne faisait pas partie de la Triade.
\end{exercice}

\begin{corrige}[Exercice : La circulation des personnes]
\textbf{1)} Migration internationale : déplacement durable d'une personne franchissant une frontière d'État pour s'installer dans un autre pays. Fuite des cerveaux : émigration des personnes qualifiées des pays pauvres vers les pays riches. Diaspora : ensemble des membres d'une communauté dispersés à l'étranger qui gardent des liens avec leur pays d'origine.
\textbf{2)} Part des migrants : $281 / 7800 \times 100 \approx \num{3,6} \%$. Commentaire : les migrations internationales, bien qu'en forte croissance, ne concernent qu'une faible minorité de l'humanité ; la circulation des personnes reste beaucoup plus limitée que celle des marchandises, des capitaux et de l'information.
\textbf{3)} Le croquis doit comporter : les taches des foyers de départ (Afrique de l'Ouest et du Nord, Mexique et Amérique centrale, Asie du Sud) et d'accueil (Amérique du Nord, Europe de l'Ouest, golfe Persique) ; des flèches d'épaisseur proportionnelle (épaisse pour Mexique $\rightarrow$ États-Unis et Afrique $\rightarrow$ Europe, moyenne pour Asie du Sud $\rightarrow$ Golfe, fine pour les flux secondaires) ; une légende numérotée, un titre daté et la flèche du nord.
\textbf{4)} Le golfe Persique attire des millions de migrants grâce à la rente pétrolière et gazière, qui finance de gigantesques chantiers (gratte-ciel, stades, infrastructures) et des services exigeant une main-d'œuvre abondante que la faible population locale ne peut fournir. Les migrants y trouvent des salaires supérieurs à ceux de leurs pays, mais sans droit à la naturalisation : c'est une immigration purement économique et temporaire, alimentée par les flux Sud-Sud depuis l'Asie du Sud.
\end{corrige}

\coursec{La circulation de l'information, des services et le commerce en ligne}

Dans la section précédente, nous avons suivi les hommes et les femmes qui traversent les frontières : migrants, touristes, étudiants, commerçants. Mais une grande partie de ce qui circule dans la mondialisation ne prend ni avion, ni bateau, ni camion : il s'agit de l'\cle{information}. Quand un étudiant de Yaoundé suit un cours en ligne, quand une cliente de Douala commande un téléphone sur Jumia, quand une banque camerounaise transfère des données vers Paris, ce sont des flux immatériels qui se déplacent à la vitesse de la lumière. Cette section étudie ces flux invisibles mais décisifs : comment l'information circule, qui en profite, et comment le commerce s'est déplacé sur Internet.

\courssub{La révolution numérique : Internet, câbles sous-marins et satellites}

\textbf{Intuition.} Imagine que tu envoies une photo à un ami à Londres. En moins d'une seconde, l'image a traversé des milliers de kilomètres. Par où est-elle passée ? Contrairement à une idée répandue, ce n'est presque jamais par satellite : elle a voyagé dans un \cle{câble sous-marin} de fibre optique posé au fond de l'océan.

\begin{definition}[La révolution numérique]
La \cle{révolution numérique} désigne l'ensemble des transformations provoquées, depuis les années 1990, par la diffusion d'Internet et des technologies de l'information et de la communication (TIC) : ordinateurs, téléphones mobiles, réseaux de fibre optique, satellites. Elle permet la circulation quasi instantanée et à très faible coût de l'information à l'échelle de la planète.
\end{definition}

L'infrastructure physique de cette révolution repose sur trois piliers :

\begin{itemize}
\item les \textbf{câbles sous-marins de fibre optique}, qui transportent l'essentiel du trafic intercontinental ;
\item les \textbf{câbles terrestres} et les réseaux mobiles (2G, 3G, 4G, et désormais 5G dans les grandes villes) qui distribuent la connexion jusqu'aux utilisateurs ;
\item les \textbf{satellites}, utiles surtout pour les zones isolées, la télévision et le positionnement (GPS), mais minoritaires dans le trafic de données.
\end{itemize}

\begin{savaistu}[Des autoroutes au fond des océans]
Plus de \textbf{95 \%} du trafic Internet intercontinental transite par des câbles sous-marins, et non par satellite. Ces câbles, de la taille d'un tuyau d'arrosage, contiennent des fibres de verre plus fines qu'un cheveu. L'Afrique est aujourd'hui entourée de plusieurs grands câbles : \textbf{SAT-3} (2001), \textbf{WACS} (2012) et \textbf{ACE} (2012) sur la façade atlantique, \textbf{EASSy} et \textbf{SEACOM} sur la façade orientale. Le Cameroun est raccordé à SAT-3, WACS et ACE via Douala et Kribi, puis un réseau national de fibre relie les grandes villes (projet de « dorsale nationale »).
\end{savaistu}

\begin{popfigure}
\includegraphics[width=\linewidth,height=9.5cm,keepaspectratio]{N12/cables.png}
\legende{Carte mondiale des câbles sous-marins de télécommunications (traits rouges) et de leurs points d'atterrissement (ronds rouges), vers 2015. On y repère la densité des liaisons de l'Atlantique Nord et de l'Europe, de l'Asie de l'Est, puis le tour de l'Afrique par ses côtes atlantique et orientale : les câbles longent le continent et touchent terre en quelques points, dont ceux du golfe de Guinée où se trouve le Cameroun. La carte ne nomme pas chaque câble (SAT-3, WACS, ACE, EASSy, SEACOM) ; on les situe à partir du cours. \\ Source : Wikimedia Commons, « Submarine cable map umap », données de câbles de Greg Mahlknecht, fond OpenStreetMap contributors, CC BY-SA 2.0.}
\end{popfigure}

\begin{attention}[L'information circule vite, mais dans des tuyaux bien réels]
Ne dis pas que l'information circule « sans support » ou « par les airs ». Internet repose sur une infrastructure physique très concrète : câbles, stations d'atterrissement, centres de données, antennes. Une coupure de câble (chalut de pêche, séisme sous-marin) peut ralentir Internet dans tout un pays. La mondialisation immatérielle dépend donc de territoires et d'équipements localisés — et ceux qui les contrôlent disposent d'un pouvoir considérable.
\end{attention}

\courssub{La fracture numérique : une mondialisation à plusieurs vitesses}

\begin{definition}[Fracture numérique]
La \cle{fracture numérique} désigne les inégalités d'accès aux technologies de l'information et de la communication (TIC), tant entre les pays qu'à l'intérieur d'un même pays (entre villes et campagnes, entre riches et pauvres, entre jeunes et personnes âgées).
\end{definition}

\begin{exemplebox}[Un monde connecté de manière très inégale]
Environ \textbf{deux tiers} de la population mondiale utilise Internet aujourd'hui (ordre de grandeur : plus de 5 milliards d'internautes). Mais cette moyenne cache de forts contrastes : plus de 90 \% de la population est connectée en Europe et en Amérique du Nord, contre \textbf{moins de 40 \% en Afrique subsaharienne}. Au Cameroun, le taux de pénétration d'Internet est de l'ordre de 40 à 45 \% (ordre de grandeur) : les grandes villes (Yaoundé, Douala, Bafoussam) sont bien couvertes en 4G, tandis que de nombreuses zones rurales de l'Extrême-Nord, de l'Est ou de l'Adamaoua restent mal desservies. La fracture numérique est donc à la fois \emph{internationale} et \emph{interne}.
\end{exemplebox}

\begin{popfigure}
\begin{center}
\begin{tikzpicture}
\begin{axis}[
  width=0.72\linewidth, height=6.2cm,
  ybar, bar width=16pt,
  ymin=0, ymax=100,
  ylabel={Part de la population utilisant Internet (\%)},
  symbolic x coords={Amérique du Nord, Europe, Asie-Pacifique, Afrique subsah.},
  xtick=data, x tick label style={font=\footnotesize, align=center},
  ytick={0,20,40,60,80,100},
  nodes near coords, nodes near coords style={font=\footnotesize},
  axis line style={draw=popline},
  tick style={draw=popline},
]
\addplot[fill=popblue, draw=popblue!60!black] coordinates {(Amérique du Nord,92) (Europe,89) (Asie-Pacifique,65) (Afrique subsah.,40)};
\end{axis}
\end{tikzpicture}
\end{center}
\legende{Taux d'utilisation d'Internet par grande région du monde (ordres de grandeur, début des années 2020, d'après l'Union internationale des télécommunications). L'écart entre l'Amérique du Nord et l'Afrique subsaharienne illustre la fracture numérique internationale.}
\end{popfigure}

Les causes de la fracture numérique sont cumulatives : coût des équipements et des abonnements, manque d'infrastructures (électricité, fibre), niveau d'éducation, langues dominantes sur le web. Ses conséquences sont graves : les populations non connectées sont exclues des services en ligne, de l'information, de la formation et d'une partie croissante des échanges économiques.

\courssub{La mondialisation des services : délocalisations et centres d'appels}

La révolution numérique a rendu « transportables » des activités autrefois immobiles : c'est la \cle{mondialisation des services}. Un service (répondre au téléphone, traiter des données, écrire un logiciel, gérer des comptes bancaires) peut désormais être produit à des milliers de kilomètres du client.

\begin{definition}[Délocalisation de services]
Une \cle{délocalisation} est le transfert d'une activité économique d'un pays vers un autre, généralement pour profiter de coûts de main-d'œuvre plus faibles. Grâce aux TIC, ce ne sont plus seulement les usines qui sont délocalisées, mais aussi les \textbf{services} : centres d'appels, saisie de données, développement de logiciels, services financiers et comptables.
\end{definition}

\begin{exemplebox}[L'Inde et les Philippines, bureaux du monde]
L'\textbf{Inde} est devenue la capitale mondiale des services informatiques délocalisés : les entreprises de Bangalore (comme Infosys ou Tata Consultancy Services) écrivent des logiciels et gèrent les systèmes d'information de firmes américaines et européennes, grâce à une main-d'œuvre anglophone, qualifiée et moins coûteuse. Les \textbf{Philippines} sont le premier pays mondial pour les centres d'appels en anglais. En Afrique, le Maroc, la Tunisie, le Sénégal et Madagascar accueillent des centres d'appels francophones pour des entreprises françaises. Le Cameroun tente de se positionner sur ce créneau (sous-traitance numérique, « business process outsourcing »), notamment à Douala et Yaoundé.
\end{exemplebox}

\courssub{Le commerce en ligne (e-commerce)}

\textbf{Situation concrète.} À Douala, Amina commande sur son téléphone une paire de chaussures sur Jumia, paie par MTN Mobile Money, et se fait livrer trois jours plus tard par un coursier à moto. Aucun billet de banque n'a changé de main, aucun magasin physique n'a été visité : pourtant, un produit bien réel a traversé la ville — et peut-être la planète.

\begin{definition}[Commerce en ligne]
Le \cle{commerce en ligne} (ou \cle{e-commerce}) est l'achat et la vente de biens et de services via Internet, sans contact physique entre le vendeur et l'acheteur. Il comprend la commande en ligne, le paiement (bancaire ou mobile) et la livraison du bien ou la fourniture du service.
\end{definition}

Le e-commerce mondial est dominé par quelques géants : \textbf{Amazon} (États-Unis), \textbf{Alibaba} (Chine), qui brassent des centaines de milliards de dollars de transactions. Mais l'Afrique connaît son propre essor : \textbf{Jumia}, créée en 2012 et présente dans une dizaine de pays africains dont le Cameroun, est souvent surnommée « l'Amazon africain ». Au Cameroun, le développement du e-commerce repose sur un atout décisif : le \cle{mobile money} (MTN Mobile Money, Orange Money), qui permet de payer sans compte bancaire, dans un pays où une grande partie de la population n'a pas de carte bancaire.

\begin{methode}[Analyser les avantages et les limites du e-commerce au Cameroun]
\begin{enumerate}
\item \textbf{Identifier les acteurs} : consommateur, vendeur, plateforme, opérateur de paiement, livreur.
\item \textbf{Lister les avantages} : choix plus large, prix comparables, gain de temps, accès à des produits absents des marchés locaux, nouveaux débouchés pour les producteurs et artisans.
\item \textbf{Lister les limites} : nécessité d'une connexion et d'un smartphone, risques d'arnaque et de contrefaçon, difficultés de livraison (adresses imprécises, routes dégradées), concurrence déloyale pour les petits commerçants, fiscalité difficile à appliquer.
\item \textbf{Conclure en nuance} : le e-commerce intègre davantage le Cameroun à la mondialisation, mais ses bénéfices restent concentrés dans les villes connectées.
\end{enumerate}
\end{methode}

\begin{attention}[Le e-commerce ne supprime pas les flux physiques]
Erreur fréquente : croire que le commerce en ligne « dématérialise » totalement les échanges. Faux ! Les marchandises commandées en ligne doivent toujours être \textbf{transportées} : conteneurs, ports (Douala, Kribi), entrepôts, camions, livreurs à moto. Le e-commerce ajoute une couche immatérielle (commande, paiement, données) au-dessus des flux physiques étudiés dans la première section ; il ne les remplace pas. La logistique et la livraison du « dernier kilomètre » sont même devenues des enjeux majeurs.
\end{attention}

\begin{exoresolu}[Calculer et commenter une part de marché]
\textbf{Énoncé.} En 2022, on estime qu'environ 5,3 milliards de personnes utilisaient Internet dans le monde, dont environ 600 millions en Afrique (ordres de grandeur). 1) Calcule la part de l'Afrique dans le nombre total d'internautes. 2) La population africaine représente environ 18 \% de la population mondiale. Compare et commente.
\tcblower
\textbf{Solution détaillée.}
1) Part de l'Afrique dans les internautes mondiaux :
\[ \frac{600}{5300} \times 100 \approx 11{,}3\,\%. \]
L'Afrique compte donc environ 11 \% des internautes du monde.
2) Comparaison : l'Afrique représente environ 18 \% de la population mondiale mais seulement 11 \% des internautes. Son « poids numérique » est donc inférieur à son poids démographique : c'est une manifestation chiffrée de la \cle{fracture numérique}. L'écart (18 \% contre 11 \%) montre que l'accès à Internet reste plus limité en Afrique que la moyenne mondiale, ce qui freine la participation du continent aux flux d'information, de services et de commerce en ligne — mais aussi que la marge de croissance y est considérable.
\end{exoresolu}

\courssub{Les flux d'information : un enjeu de pouvoir}

Qui contrôle l'information contrôle une part du pouvoir mondial. Trois acteurs dominent ces flux :

\begin{itemize}
\item les \textbf{médias mondiaux} (chaînes d'information continues comme CNN, BBC, Al Jazeera, France 24 ; agences de presse) qui construisent l'image du monde vue par des milliards de personnes ;
\item les \textbf{réseaux sociaux} (Facebook, WhatsApp, TikTok, X), devenus la première source d'information de nombreux jeunes, y compris au Cameroun, mais aussi des vecteurs de rumeurs et de désinformation ;
\item les \textbf{géants du numérique} (Google, Meta, Amazon, Microsoft, Apple — les « GAFAM »), qui collectent et exploitent les \cle{données} personnelles, nouvelle matière première de l'économie mondiale.
\end{itemize}

\begin{aretenir}[L'essentiel de la section]
\begin{itemize}
\item La révolution numérique (Internet, câbles sous-marins, satellites) permet la circulation instantanée de l'information ; plus de 95 \% du trafic intercontinental passe par les câbles sous-marins (SAT-3, WACS, ACE pour l'Afrique de l'Ouest).
\item La \cle{fracture numérique} oppose pays riches et pays pauvres, mais aussi villes et campagnes au sein du Cameroun.
\item Les services se mondialisent : délocalisations, centres d'appels, logiciels (Inde, Philippines).
\item Le \cle{commerce en ligne} (Amazon, Alibaba, Jumia, mobile money) transforme les échanges, sans supprimer les flux physiques de marchandises.
\item Les flux d'information (médias, réseaux sociaux, données) constituent un enjeu de pouvoir majeur entre États et entreprises.
\end{itemize}
\end{aretenir}

\begin{exercice}[Pour t'entraîner]
1) Définis la fracture numérique et illustre-la par deux exemples camerounais. 2) Construis un croquis simple montrant comment une commande passée sur Jumia à Yaoundé mobilise à la fois des flux immatériels et des flux physiques. 3) Explique pourquoi les câbles sous-marins sont des enjeux stratégiques pour les États.
\end{exercice}

\coursec{Les mouvements financiers et les IDE}

Dans la section précédente, nous avons vu que l'information et les services circulent désormais à la vitesse de la lumière grâce aux câbles sous-marins et à l'Internet : un commerçant de Douala commande ses marchandises à Guangzhou en quelques clics. Mais derrière chaque commande, chaque chantier, chaque transfert d'argent envoyé par un parent installé à Paris ou à Washington, il y a des \cle{flux financiers}. La mondialisation, c'est aussi — et peut-être d'abord — une immense circulation de capitaux. C'est ce que nous allons étudier maintenant, en partant d'une situation bien camerounaise : le port en eau profonde de Kribi, construit avec des capitaux chinois.

\courssub{Les types de flux financiers : fonds bilatéraux et fonds multilatéraux}

\textbf{Intuition.} Quand l'État camerounais veut construire une route ou financer son budget, il peut recevoir de l'argent de l'extérieur selon deux grands circuits : soit directement d'un autre État (la France, la Chine, le Japon), soit d'une organisation internationale regroupant de nombreux États. C'est toute la différence entre le \emph{bilatéral} (deux partenaires) et le \emph{multilatéral} (plusieurs partenaires).

\begin{definition}[Fonds bilatéraux et fonds multilatéraux]
Les \cle{fonds bilatéraux} sont des transferts financiers versés directement d'un État à un autre État : dons, prêts à taux préférentiels, annulations de dette. Ils relèvent de l'\cle{aide publique au développement} (APD) et sont souvent liés à des relations historiques (coopération franco-camerounaise, sino-camerounaise, japonaise).
\par
Les \cle{fonds multilatéraux} proviennent d'organisations internationales regroupant plusieurs États membres :
\begin{itemize}
\item la \cle{Banque mondiale} (et sa filiale IDA, Association internationale de développement), qui finance des projets de développement (routes, éducation, santé) ;
\item le \cle{FMI} (Fonds monétaire international), qui intervient pour stabiliser les balances des paiements et accompagne les programmes d'ajustement structurel ;
\item la \cle{BAD} (Banque africaine de développement), institution panafricaine basée à Abidjan, qui cofinance de grands projets d'infrastructures.
\end{itemize}
\end{definition}

\begin{exemplebox}[Des financements concrets au Cameroun]
La BAD et la Banque mondiale ont cofinancé l'autoroute Yaoundé--Nsimalen et des projets hydroélectriques ; le FMI a conclu avec le Cameroun des accords au titre de la Facilité élargie de crédit (notamment en 2017, pour un montant de l'ordre de 666 millions de dollars — ordre de grandeur à retenir). Côté bilatéral, la Chine a accordé des prêts et des dons pour le palais des congrès de Yaoundé ou le stade de Bafoussam ; le Japon (via la JICA) a financé le second pont sur le Wouri et des rocades urbaines ; la France reste un partenaire historique via l'Agence française de développement (AFD).
\end{exemplebox}

\begin{attention}[Ne pas tout mélanger]
L'aide publique au développement est un \emph{transfert sans contrepartie de propriété} : le pays donneur ne devient pas propriétaire de l'entreprise camerounaise. À l'inverse, un IDE implique une prise de contrôle durable (voir plus loin). Confondre aide et investissement est l'erreur la plus fréquente dans les copies du Baccalauréat.
\end{attention}

\courssub{Les transferts de la diaspora : un fleuve financier discret}

\textbf{Observation.} À Douala comme à Bafoussam, combien de familles vivent grâce à l'argent envoyé par un oncle de Lyon, une sœur de Montréal ou un cousin de Dubaï ? Ces envois, effectués via Western Union, MoneyGram, les banques ou les applications mobiles (Mobile Money, Wave), constituent les \cle{transferts de la diaspora} (ou transferts des migrants).

\begin{definition}[Transferts de la diaspora]
Les \cle{transferts de la diaspora} sont les sommes d'argent envoyées par les migrants vers leur pays d'origine, destinées à la consommation familiale (nourriture, loyers), à l'éducation (frais de scolarité), à la santé et à l'investissement (construction de maisons, création de petites entreprises).
\end{definition}

\begin{exemplebox}[L'importance des transferts pour le Cameroun]
La diaspora camerounaise est estimée à plusieurs millions de personnes (Afrique du Sud, France, États-Unis, Allemagne, Canada, pays du Golfe...). Ses transferts sont \textbf{estimés entre 200 et 400 milliards de FCFA par an} selon les années et la prise en compte des canaux informels (soit environ 1 à 2 \% du PIB — ordres de grandeur prudents, les chiffres officiels de la Banque mondiale ne comptabilisant que les canaux formels, autour de 150 à 200 milliards FCFA). À l'échelle de l'Afrique subsaharienne, ces transferts \textbf{dépassent souvent l'aide publique au développement} reçue. Ils financent massivement la construction immobilière dans les villes camerounaises et la scolarisation des enfants restés au pays.
\end{exemplebox}

\begin{savaistu}[Pourquoi les transferts de la diaspora sont-ils si précieux ?]
Contrairement aux capitaux spéculatifs, les transferts des migrants sont \emph{stables} : même en temps de crise, la diaspora continue d'envoyer de l'argent à sa famille, parfois davantage. Les économistes disent qu'ils sont « contra-cycliques ». C'est pourquoi certains États africains émettent des « bons de la diaspora » pour canaliser cette épargne vers de grands projets nationaux.
\end{savaistu}

\courssub{Les IDE : définition et distinction avec les placements de portefeuille}

\textbf{Intuition.} Quand la compagnie chinoise CHEC (China Harbour Engineering Company) construit et exploite en partie le port de Kribi, elle ne se contente pas de prêter de l'argent : elle s'installe durablement, prend des participations, dirige. C'est un investissement direct. Quand un fonds achète simplement des actions d'une entreprise pour les revendre six mois plus tard, c'est un placement de portefeuille.

\begin{definition}[Investissement Direct Étranger]
Un \cle{IDE} (Investissement Direct Étranger) est un investissement réalisé par une entreprise étrangère pour \textbf{créer, acquérir ou contrôler} une activité durable dans un pays d'accueil. Par convention internationale (OCDE, CNUCED), on parle d'investissement direct lorsque l'investisseur détient au moins \textbf{10 \%} du capital de l'entreprise, seuil censé marquer une volonté de contrôle durable. Il prend trois formes principales : création d'une filiale (\emph{greenfield}), rachat d'une entreprise existante (fusion-acquisition), extension des capacités d'une filiale déjà implantée.
\end{definition}

\begin{propriete}[IDE et placements de portefeuille : la distinction exigible]
\begin{itemize}
\item \textbf{IDE} : investissement productif, durable, avec contrôle de la gestion (usine, mine, plantation, filiale bancaire, infrastructure en concession). Il génère des emplois et des transferts de technologie.
\item \textbf{Placements de portefeuille} : achat d'actions ou d'obligations \emph{sans} volonté de contrôle, motivés par le rendement financier à court terme ; ils sont mobiles et peuvent quitter un pays en quelques heures (« capitaux volatils » ou \emph{hot money}).
\end{itemize}
Cette distinction est exigible au Baccalauréat : elle permet de comprendre pourquoi les IDE sont recherchés par les États (stabilité, emplois, transferts de technologie) alors que les capitaux de portefeuille peuvent déstabiliser une économie en cas de fuite brutale.
\end{propriete}

\courssub{La géographie des IDE : concentration et nouvelles dynamiques}

\textbf{Une géographie très inégale.} Les IDE restent massivement concentrés au sein de la \cle{Triade} (Amérique du Nord, Union européenne, Asie de l'Est) : les pays riches investissent d'abord les uns chez les autres, car ils y trouvent des marchés solvables, des infrastructures de qualité et une sécurité juridique. Les États-Unis et la Chine figurent parmi les premiers destinataires et les premiers investisseurs mondiaux.

\textbf{La montée des IDE vers l'Afrique.} Depuis les années 2000, l'Afrique attire davantage de capitaux, mais de manière très sélective :
\begin{itemize}
\item le \textbf{pétrole} et le \textbf{gaz} (Nigeria, Angola, Mozambique) ;
\item les \textbf{mines} (RDC pour le cobalt/cuivre, Zambie, Guinée pour la bauxite) ;
\item les \textbf{télécommunications} et la banque (MTN, Orange Money, fintechs) ;
\item les \textbf{infrastructures} financées par la Chine dans le cadre de l'initiative « Ceinture et Route » (nouvelles routes de la soie).
\end{itemize}

\begin{exemplebox}[Les IDE chinois au Cameroun]
La Chine est devenue le premier partenaire commercial du Cameroun et un investisseur majeur :
\begin{itemize}
\item le \textbf{port en eau profonde de Kribi} (Région du Sud), construit par China Harbour Engineering Company (CHEC), opérationnel depuis 2015 pour sa première phase (concession 25 ans) ;
\item les \textbf{barrages hydroélectriques} de Memve'ele (Sud, environ 211 MW, mis en service 2019) et de Mekin (Nord, 15 MW), financés en grande partie par Eximbank of China ;
\item de nombreuses routes (ex. : Batchenga-Ntui-Yoko, Kumba-Mamfé) et le palais des congrès de Yaoundé.
\end{itemize}
Ces projets illustrent la stratégie chinoise : sécuriser l'accès aux matières premières (bois, coton, minerais via Kribi) et ouvrir des débouchés à ses entreprises de travaux publics (BTP).
\end{exemplebox}

\begin{popfigure}
\begin{center}
\begin{tikzpicture}[scale=1.0,>=Stealth, font=\small]
% Styles
\tikzstyle{bloc}=[rectangle, rounded corners=3pt, draw, thick, align=center, minimum height=1.8cm, minimum width=3.2cm]
\tikzstyle{flux}=[->, line width=1.8pt, shorten >=2pt, shorten <=2pt]
% Colonnes
% GAUCHE : Investisseur Chinois
\node[bloc, fill=poporangeL, draw=poporange, align=center] (invest) at (-4.5, 2.5){FTN Chinoise\\ (CHEC / Eximbank)\\\textbf{Apporteur de capitaux \& technologie}};
\node[bloc, fill=popgreenL, draw=popgreen, align=center] (construct) at (-4.5, -0.5){Phase Construction\\\textbf{Ouvriers \& matériels chinois + locaux}\\\textit{Création d'emplois temporaires}};
% DROITE : Cameroun / Kribi
\node[bloc, fill=popblueL, draw=popblue, align=center] (port) at (4.5, 2.5){Port en eau profonde\\ de Kribi (Cameroun)\\\textbf{Infrastructure stratégique}};
\node[bloc, fill=poppurpleL, draw=poppurple, align=center] (exploit) at (4.5, -0.5){Phase Exploitation\\\textbf{Conteneurs, export bois/coton/cacao}\\\textit{Revenus portuaires, taxes}};
\node[bloc, fill=popgoldL, draw=popgold, align=center] (retomb) at (4.5, -3.5){Retombées territoriales\\\textbf{Emplois durables, zone industrielle,}\\\textbf{urbanisation de Kribi, fiscalité}};
% FLÈCHES PRINCIPALES
\draw[flux, poporange] (invest.east) -- node[above, align=center] {1. IDE : Prêt Eximbank\\ + Participation CHEC\\(Concession 25 ans)} (port.west);
\draw[flux, popblue] (port.south) -- node[right, align=center]{2. Mise à disposition\\ de l'infrastructure} (exploit.north);
\draw[flux, poppurple] (exploit.south) -- node[right, align=center]{3. Activité économique\\ \& génération de valeur} (retomb.north);
\draw[flux, poporange, dashed] (retomb.west) -- node[below, align=center, xshift=-1cm] {4. Remboursement dette\\ + Revenus concession} ++(-9,0) |- (invest.south);
% FLÈCHES TRANSVERSALES (Construction)
\draw[flux, popgreen] (construct.east) -- node[above, align=center, yshift=0.5cm] {Chantier\\ (Transfert de compétences)} (exploit.west);
% LEGENDES INTERNES
\node[align=center, font=\tiny, text=popmuted] at (-4.5, -2.2) {Main-d'œuvre locale formée};
\node[align=center, font=\tiny, text=popmuted] at (4.5, -5.0) {Effet d'entraînement\\ (logistique, commerce)};
\end{tikzpicture}
\end{center}
\legende{Schéma du circuit d'un IDE chinois dans le port en eau profonde de Kribi : 1. Apport de capitaux (prêt concessionnel Eximbank + participation au capital CHEC) ; 2. Construction (emplois, transfert de compétences) ; 3. Exploitation portuaire (exportations camerounaises) ; 4. Retour financier vers l'investisseur (remboursement, dividendes) et retombées locales durables (emplois, zone industrielle, fiscalité).}
\end{popfigure}

\courssub{Comparer les grands flux financiers vers l'Afrique subsaharienne}

Pour mesurer le poids relatif des différents flux, les géographes comparent les IDE, l'aide publique au développement et les transferts de la diaspora sur un même graphique.

\begin{popfigure}
\begin{center}
\begin{tikzpicture}
\begin{axis}[
  ybar, bar width=0.6cm,
  width=0.95\linewidth, height=6.5cm,
  ymin=0, ymax=100,
  ylabel={Milliards de dollars US (ordres de grandeur annuels récents)},
  xtick={1,2,3},
  xticklabels={IDE, APD, Transferts\\ diaspora},
  x tick label style={font=\small, align=center},
  ymajorgrids=true, grid style={popline!40},
  nodes near coords, nodes near coords style={font=\small\bfseries},
  every axis plot/.append style={fill=popblue!70, draw=popdark}
]
\addplot coordinates {(1,45) (2,35) (3,55)};
\end{axis}
\end{tikzpicture}
\end{center}
\legende{Comparaison des principaux flux financiers entrants nets en Afrique subsaharienne (ordres de grandeur annuels récents, en milliards de dollars US ; sources : Banque mondiale, CNUCED, OCDE — les valeurs exactes varient selon les années). Les transferts de la diaspora dépassent désormais l'aide publique au développement (APD) et rivalisent avec les IDE.}
\end{popfigure}

\begin{methode}[Construire un diagramme en bâtons pour comparer des flux financiers]
\begin{enumerate}
\item \textbf{Choisir une même unité} pour tous les flux (par exemple milliards de dollars US courants) et une même année de référence.
\item \textbf{Ordonner les catégories} sur l'axe horizontal (ex. : IDE, APD, Transferts diaspora) et graduer l'axe vertical à partir de zéro.
\item \textbf{Tracer des bâtons de largeur égale}, de couleur distincte si l'on compare plusieurs années (bâtons juxtaposés), ou de même couleur pour une année unique.
\item \textbf{Légender et dater} : titre explicite, source précise (Banque mondiale, CNUCED, année), unité.
\item \textbf{Conclure} par une phrase de comparaison chiffrée : « En 2022, les transferts de la diaspora vers l'Afrique subsaharienne (≈ 55 Md \$) ont représenté 1,6 fois l'APD (≈ 35 Md \$) ».
\end{enumerate}
\end{methode}

\courssub{Le rôle des FTN et l'instabilité financière mondiale}

\textbf{Les FTN, chefs d'orchestre des flux.} Les \cle{firmes transnationales} (FTN) — TotalEnergies, MTN, Huawei, Bolloré Logistics, CHEC, Dangote Cement — décident de la localisation de leurs investissements selon trois critères majeurs : l'accès aux matières premières, la proximité des marchés de débouché et le coût des facteurs de production (main-d'œuvre, fiscalité, énergie). En fragmentant leurs chaînes de valeur mondiales entre plusieurs pays, elles orientent les IDE et, avec eux, les emplois, les technologies et les devises. Une grande FTN peut ainsi peser plus lourd qu'un État africain moyen : le chiffre d'affaires des plus grandes d'entre elles (Walmart, Apple, Amazon, Sinopec) dépasse le PIB de nombreux pays du continent.

\textbf{La spéculation, revers de la liberté de circulation des capitaux.} Depuis la dérégulation financière des années 1980, les capitaux circulent 24 h/24 entre les places financières majeures (New York, Londres, Tokyo, Hong Kong, Shanghai). Une part croissante de ces mouvements est \emph{spéculative} : achats et ventes rapides de devises, d'actions ou de matières premières pour réaliser des profits immédiats (arbitrages, \emph{carry trade}). Cette finance « de casino » provoque des instabilités majeures : crise asiatique de 1997 (fuite des capitaux), crise mondiale de 2008 (titrisation), envolées spéculatives des prix alimentaires (2008, 2011, 2022) qui aggravent l'insécurité alimentaire dans les pays pauvres importateurs. C'est un complément utile pour nuancer une copie du Bac : la mondialisation financière enrichit les pôles dynamiques, mais elle est structurellement fragile et inégalitaire.

\begin{exoresolu}[Distinguer et hiérarchiser les flux financiers]
\textbf{Énoncé.} Voici quatre flux financiers reçus par le Cameroun : (a) un prêt de la Banque mondiale (IDA) pour un programme d'électrification rurale ; (b) l'argent envoyé par des Camerounais de Paris à leurs familles de Yaoundé via Mobile Money ; (c) la construction d'une cimenterie par un groupe marocain qui en détient 80 \% du capital ; (d) l'achat d'actions d'une banque camerounaise cotée à la BVMAC par un fonds de pension américain, revendues trois mois plus tard. Pour chaque flux, indique sa nature (fonds multilatéral, transfert de la diaspora, IDE, placement de portefeuille) en justifiant, puis explique lequel est le plus stable à long terme.
\tcblower
\textbf{Solution détaillée.}
\begin{itemize}
\item (a) \textbf{Fonds multilatéral} : la Banque mondiale (via l'IDA) est une organisation internationale regroupant de nombreux États membres ; le prêt concessionnel finance un service public (électrification) sans prise de contrôle d'entreprises privées.
\item (b) \textbf{Transfert de la diaspora} : envoi privé de migrants vers leur pays d'origine, destiné à la consommation et à l'investissement familial ; flux contra-cyclique et stable.
\item (c) \textbf{IDE} : le groupe marocain crée une activité productive durable (cimenterie) et détient 80 \% du capital, donc le contrôle effectif (seuil de 10 \% largement dépassé) ; flux stable, créateur d'emplois et de technologie.
\item (d) \textbf{Placement de portefeuille} : achat d'actions cotées sans volonté de contrôle de gestion, horizon de placement court (3 mois), motivation purement spéculative ; capital volatile (\"hot money\").
\end{itemize}
\textbf{Hiérarchie de la stabilité (du plus stable au plus volatile) :}
1. \textbf{IDE (c)} : ancrage territorial, actifs fixes non délocalisables à court terme.
2. \textbf{Transferts diaspora (b)} : motivation affective/familiale, résilients aux chocs.
3. \textbf{Fonds multilatéraux (a)} : engagements pluriannuels, conditionnalité mais prévisibilité.
4. \textbf{Placements de portefeuille (d)} : fuite possible en quelques heures (\"sudden stop\"), risque de crise de change et boursière.
\end{exoresolu}

\begin{attention}[Les trois pièges classiques de la copie]
\begin{itemize}
\item Confondre \cle{IDE} et \cle{aide publique au développement} : l'IDE donne un droit de propriété et de contrôle (actionnariat $\ge$ 10 \%) ; l'APD est un transfert (don/prêt) sans contrepartie de propriété.
\item Croire que les IDE vont d'abord vers les pays pauvres : ils se concentrent dans la Triade (≈ 70 \% du stock mondial) ; l'Afrique ne reçoit qu'une faible part des flux mondiaux (≈ 3 à 5 \%), ciblée sur les secteurs extractifs (pétrole, mines) et les services (télécoms, finance).
\item Oublier de dater et de sourcer les chiffres dans un commentaire de tableau ou de graphique : un flux financier sans année ni source (Banque mondiale, CNUCED, FMI, BEAC) perd toute valeur scientifique.
\end{itemize}
\end{attention}

\begin{aretenir}[L'essentiel de la section]
La mondialisation financière repose sur quatre grands flux : les \cle{fonds bilatéraux} (d'État à État : France, Chine, Japon via AFD, Eximbank, JICA), les \cle{fonds multilatéraux} (Banque mondiale/IDA, FMI, BAD), les \cle{transferts de la diaspora} (estimés 200--400 milliards FCFA/an pour le Cameroun, supérieurs à l'APD en Afrique subsaharienne) et les \cle{IDE}, investissements durables avec contrôle (seuil 10 \% capital), à distinguer des placements de portefeuille volatils. Les IDE restent concentrés dans la Triade, mais la Chine s'impose en Afrique et au Cameroun (port de Kribi, barrages Memve'ele/Mekin, routes). Les FTN orientent ces flux, tandis que la spéculation financière nourrit l'instabilité mondiale.
\end{aretenir}

\coursec{Dossier 3 : Le commerce équitable (activité guidée)}

Dans la section précédente, nous avons vu que la mondialisation est aussi un gigantesque mouvement de capitaux : fonds bilatéraux et multilatéraux, transferts de la diaspora, investissements directs étrangers. Mais cet argent qui circule profite-t-il à tous ? Pour le comprendre, descendons de la sphère financière vers le champ d'un planteur camerounais. Ce dossier t'invite à enquêter sur une injustice économique bien concrète, puis à découvrir une réponse organisée : le \cle{commerce équitable}.

\courssub{Mise en situation : le paradoxe du cacao de Sangmélima}

\textbf{Une scène de la vie quotidienne.}\par
À Sangmélima, dans la région du Sud, un producteur de cacao vend sa récolte de fèves à un acheteur agréé. Pour produire un kilogramme de fèves, il a travaillé plusieurs mois : défrichage, entretien des cacaoyers, récolte des cabosses, écabossage, fermentation, séchage. Pourtant, le prix qu'il reçoit couvre à peine ses coûts de production et les besoins de sa famille. Quelques mois plus tard, ces mêmes fèves, transformées en Europe, se retrouvent dans une tablette de chocolat vendue l'équivalent de plusieurs milliers de francs CFA dans un supermarché de Paris ou de Bruxelles.

\textbf{La question qui fâche.}\par
Comment se fait-il que celui qui produit la matière première essentielle — le cacao — perçoive la plus petite part du prix final ? Où va la valeur ? Qui la capte ? Pour répondre, observons un document statistique : la répartition approximative du prix final d'une tablette de chocolat.

\begin{popfigure}
\begin{center}
\begin{tikzpicture}
\begin{axis}[
    ybar,
    bar width=0.55cm,
    width=\linewidth,
    height=6.5cm,
    ymin=0, ymax=45,
    ylabel={Part du prix final (en \%)},
    symbolic x coords={Producteur, Intermédiaires, Transformateur, Distributeur},
    xtick=data,
    x tick label style={rotate=15, anchor=east, font=\small},
    ytick={0,5,10,15,20,25,30,35,40},
    nodes near coords,
    every node near coord/.append style={font=\small\bfseries, color=popdark},
    axis line style={popline},
    tick style={popline},
]
\addplot[fill=popgreen, draw=popdark] coordinates {(Producteur,6)};
\addplot[fill=popgold, draw=popdark] coordinates {(Intermédiaires,12)};
\addplot[fill=poporange, draw=popdark] coordinates {(Transformateur,40)};
\addplot[fill=popblue, draw=popdark] coordinates {(Distributeur,42)};
\end{axis}
\end{tikzpicture}
\end{center}
\legende{Répartition approximative du prix final d'une tablette de chocolat vendue en Europe (ordres de grandeur, sources : rapports Fairtrade et ONG du cacao, années 2010-2020). Le producteur de fèves ne perçoit souvent que 5 à 7~\% du prix final, tandis que la transformation et la distribution captent l'essentiel de la valeur ajoutée.}
\end{popfigure}

\textbf{Lecture du document.}\par
Le diagramme en bâtons montre une répartition spectaculairement inégale : le producteur reçoit environ 6~\% du prix final, les intermédiaires (acheteurs, transporteurs, exportateurs) environ 12~\%, tandis que le transformateur (chocolatier) et le distributeur (grande distribution) se partagent plus de 80~\% de la valeur. Autrement dit, plus on s'éloigne du champ, plus la part captée est grande. Cette observation est le point de départ de toute la réflexion sur le commerce équitable.

\begin{exemplebox}[Un chiffre qui résume tout]
Pour une tablette de chocolat vendue 1\,500~FCFA en Europe, le producteur de cacao de Sangmélima ou de Kumba perçoit souvent \textbf{moins de 7~\%} du prix final, soit moins de 105~FCFA, alors que sa fève est l'ingrédient indispensable du produit. Sur une année, un planteur camerounais de cacao gagne fréquemment moins que le seuil de pauvreté monétaire, alors que l'industrie mondiale du chocolat pèse des dizaines de milliards de dollars.
\end{exemplebox}

\courssub{Le concept de commerce équitable : définition et origine}

\textbf{L'intuition d'abord.}\par
Face à ce déséquilibre, une idée simple est née : et si, au lieu de laisser le seul jeu du marché fixer les prix, producteurs et acheteurs signaient un \emph{contrat moral et commercial} garantissant un prix juste, stable, versé directement à des producteurs organisés ? C'est exactement le pari du commerce équitable.

\begin{definition}[Le commerce équitable]
Le \cle{commerce équitable} est un \textbf{partenariat commercial} entre des producteurs organisés du Sud et des acheteurs du Nord, fondé sur le dialogue, la transparence et le respect, qui garantit aux producteurs un \textbf{prix rémunérateur stable}, des \textbf{conditions de travail dignes} et un \textbf{développement durable} de leurs communautés. Il vise à corriger les injustices du commerce international conventionnel, et non à faire la charité.
\end{definition}

\textbf{Un peu d'histoire.}\par
Le commerce équitable naît progressivement entre les années 1960 et 1980. Dans les années 1960, des ONG et des mouvements de solidarité internationale (notamment en Europe du Nord) ouvrent des « magasins du monde » vendant des produits artisanaux achetés directement aux producteurs du Sud. Le slogan de l'époque, « Trade, not aid » (« du commerce, pas de l'aide »), résume la philosophie du mouvement. En 1988, aux Pays-Bas, le prêtre néerlandais Frans van der Hoff et l'activiste Nico Roozen créent le premier label de certification, \cle{Max Havelaar}, pour le café mexicain : c'est une révolution, car le produit équitable entre désormais dans les supermarchés ordinaires. Ce label deviendra la base du système international \cle{Fairtrade} (Fairtrade International, ou FLO), qui certifie aujourd'hui le café, le cacao, la banane, le thé, le coton, le sucre et bien d'autres produits.

\begin{savaistu}[Le label Max Havelaar / Fairtrade]
Créé en 1988 aux Pays-Bas, le label \textbf{Max Havelaar} doit son nom au héros d'un roman néerlandais du XIX\textsuperscript{e} siècle qui dénonçait l'exploitation des paysans colonisés. Devenu le cœur du système \textbf{Fairtrade}, il est aujourd'hui présent dans \textbf{plus de 70 pays producteurs} et concerne plus d'un million et demi de producteurs et travailleurs agricoles dans le monde. Le Cameroun fait partie des pays où des coopératives de café, de cacao et de banane ont obtenu des certifications équitables.
\end{savaistu}

\courssub{Le déséquilibre économique mondial : pourquoi les producteurs s'appauvrissent}

\textbf{Premier mécanisme : la baisse et l'instabilité des cours.}\par
Les matières premières agricoles (cacao, café, coton) et minières voient leurs cours fixés sur les bourses mondiales (Londres, New York), loin des champs. Ces cours connaissent de fortes fluctuations et une tendance à la baisse à long terme, en raison de la surproduction, de la spéculation et de la concurrence entre pays exportateurs. Le producteur camerounais, qui ne stocke pas et doit vendre vite, subit ces variations sans aucune protection : une chute de 30~\% du cours mondial du cacao peut réduire son revenu annuel d'autant.

\textbf{Deuxième mécanisme : la détérioration des termes de l'échange.}\par

\begin{definition}[La détérioration des termes de l'échange]
La \cle{détérioration des termes de l'échange} désigne la \textbf{baisse du pouvoir d'achat des exportations de matières premières par rapport aux importations de produits manufacturés}. Concrètement : avec la même quantité de cacao ou de café exportée, un pays du Sud achète de moins en moins de machines, de médicaments ou de véhicules importés, car les prix des produits manufacturés augmentent plus vite que ceux des matières premières.
\end{definition}

Ce mécanisme, théorisé dès les années 1950 par les économistes Raúl Prebisch et Hans Singer, piège les pays exportateurs de produits primaires : plus ils exportent, plus ils doivent exporter pour importer la même quantité de biens indispensables. Le Cameroun, qui exporte du pétrole brut, du cacao, du café et du bois peu transformé, et qui importe des produits manufacturés et des équipements, connaît bien cette contrainte.

\textbf{Troisième mécanisme : la part inégale de la valeur ajoutée.}\par
La \cle{valeur ajoutée} — la richesse créée à chaque étape de transformation — est captée là où se trouvent la transformation, la marque et la distribution, c'est-à-dire au Nord. Le diagramme en bâtons de la mise en situation l'a montré : torréfier, conchier, emballer et vendre du chocolat rapporte infiniment plus que cultiver la fève. Tant que le Cameroun exporte des fèves brutes plutôt que du chocolat fini, il reste en bas de l'échelle de la valeur.

\courssub{Les normes du commerce équitable}

Le commerce équitable n'est pas une vague promesse : il repose sur des \textbf{normes précises}, contrôlées par des organismes de certification indépendants. Les voici.

\begin{enumerate}
\item \textbf{Un prix minimum garanti.} L'acheteur s'engage à payer un prix plancher qui couvre les coûts de production durables, même si le cours mondial s'effondre. Pour le café ou le cacao labellisés Fairtrade, ce prix minimum est fixé par l'organisation internationale et révisé régulièrement.
\item \textbf{Une prime de développement.} En plus du prix, l'acheteur verse une \cle{prime} (par exemple, un montant fixe par tonne) à la coopérative. Cette prime finance des projets collectifs décidés démocratiquement : écoles, forages, centres de santé, routes, équipements de transformation.
\item \textbf{Des relations commerciales durables.} Les contrats sont de long terme, avec un préfinancement partiel possible de la récolte, ce qui libère le producteur de l'endettement auprès des intermédiaires.
\item \textbf{Le respect des droits des travailleurs.} Liberté syndicale, salaires décents, sécurité au travail, absence de discrimination.
\item \textbf{L'interdiction du travail des enfants} et du travail forcé : les enfants doivent aller à l'école, pas aux champs.
\item \textbf{Le respect de l'environnement.} Limitation des pesticides dangereux, protection des sols et de la biodiversité, gestion durable de l'eau ; le commerce équitable va souvent de pair avec la certification biologique.
\end{enumerate}

\begin{attention}[Le commerce équitable n'est pas de la charité]
Erreur fréquente à l'examen : présenter le commerce équitable comme une \emph{aide} ou un \emph{don} aux pays pauvres. Faux ! C'est un \textbf{échange commercial} régi par des \textbf{normes} : le producteur vend un produit de qualité, l'acheteur paie un prix juste. Il n'y a ni mendiant ni bienfaiteur, mais deux partenaires. La formule fondatrice « Trade, not aid » (« du commerce, pas de l'aide ») doit guider ta rédaction.
\end{attention}

\courssub{Les acteurs du commerce équitable}

\begin{itemize}
\item \textbf{Les coopératives de producteurs} : c'est la condition de base. Les petits planteurs doivent s'organiser en coopératives démocratiques pour négocier collectivement, recevoir la prime et être certifiés. Au Cameroun, des coopératives de caféiculteurs de l'Ouest (région de Bafoussam, Dschang) et de cacaoculteurs du Centre et du Sud ont obtenu des certifications équitables ; la filière banane d'exportation (plantations de la plaine de Njombé-Penja) a également été concernée par des démarches de certification sociale et équitable.
\item \textbf{Les labels et organismes de certification} : Fairtrade International (label Max Havelaar), mais aussi d'autres référentiels (Rainforest Alliance, labels biologiques combinés). Ils définissent les normes, contrôlent leur application et autorisent l'apposition du logo sur les produits.
\item \textbf{Les ONG et les plateformes nationales} : elles sensibilisent, accompagnent les coopératives et défendent le mouvement.
\item \textbf{Les consommateurs engagés} : en acceptant de payer un peu plus cher un café ou un chocolat labellisé, le consommateur du Nord (et de plus en plus du Sud) devient un acteur du développement.
\end{itemize}

\begin{popfigure}
\begin{center}
\begin{tikzpicture}[
    font=\small,
    box/.style={draw=popdark, rounded corners=2pt, align=center, minimum height=0.9cm, text width=2.6cm},
    flux/.style={-{Stealth[length=2.5mm]}, very thick},
]
% Titre circuit conventionnel
\node[font=\small\bfseries, popdark] at (3.2,4.6) {Circuit conventionnel};
\node[box, fill=popgreenL, align=center] (p1) at (0.6,3.4){Producteur\\(café, Ouest Cameroun)};
\node[box, fill=popgoldL, align=center] (i1) at (3.2,3.4){Acheteurs et\\intermédiaires multiples};
\node[box, fill=poporangeL, align=center] (e1) at (5.8,3.4){Exportateur\\(port de Douala)};
\node[box, fill=popblueL, align=center] (t1) at (8.4,3.4){Torréfacteur\\et distributeur (Nord)};
\draw[flux, popmuted] (p1) -- (i1);
\draw[flux, popmuted] (i1) -- (e1);
\draw[flux, popmuted] (e1) -- (t1);
\node[align=center, font=\footnotesize, popmuted, text width=8cm] at (4.5,2.5) {Chaîne longue : chaque intermédiaire prélève une marge ;\\le producteur subit le prix imposé et l'instabilité des cours.};

% Titre circuit équitable
\node[font=\small\bfseries, popdark] at (3.2,1.4) {Circuit équitable};
\node[box, fill=popgreen] (p2) at (0.6,0.2) {\textbf{Coopérative} de producteurs};
\node[box, fill=poporange, align=center] (e2) at (4.5,0.2){Acheteur engagé\\(contrat long terme)};
\node[box, fill=popblue, align=center] (t2) at (8.4,0.2){Consommateur\\engagé};
\draw[flux, popgreen!70!black] (p2) -- (e2) node[midway, above, font=\footnotesize, popdark]{prix minimum garanti + prime};
\draw[flux, popgreen!70!black] (e2) -- (t2) node[midway, above, font=\footnotesize, popdark]{produit labellisé};
\draw[flux, poppurple, dashed] (e2.south) .. controls (4.5,-1.1) and (2.5,-1.2) .. (p2.south) node[midway, below, font=\footnotesize, popdark]{prime de développement : écoles, forages, santé};
\end{tikzpicture}
\end{center}
\legende{Schéma comparatif des circuits conventionnel et équitable du café camerounais. En haut, la chaîne longue multiplie les intermédiaires et écrase le prix payé au producteur. En bas, le circuit équitable raccourcit la chaîne, garantit un prix minimum et reverse une prime collective à la coopérative.}
\end{popfigure}

\courssub{Les critiques du commerce équitable}

Un bon géographe ne se contente pas d'admirer : il évalue. Le commerce équitable fait l'objet de critiques sérieuses qu'il faut connaître pour le débat.

\begin{enumerate}
\item \textbf{Une part de marché marginale.} Les produits équitables représentent une part très faible du commerce mondial (de l'ordre de quelques pourcents, souvent moins de 2~\% pour le café ou le cacao selon les marchés). L'immense majorité des producteurs reste dans le circuit conventionnel.
\item \textbf{Le coût de la certification.} Audits, contrôles, frais administratifs : obtenir et conserver un label coûte cher aux coopératives, parfois plusieurs milliers d'euros par an, ce qui exclut les plus petites et les plus pauvres.
\item \textbf{Le risque d'effet de vitrine (le \emph{fairwashing}).} Certaines grandes entreprises labellisent une infime partie de leurs produits pour soigner leur image, tout en maintenant l'essentiel de leurs achats dans le circuit conventionnel : c'est une stratégie de communication plus qu'un engagement réel.
\item \textbf{La dépendance aux consommateurs du Nord.} Le système repose sur la bonne volonté des consommateurs des pays riches ; en cas de crise économique ou de désintérêt, la demande de produits équitables, plus chers, peut chuter.
\item \textbf{Une réponse partielle.} Le commerce équitable corrige les symptômes sans transformer les structures : il ne remet pas en cause la spéculation sur les cours, ni l'absence de transformation locale des matières premières au Cameroun.
\end{enumerate}

\courssub{Travail guidé : analyser un tableau statistique, puis débattre}

\begin{methode}[Analyser un tableau statistique en cinq étapes]
\begin{enumerate}
\item \textbf{Lire le titre} : que mesure le tableau ? (ici, des prix du cacao).
\item \textbf{Lire les unités et les dates} : en quelle monnaie, pour quelle période ?
\item \textbf{Repérer les valeurs extrêmes} : la plus forte, la plus faible.
\item \textbf{Calculer des rapports} : écarts absolus, rapports (combien de fois plus ?), pourcentages.
\item \textbf{Conclure en une phrase} qui répond au problème posé, sans recopier le tableau.
\end{enumerate}
\end{methode}

\begin{exoresolu}[Prix conventionnel et prix équitable du cacao]
\textbf{Énoncé.} Voici un tableau simplifié (ordres de grandeur réalistes) des prix perçus pour une tonne de fèves de cacao :

\begin{center}
\begin{tabularx}{\linewidth}{|l|X|X|}\hline
\rowcolor{popblueL}
\textbf{Situation} & \textbf{Prix perçu par la coopérative (USD/t)} & \textbf{Compléments} \\ \hline
Marché conventionnel (cours mondial bas) & 1\,800 & aucun ; prix imposé, instable \\ \hline
Commerce équitable (prix minimum Fairtrade) & 2\,400 & + prime d'environ 240 USD/t pour des projets collectifs \\ \hline
\end{tabularx}
\end{center}

1. Calcule l'écart de prix entre les deux situations, en valeur absolue puis en pourcentage par rapport au prix conventionnel.
2. Calcule le gain total par tonne pour une coopérative équitable (prix + prime).
3. Conclus en une phrase sur l'intérêt du commerce équitable pour le producteur.
\tcblower
\textbf{Solution détaillée.}
\begin{enumerate}
\item Écart absolu : $2\,400 - 1\,800 = 600$ USD/t. En pourcentage : $\dfrac{600}{1\,800} \times 100 \approx 33~\%$. Le prix équitable est donc supérieur d'environ un tiers au prix conventionnel.
\item Gain total : $2\,400 + 240 = 2\,640$ USD/t, soit un gain de $2\,640 - 1\,800 = 840$ USD/t, c'est-à-dire environ $47~\%$ de plus que le circuit conventionnel.
\item Conclusion : pour une tonne de cacao, le commerce équitable rapporte à la coopérative près de la moitié de revenu en plus, dont une prime collective qui finance écoles, forages et centres de santé : il améliore à la fois le revenu individuel et le développement du village.
\end{enumerate}
\end{exoresolu}

\textbf{Débat argumenté (activité de classe).}\par
Par groupes, préparez puis soutenez l'une des deux thèses suivantes, en mobilisant au moins trois arguments tirés du dossier : « Le commerce équitable est une véritable solution à la pauvreté des producteurs camerounais » contre « Le commerce équitable est une goutte d'eau qui ne change pas les règles du commerce mondial ». Chaque groupe doit anticiper et répondre à une objection de l'autre camp. Le jury (le reste de la classe) évalue la solidité des arguments, la précision des chiffres et la qualité de l'expression.

\begin{aretenir}[L'essentiel du dossier 3]
\begin{itemize}
\item Le \cle{commerce équitable} est un partenariat commercial garantissant un prix rémunérateur stable, des conditions de travail dignes et un développement durable ; il naît des années 1960-1980, avec le label Max Havelaar/Fairtrade créé en 1988.
\item Il répond à un déséquilibre mondial : baisse des cours, \cle{détérioration des termes de l'échange}, part inégale de la valeur ajoutée (le producteur de cacao capte souvent moins de 7~\% du prix final du chocolat).
\item Ses normes : prix minimum garanti, prime de développement, contrats durables, droits des travailleurs, interdiction du travail des enfants, respect de l'environnement.
\item Ses limites : faible part du marché, coût de la certification, \emph{fairwashing}, dépendance au consommateur du Nord.
\end{itemize}
\end{aretenir}

\begin{exercice}[À toi de jouer]
1. Définis en deux phrases la détérioration des termes de l'échange et explique pourquoi elle pénalise le Cameroun.
2. Rédige un paragraphe argumenté (10 à 15 lignes) répondant à la question : « Le commerce équitable peut-il, à lui seul, sortir les producteurs camerounais de la pauvreté ? »
3. Sur le schéma comparatif des circuits du café, identifie deux différences fondamentales entre le circuit conventionnel et le circuit équitable.
\end{exercice}

\begin{corrige}[Exercice : commerce équitable]
1. La détérioration des termes de l'échange est la baisse du pouvoir d'achat des exportations de matières premières par rapport aux importations de produits manufacturés : les prix des premières stagnent ou baissent tandis que ceux des seconds augmentent. Le Cameroun, qui exporte du cacao, du café et du pétrole bruts et importe des machines et des produits manufacturés, doit exporter de plus en plus pour importer la même quantité de biens, ce qui appauvrit l'État et les producteurs.
2. Paragraphe attendu : le commerce équitable améliore réellement le revenu des producteurs certifiés (prix minimum, prime, stabilité), mais il ne peut suffire : il ne touche qu'une faible part du marché, il coûte cher aux coopératives, il dépend des consommateurs du Nord et ne règle ni la spéculation sur les cours ni l'absence de transformation locale. La solution complète suppose aussi l'industrialisation (transformation du cacao sur place), l'organisation des producteurs et des politiques publiques de soutien.
3. Deux différences : la chaîne équitable est courte (coopérative, acheteur engagé, consommateur) alors que la chaîne conventionnelle multiplie les intermédiaires ; le circuit équitable garantit un prix minimum et reverse une prime collective, alors que le circuit conventionnel impose au producteur un prix instable sans retour pour la communauté.
\end{corrige}

\coursec{Dossier 4 : Les problèmes liés à la mondialisation (activité guidée)}

Le Dossier 3 nous a montré que le commerce équitable tente de corriger certains déséquilibres du commerce mondial. Mais la mondialisation ne produit pas seulement des inégalités : elle engendre aussi de véritables \cle{problèmes} — drames humains, criminalité transnationale, conflits, marginalisation. Ce dossier t'invite à agir en géographe : observer des documents, les classer, puis rédiger un paragraphe argumenté. La consigne officielle est claire : \emph{inventorier, classer et interpréter des documents}.

\courssub{Situation-problème : que voit-on dans la presse ?}

Ouvre un journal camerounais ou un site d'information : « Une embarcation de migrants chavire au large de la Libye », « Un réseau de cyberescrocs démantelé à Douala », « Des groupes armés pillent l'or et le diamant en Centrafrique », « Les pays pauvres absents du numérique mondial ». Ces faits divers, en apparence isolés, ont un point commun : ils sont tous rendus possibles — ou aggravés — par la \cle{mondialisation}, c'est-à-dire par la libre circulation accrue des personnes, des biens, des capitaux et de l'information. Notre tâche : comprendre chaque problème, ses causes, ses manifestations, ses conséquences, et les réponses possibles.

\courssub{L'immigration clandestine : le rêve qui tourne au drame}

\textbf{Intuition.} Dans beaucoup de quartiers de Yaoundé, de Douala ou de Bafoussam, des jeunes parlent de « tomber » (partir) vers l'Europe, quitte à traverser le désert et la mer sur des embarcations de fortune. Pourquoi risquer sa vie ? C'est toute la question de l'immigration clandestine.

\begin{definition}[L'immigration clandestine]
L'\cle{immigration clandestine} (ou migration irrégulière) est l'entrée ou le séjour d'une personne sur le territoire d'un État \textbf{sans autorisation légale} (sans visa, sans titre de séjour), ou l'entrée par des voies non contrôlées.
\end{definition}

\textbf{Les causes.} Elles sont de deux types, et les géographes parlent de facteurs de \emph{répulsion} (\emph{push}) et d'\emph{attraction} (\emph{pull}) :
\begin{itemize}
\item la \textbf{pauvreté} et le \textbf{chômage} des jeunes dans les pays de départ ;
\item les \textbf{conflits armés} et l'insécurité (Boko Haram dans l'Extrême-Nord du Cameroun, crises au Sahel, guerres civiles) ;
\item les effets du \textbf{changement climatique} (sécheresses récurrentes au Sahel, dégradation des sols) ;
\item l'attraction des pays riches : emplois, salaires élevés, image de l'« eldorado » européen véhiculée par internet et par les réussites visibles de la diaspora.
\end{itemize}

\textbf{Les routes et les drames.} Trois grandes routes relient l'Afrique à l'Europe : la route de la \textbf{Méditerranée centrale} (Libye vers l'Italie), la route de la \textbf{Méditerranée occidentale} (Maroc vers l'Espagne) et la route de l'\textbf{Atlantique} (côtes ouest-africaines vers les îles Canaries). Toutes passent d'abord par le \textbf{Sahara}, où les migrants sont à la merci des passeurs, de la soif et des violences.

\begin{popfigure}
\includegraphics[width=\linewidth,height=9.5cm,keepaspectratio]{N12/routes_migr.png}
\legende{Les routes de l'immigration africaine vers l'Europe (traits rouges, carte simplifiée en français ; les trois routes ne sont pas numérotées). On y lit : les axes qui convergent vers Agadès (Niger), grande ville-étape du Sahara ; de là, la piste vers la Libye puis Tripoli et Malte (Méditerranée centrale) ; les pistes vers l'Algérie et le Maroc (Ceuta, Melilla) ; et la route côtière de l'Atlantique de Dakar vers le Maroc et les îles Canaries. Le Cameroun, en bas de la carte, est relié à ce réseau par le Nigeria et le Tchad. \\ Source : Wikimedia Commons, « Carte des routes d'immigration africaine vers l'Europe », historicair, Manll (traduction), CC BY-SA 3.0.}
\end{popfigure}

\begin{savaistu}[La route la plus meurtrière du monde]
La \textbf{Méditerranée} est la route migratoire la plus meurtrière du monde : selon l'Organisation internationale pour les migrations (OIM, projet \emph{Missing Migrants}), plus de \textbf{28 000 migrants y ont perdu la vie depuis 2014} (ordre de grandeur à manier avec prudence car beaucoup de naufrages « invisibles » ne sont pas signalés). Au Cameroun, le phénomène des jeunes candidats au départ — parfois surnommés les « aventuriers » — illustre le désespoir d'une jeunesse souvent diplômée mais sans emploi.
\end{savaistu}

\courssub{La cybercriminalité : le crime sans frontières}

\textbf{Intuition.} Un message promettant un gain à un concours, un faux site bancaire, un compte de réseau social piraté : qui n'a jamais reçu une de ces tentatives d'escroquerie ? La mondialisation de l'information a aussi mondialisé le crime.

\begin{definition}[La cybercriminalité]
La \cle{cybercriminalité} regroupe toutes les \textbf{infractions commises au moyen des réseaux informatiques} : escroqueries en ligne (hameçonnage, \emph{phishing}), piratage de systèmes et de comptes, vol de données personnelles et bancaires, chantage numérique (\emph{rançongiciels}), \textbf{blanchiment} d'argent (recyclage de fonds d'origine criminelle pour leur donner une apparence légale).
\end{definition}

\textbf{Manifestations.} Au Cameroun, le phénomène de la \cle{feymania} (argot local dérivé de « feyman », arnaqueur) désigne les escroqueries par internet ou par téléphone (faux héritages, fausses annonces de vente de véhicules, arnaques aux transferts d'argent type \emph{Mobile Money}). À l'échelle mondiale, le coût de la cybercriminalité est estimé à \textbf{plusieurs milliers de milliards de dollars par an} (les études de référence, comme celles de \emph{Cybersecurity Ventures}, projettent un coût mondial de \num{10500} milliards de dollars en 2025 ; ordre de grandeur à manier avec prudence car les estimations varient selon les méthodologies). La cybercriminalité prospère parce que les réseaux ignorent les frontières, alors que les polices et les législations restent largement nationales : un escroc peut opérer depuis un continent contre des victimes d'un autre.

\courssub{Les conflits armés : une économie de guerre mondialisée}

Les conflits armés contemporains sont profondément liés à la mondialisation, par trois mécanismes :
\begin{itemize}
\item le \textbf{trafic d'armes} : les réseaux transnationaux acheminent des armes légères et de petit calibre vers les zones de guerre, alimentant les rebellions (Sahel, Centrafrique, bassin du lac Tchad, Nord-Ouest et Sud-Ouest du Cameroun) ;
\item le \textbf{pillage des ressources} : certains minerais extraits en zone de conflit financent les belligérants — on parle de \cle{minerais de sang} (ou « minerais de conflit »), comme le coltan, la cassitérite (étain), l'or ou les diamants, qui entrent ensuite dans les filières mondiales (industrie électronique, bijouterie, automobile) ;
\item le \textbf{financement transnational} des groupes armés : rançons (piraterie, enlèvements), trafics de drogue et d'êtres humains, transferts d'argent par des réseaux informels (\emph{hawala}) ou détournés via des paradis fiscaux.
\end{itemize}
Ainsi, un téléphone portable acheté à Douala peut contenir, à l'origine, un minerai ayant financé une guerre à des milliers de kilomètres : la mondialisation relie le consommateur final au conflit.

\courssub{La marginalisation des pays pauvres}

\begin{exemplebox}[Les PMA : présents en hommes, absents des échanges]
Les \cle{PMA} (Pays les Moins Avancés), une quarantaine d'États (dont 33 en Afrique) regroupant environ \textbf{1,1 milliard d'habitants} (ordre de grandeur 2023), ne représentent qu'\textbf{environ 1 \% du commerce mondial de marchandises} (données OMC/UNCTAD, part stable depuis des décennies). Ils attirent aussi une part très faible des \textbf{IDE} (Investissements Directs Étrangers), concentrés à plus de 70 \% sur les pays développés et les grandes économies émergentes (Chine, Inde, Brésil). Enfin, la \textbf{fracture numérique} les prive d'une partie des flux d'information et de services : faible débit internet, coût élevé de la connexion par rapport au revenu, retard d'équipement scolaire et administratif.
\end{exemplebox}

Cette marginalisation s'auto-entretient : pauvres, ces pays n'exportent que des matières premières peu transformées (pétrole, cacao, bois, coton) à faible valeur ajoutée ; peu intégrés aux chaînes de valeur mondiales, ils restent pauvres et vulnérables aux chocs de prix. C'est le cercle vicieux de la marginalisation.

\courssub{Autres problèmes : le complément utile au Bac}

\begin{itemize}
\item L'\textbf{évasion fiscale} et l'\textbf{optimisation agressive} : les firmes multinationales déplacent artificiellement leurs profits vers des paradis fiscaux (via les prix de transfert), privant les États (surtout les plus pauvres) de recettes fiscales vitales pour l'éducation, la santé et les infrastructures.
\item Les \textbf{délocalisations} et le \textbf{dumping social/environnemental} : les entreprises s'installent là où les salaires, les protections sociales et les normes environnementales sont les plus faibles, créant une concurrence déloyale défavorable aux travailleurs et aux écosystèmes locaux.
\item L'\textbf{uniformisation culturelle} : la domination de quelques industries culturelles (cinéma, musique, plateformes de streaming, réseaux sociaux) menace la diversité des langues, des cultures locales et des savoirs traditionnels.
\item Les \textbf{risques sanitaires mondiaux} : la circulation rapide des personnes (transport aérien) accélère la diffusion des épidémies (Ebola, Covid-19, Mpox), transformant une crise locale en pandémie mondiale en quelques semaines.
\end{itemize}

\courssub{Les réponses : réguler plutôt que subir}

Face à ces problèmes, plusieurs réponses existent, à des échelles variées :
\begin{itemize}
\item la \textbf{régulation internationale} : processus de Kimberley (diamants), directive européenne et loi américaine Dodd-Frank (minerais de conflit : étain, tantale, tungstène, or), conventions de l'OCDE et du G20 sur l'érosion de la base d'imposition (BEPS), lutte contre le blanchiment (GAFI) ;
\item la \textbf{coopération policière et judiciaire} : \textbf{Interpol} (Organisation internationale de police criminelle, 196 pays membres) relie les polices pour traquer les cybercriminels, les trafiquants d'êtres humains et de drogue ; mandats d'arrêt internationaux, extradition ;
\item l'\textbf{intégration régionale} : des organisations comme la \textbf{CEEAC} (Communauté Économique des États de l'Afrique Centrale) ou la \textbf{CEMAC} (Communauté Économique et Monétaire de l'Afrique Centrale) coordonnent la gestion des frontières, la libre circulation des personnes (dans la théorie) et la lutte contre le terrorisme (Force multinationale mixte Bassin du Lac Tchad) ;
\item la \textbf{labellisation et la valorisation des produits locaux} : valoriser les produits camerounais à forte identité (café arabica des hauts plateaux, cacao, poivre de Penja \emph{IGP}, miel de l'Adamaoua, coton) par des labels de qualité (IGP, AOP, Bio, Commerce équitable) permet de mieux insérer les producteurs dans la mondialisation — c'est le lien direct avec le Dossier 3.
\end{itemize}

\begin{attention}[Le piège du procès à charge]
Au Bac, ne présente \textbf{jamais} la mondialisation comme uniquement négative. Le correcteur attend une analyse \textbf{nuancée} : la mondialisation crée des opportunités (accès aux marchés, investissements, information, transferts de la diaspora qui financent des milliers de familles camerounaises) \emph{et} des problèmes (ceux de ce dossier). Une copie qui ne voit que les drames perd des points ; une copie qui hiérarchise, nuance et propose des pistes de régulation en gagne.
\end{attention}

\courssub{Travail guidé : classer les documents, puis argumenter}

\begin{methode}[Classer des documents par type de problème]
\begin{enumerate}
\item \textbf{Lis} chaque document (photographie, extrait de presse, tableau de données, carte) et souligne les mots-clés : « naufrage », « piratage », « milice », « 1 \% du commerce », « rançongiciel ».
\item \textbf{Identifie} le problème principal correspondant : immigration clandestine, cybercriminalité, conflit armé, marginalisation.
\item \textbf{Classe} les documents dans un tableau à quatre colonnes, en notant pour chacun la cause et la conséquence principales.
\item \textbf{Vérifie} qu'aucun document ne relève de deux catégories sans justification (ex. : le trafic de migrants relève à la fois de l'immigration clandestine et de la criminalité transnationale organisée ; le pillage de l'or en zone de guerre relève du conflit armé et de la marginalisation économique : dis-le explicitement !).
\end{enumerate}
\end{methode}

Voici le tableau-synthèse à compléter pendant l'activité (tu peux le reproduire sur ta copie) :

\begin{center}
\begin{tabularx}{\linewidth}{|l|X|X|X|X|}
\hline
\rowcolor{popblueL} \textbf{Problème} & \textbf{Causes principales} & \textbf{Manifestations concrètes} & \textbf{Conséquences majeures} & \textbf{Pistes de solutions / Régulation} \\
\hline
Immigration clandestine & Pauvreté, chômage, conflits, climat & Traversées Sahara + Méditerranée (3 routes), passeurs & Morts en mer/désert (OIM : 28 000+ depuis 2014), trafic d'êtres humains, diasporas fragilisées & Coopération UE-Afrique, développement pays de départ, voies légales, protection réfugiés \\
\hline
Cybercriminalité & Mondialisation réseaux, anonymat, impunité & Feymania (Cameroun), phishing, rançongiciels, blanchiment & Pertes \num{>10000} Md\$ /an (monde), perte confiance numérique, victimes ruinées & Interpol, législations nationales (loi camerounaise 2010), sensibilisation, cybersécurité \\
\hline
Conflits armés & Trafics d'armes, convoitise ressources, financement transnational & Minerais de sang (coltan, or), groupes armés, enfants soldats & Insécurité, déplacés internes/réfugiés, destruction économie, contamination filières mondiales & Traçabilité minerais (Kimberley, Dodd-Frank), embargo ONU, DDR (Désarmement-Démobilisation-Réinsertion) \\
\hline
Marginalisation des PMA & Faible industrialisation, fracture numérique, termes de l'échange défavorables & $\approx$ 1 \% commerce mondial, IDE négligeables, export matières premières brutes & Pauvreté persistante, dépendance aide, départs migratoires, instabilité politique & IDE ciblés, commerce équitable, labellisation (IGP), transfert technologie, annulation dette \\
\hline
\end{tabularx}
\end{center}

\begin{methode}[Rédiger un paragraphe argumenté en géographie (niveau Bac)]
\begin{enumerate}
\item \textbf{Annonce} l'idée principale dans une phrase claire (la phrase d'introduction du paragraphe). Elle répond directement à la question.
\item \textbf{Prouve}-la par un exemple \textbf{chiffré} (donnée statistique, pourcentage, montant) \textbf{ET/OU localisé} (un lieu précis, une route, une région, une organisation).
\item \textbf{Conclus} en montrant la portée de l'idée (mécanisme, conséquence) ou en ouvrant vers l'idée suivante (transition).
\end{enumerate}
Un bon paragraphe au Bac fait 4 à 6 lignes manuscrites : ni une simple liste de faits, ni un discours général sans preuve géographique.
\end{methode}

\begin{exoresolu}[Rédiger un paragraphe argumenté sur la marginalisation des pays pauvres]
\textbf{Énoncé.} À partir des documents étudiés, rédige un paragraphe argumenté (4 à 6 lignes) montrant que la mondialisation marginalise les pays les plus pauvres.
\tcblower
\textbf{Solution détaillée.} \emph{Annonce} : Loin de profiter à tous de manière équitable, la mondialisation marginalise structurellement les pays les plus pauvres, qui restent confinés à la périphérie des grands flux d'échanges et de capitaux. \emph{Preuve chiffrée} : les Pays les Moins Avancés (PMA), qui rassemblent pourtant plus de 1 milliard d'êtres humains, ne pèsent qu'environ 1 \% du commerce mondial de marchandises et ne captent qu'une part infime des Investissements Directs Étrangers (IDE), massivement concentrés sur la Triade et les grands émergents. \emph{Preuve localisée} : en Afrique centrale, l'essentiel des exportations (pétrole au Tchad et en Guinée équatoriale, bois et cacao au Cameroun, coton en RCA) reste constitué de matières premières brutes, peu transformées et donc peu rémunératrices, tandis que la fracture numérique exclut ces territoires de l'économie de la connaissance. \emph{Conclusion} : enfermés dans ce cercle vicieux de la non-diversification et de la dépendance, ces pays subissent la mondialisation plus qu'ils n'en bénéficient, ce qui nourrit à son tour l'instabilité et les départs migratoires clandestins.
\end{exoresolu}

\begin{experience}[TP : Construire un croquis synthétique « La mondialisation : flux et problèmes »]
\textbf{Objectif.} Réaliser un croquis légendé sur une carte du monde (ou de l'Afrique) qui synthétise le fonctionnement de la mondialisation (flux) et ses problèmes (Dossiers 3 \& 4).
\textbf{Matériel.} Fond de carte monde (projection Robinson ou Mercator) ou Afrique (A4), crayons de couleur, feutres fins, règle, gomme.
\textbf{Consignes.}
\begin{enumerate}
\item \textbf{Fond de carte :} Trace le contour des continents (monde) ou du continent africain. Repère et nomme : Triade (USA, UE, Asie de l'Est), PMA (zone Afrique subsaharienne), grands ports (Shanghai, Rotterdam, Los Angeles, Douala, Lagos, Mombasa).
\item \textbf{Flux de marchandises (mondialisation qui fonctionne) :} Dessine 3 flèches principales (traits pleins, largeur $\propto$ volume) : Asie $\to$ USA/UE (produits manufacturés) ; Afrique/Amérique du Sud $\to$ Monde (matières premières) ; UE/USA $\to$ Afrique (produits manufacturés, IDE). Légende : \emph{Flux commerciaux dominants}.
\item \textbf{Flux financiers et IDE :} Flèches pointillées (capitaux) : Triade $\to$ Triade (majoritaire) ; Triade $\to$ Émergents ; Diaspora $\to$ Pays d'origine (transferts). Légende : \emph{Flux financiers / IDE / Transferts diaspora}.
\item \textbf{Problèmes (Dossier 4) :} Utilise des \textbf{symboles ponctuels} (pas des flèches) :
\begin{itemize}
\item Zone dangereuse sur Méditerranée (routes 1, 2) et Atlantique (route 3) : \emph{Immigration clandestine / Morts}.
\item Cyberattaques sur hubs numériques (Lagos, Douala, Nairobi, Europe, USA) : \emph{Cybercriminalité}.
\item Pictogrammes « conflit » et « minerais » sur Est RDC, Centrafrique, Sahel : \emph{Conflits / Minerais de sang}.
\item $\downarrow$ (flèche descendante) sur zone PMA (Afrique subsaharienne, Haïti, Afghanistan, etc.) : \emph{Marginalisation (1 \% commerce)}.
\end{itemize}
\item \textbf{Réponses (Dossiers 3 \& 4) :} Symboles positifs : Pictogramme « environnement » sur zones labellisées (Penja, Kilimandjaro, Café Arabica) : \emph{Commerce équitable / Labels} ; (pictogramme « poignée de main ») sur sièges organisations (Lyon/Interpol, Bruxelles/UE, Yaoundé/CEMAC, New York/ONU) : \emph{Régulation / Coopération}.
\item \textbf{Légende organisée :} Titre précis, 3 colonnes : \textbf{Flux de la mondialisation} / \textbf{Problèmes générés} / \textbf{Réponses et régulation}. Échelle, Nord, Auteur, Date, Source.
\end{enumerate}
\textbf{Critères de réussite.} Lisibilité (hiérarchie traits/pointillés/symboles), légende structurée en 3 parties, localisation précise (noms de lieux), titre complet.
\end{experience}

\begin{exercice}[À ton tour — Entraînement Bac]
\begin{enumerate}
\item Classe les quatre documents suivants dans le tableau-synthèse (colonne « Problème ») en justifiant d'une phrase chacun :
\begin{itemize}
\item (a) Photo d'une embarcation de migrants surchargée secourue au large de Lampedusa (Italie).
\item (b) Article de presse : « Douala : démantèlement d'un réseau de "feymen" spécialisés dans l'arnaque au faux visa Schengen ».
\item (c) Tableau statistique OMC 2023 : « Part des PMA dans les exportations mondiales de marchandises : 0,9 \% ».
\item (d) Reportage vidéo : « Centrafrique : l'or du sang finance les groupes armés à Bambari ».
\end{itemize}
\item Rédige un paragraphe argumenté (4 à 6 lignes) répondant à la question : « La mondialisation est-elle seulement une chance pour les pays en développement ? » (Pense à la nuance exigée au Bac : thèse / antithèse / synthèse dans un seul paragraphe dense).
\end{enumerate}
\end{exercice}

\begin{corrige}[Exercice : classification et argumentation]
\begin{enumerate}
\item (a) \textbf{Immigration clandestine} — Route de la Méditerranée centrale (Libye/Tunisie $\to$ Italie/Lampedusa). Justification : embarcations de fortune, migrants sans visa, drame humanitaire en mer.
\item (b) \textbf{Cybercriminalité} — Phénomène de la \emph{feymania} camerounaise (arnaque au visa, usurpation d'identité). Justification : usage d'internet/téléphone pour escroquer, réseaux transnationaux, impunité relative.
\item (c) \textbf{Marginalisation des pays pauvres (PMA)}. Justification : donnée chiffrée officielle (0,9 \% $\approx$ 1 \%), preuve quantitative de l'exclusion des flux commerciaux mondiaux.
\item (d) \textbf{Conflits armés} — Pillage des ressources / Minerais de sang (or). Justification : exploitation artisanale contrôlée par groupes armés, financement de la guerre, entrée dans filière mondiale (joaillerie, banque centrale).
\item \textbf{Paragraphe type (nuancé) :} \emph{Annonce} : La mondialisation est un processus ambivalent pour les pays en développement : source d'opportunités réelles, elle génère simultanément des vulnérabilités structurelles. \emph{Preuves chance} : elle ouvre l'accès aux marchés mondiaux (export cacao, bois, coton camerounais), attire des IDE (ex. : port de Kribi, barrage de Nachtigal, usines de transformation), permet les transferts massifs de la diaspora (plus de \num{300} milliards FCFA/an pour le Cameroun, première source de devises devant l'export) et diffuse les connaissances (télémédecine, formation en ligne). \emph{Preuves problèmes} : elle expose à la volatilité des cours (choc cacao 2024), favorise l'immigration clandestine meurtrière (routes sahéliennes), propage la cybercriminalité (\emph{feymania}) et maintient les PMA dans l'export de matières premières brutes (1 \% commerce mondial). \emph{Synthèse} : La mondialisation n'est pas une chance « donnée » mais un rapport de forces à réguler : l'enjeu pour le Cameroun, comme pour les PMA, est de passer de l'intégration subie (export brut, import manufacturé) à l'intégration choisie (transformation locale, labellisation IGP/équitable, politique industrielle, maîtrise du numérique) pour en capter la valeur ajoutée.
\end{enumerate}
\end{corrige}

\begin{aretenir}[L'essentiel du Dossier 4 — Fiche de révision Bac]
\begin{itemize}
\item L'\cle{immigration clandestine} (entrée/séjour sans titre légal) emprunte 3 routes meurtrières : Méditerranée centrale (Libye $\to$ Italie), Méditerranée occidentale (Maroc $\to$ Espagne), Atlantique (Afrique de l'Ouest $\to$ Canaries). Bilan OIM : 28 000+ morts en Méditerranée depuis 2014.
\item La \cle{cybercriminalité} (escroqueries, piratage, rançongiciels, blanchiment) exploite l'anonymat et la transnationalité des réseaux. Au Cameroun : \cle{feymania}. Coût mondial : plusieurs milliers de milliards \$/an.
\item Les conflits armés sont alimentés par le trafic d'armes et les \cle{minerais de sang} (coltan, or, diamants, étain) qui financent les belligérants et contaminent les chaînes de valeur mondiales (électronique).
\item Les \cle{PMA} ($\approx$ 1,1 Md d'habitants, 46 pays) ne pèsent qu'\textbf{$\approx$ 1 \% du commerce mondial} et attirent peu d'IDE : c'est la \cle{marginalisation}, cercle vicieux (pauvreté $\to$ export brut $\to$ pauvreté).
\item Réponses multiniveaux : Régulation internationale (Kimberley, BEPS, OMC), Coopération policière (Interpol), Intégration régionale (CEMAC, CEEAC, UA), Valorisation locale (Labels IGP, Commerce équitable, Transformation).
\item \textbf{Au Bac :} Analyse \textbf{nuancée} obligatoire (chance + problèmes + régulation). Pas de procès à charge, pas d'angélisme. Preuves chiffrées et localisées exigées.
\end{itemize}
\end{aretenir}

\coursec{TP : Cartographier la mondialisation et agir compétent}

Après avoir identifié, dans le dossier 4, les ombres de la mondialisation (immigration clandestine, cybercriminalité, conflits, marginalisation des pays pauvres), il est temps de passer de la lecture du monde à sa \cle{représentation}. Un géographe ne se contente pas de décrire : il \textbf{cartographie}, il \textbf{schématise}, il \textbf{calcule}. Ces travaux pratiques te placent dans la peau d'un cartographe et d'un acteur économique camerounais : tu vas légender une carte mondiale, construire un croquis du Cameroun inséré dans les flux mondiaux, traduire des statistiques en diagramme, puis proposer concrètement un label pour valoriser les produits du terroir. C'est l'aboutissement de la famille de situations « la mondialisation » : \emph{agir compétent}.

\courssub{TP 1 : Légender une carte mondiale des flux}

\textbf{Situation de départ.} Ton professeur te remet un planisphère muet (continents esquissés, principales villes-ports repérées par des points) et une série de données : les grandes routes maritimes du pétrole et des conteneurs, les principaux couloirs migratoires, les flux d'IDE, le tracé des câbles sous-marins de fibre optique. Ta mission : transformer ce fond de carte muet en \cle{carte de la mondialisation}.

\begin{methode}[Légender une carte de flux en cinq étapes]
\begin{enumerate}
\item \textbf{Lire le fond de carte} : repérer les continents, les océans, les points (villes, ports) déjà placés. Orienter la carte (flèche du nord).
\item \textbf{Classer l'information} en 3 à 5 rubriques maximum : par exemple (1) flux de marchandises, (2) flux de personnes, (3) flux de capitaux (IDE), (4) flux d'information (câbles).
\item \textbf{Choisir des figurés adaptés} : flèches d'épaisseur proportionnelle à l'importance du flux, couleurs distinctes par nature de flux, points ou étoiles pour les pôles (Triade : Amérique du Nord, Europe occidentale, Asie de l'Est).
\item \textbf{Tracer avec sobriété} : les grandes artères d'abord (détroit de Malacca, canal de Suez, détroit d'Ormuz, route transpacifique), puis les flux secondaires.
\item \textbf{Rédiger la légende} dans un cadre, numérotée ou fléchée, et donner un \textbf{titre} précis à la carte.
\end{enumerate}
\end{methode}

\begin{attention}[Le piège de la surcharge]
L'erreur la plus fréquente au Baccalauréat est de vouloir tout faire figurer : la carte devient illisible et la légende interminable. Une bonne carte de Terminale comporte \textbf{3 à 5 rubriques de légende maximum}. Regroupe : « flux commerciaux majeurs » plutôt que dix routes séparées ; « migrations Sud--Nord » plutôt que chaque couloir. La qualité d'une carte se juge à sa \cle{lisibilité}, pas à son exhaustivité.
\end{attention}

\textbf{Résultat attendu.} La carte légendée doit faire apparaître visuellement l'\cle{inégalité} de la mondialisation : une toile dense de flux entre les trois pôles de la Triade, des routes secondaires vers les pays émergents (Chine, Inde, Brésil, golfe de Guinée pour le pétrole), et des marges quasi vides (une grande partie de l'Afrique intérieure, de l'Asie centrale). Cette simple observation visuelle prépare l'interprétation : la mondialisation est un phénomène \textbf{sélectif}.

\courssub{TP 2 : Construire un croquis « Le Cameroun dans la mondialisation »}

\textbf{Consigne.} À partir du fond de carte du Cameroun ci-dessous, construis un croquis montrant l'insertion du pays dans les flux mondiaux : ports et axes d'exportation, zones de production, origine des IDE, rayonnement de la diaspora.

\begin{popfigure}
\includegraphics[width=\linewidth,height=11cm,keepaspectratio]{N01/cmr_map.jpg}
\legende{Fond de carte du Cameroun (régions, capitales, villes, routes, voies ferrées, fleuves). On y trouve Douala (Littoral), Kribi (Sud) et Yaoundé. Travail demandé : (1) localiser et nommer les ports de Douala et de Kribi ; (2) hachurer ou colorier les zones de production (cacao-café-banane au Sud-Ouest et au Centre, coton et élevage au Nord, pétrole au large) ; (3) tracer les axes d'exportation vers la mer ; (4) indiquer par des flèches entrantes l'origine des IDE (Chine, France, Union européenne) ; (5) figurer la diaspora par des flèches sortantes vers l'Europe et l'Amérique du Nord. \\ Source : Wikimedia Commons, « Cameroon Map », CIA, domaine public.}
\end{popfigure}

\begin{methode}[Les quatre obligations d'un croquis géographique]
Un croquis géographique comporte \textbf{toujours} : (1) un \textbf{titre} qui annonce le sujet et l'espace étudié ; (2) une \textbf{légende organisée} en rubriques hiérarchisées (de la plus importante à la plus secondaire) ; (3) une \textbf{orientation} (flèche du nord) et, si possible, une échelle indicative ; (4) des \textbf{éléments localisés avec sobriété} : chaque figuré est placé au bon endroit, sans surcharge. Un croquis sans titre ou sans légende est sanctionné, même s'il est exact.
\end{methode}

\begin{savaistu}[Kribi, vitrine des IDE chinois au Cameroun]
Le \cle{port en eau profonde de Kribi} (région du Sud), inauguré en 2018, a été financé en grande partie par des prêts et des investissements chinois (China Harbour Engineering Company). Capable d'accueillir de grands porte-conteneurs que Douala ne peut recevoir, il ambitionne de faire du Cameroun une \textbf{plateforme logistique de l'Afrique centrale}, au service du Tchad et de la Centrafrique enclavés. C'est un exemple parfait d'IDE transformant la géographie d'un pays.
\end{savaistu}

\courssub{TP 3 : Construire un diagramme circulaire des exportations camerounaises}

\textbf{Données.} Le tableau suivant donne, en ordres de grandeur, les destinations des exportations du Cameroun (parts approximatives, à titre pédagogique ; les valeurs réelles fluctuent selon les années).

\begin{center}
\begin{tabularx}{\linewidth}{|l|X|X|}\hline
\rowcolor{popblueL} \textbf{Destination} & \textbf{Part approximative (\%)} & \textbf{Angle calculé (degrés)} \\ \hline
Union européenne (Pays-Bas, France, Italie, Espagne) & 45 & 162 \\ \hline
Chine et Asie & 25 & 90 \\ \hline
Afrique (CEMAC, Nigeria) & 15 & 54 \\ \hline
Reste du monde & 15 & 54 \\ \hline
\textbf{Total} & \textbf{100} & \textbf{360} \\ \hline
\end{tabularx}
\end{center}

\begin{methode}[Convertir des pourcentages en degrés]
Pour construire un diagramme circulaire, chaque part doit être convertie en angle au centre :
\[ \text{angle} = \frac{\text{part} \times 360}{\text{total}} \]
Exemple : pour l'Union européenne, $\dfrac{45 \times 360}{100} = 162$ degrés. Vérifie toujours que la somme des angles fait bien $360$ degrés avant de tracer les secteurs au rapporteur.
\end{methode}

\begin{popfigure}
\begin{center}
\begin{tikzpicture}[scale=1.6]
% UE : 162 deg, Asie : 90, Afrique : 54, Reste : 54
\fill[popblue] (0,0) -- (90:1.4) arc (90:252:1.4) -- cycle;
\fill[poporange] (0,0) -- (252:1.4) arc (252:342:1.4) -- cycle;
\fill[popgreen] (0,0) -- (342:1.4) arc (342:396:1.4) -- cycle;
\fill[popgold] (0,0) -- (396:1.4) arc (396:450:1.4) -- cycle;
\draw[popdark,thick] (0,0) circle (1.4);
\node[font=\scriptsize\bfseries,white] at (170:0.85) {UE 45\%};
\node[font=\scriptsize\bfseries,white] at (295:0.85) {Asie 25\%};
\node[font=\scriptsize\bfseries,popdark] at (15:0.9) {15\%};
\node[font=\scriptsize\bfseries,popdark] at (60:0.9) {15\%};
% Légende
\fill[popblue] (2.0,0.9) rectangle (2.3,1.1); \node[right,font=\scriptsize] at (2.3,1.0) {Union européenne};
\fill[poporange] (2.0,0.4) rectangle (2.3,0.6); \node[right,font=\scriptsize] at (2.3,0.5) {Chine et Asie};
\fill[popgreen] (2.0,-0.1) rectangle (2.3,0.1); \node[right,font=\scriptsize] at (2.3,0.0) {Afrique (CEMAC, Nigeria)};
\fill[popgold] (2.0,-0.6) rectangle (2.3,-0.4); \node[right,font=\scriptsize] at (2.3,-0.5) {Reste du monde};
\end{tikzpicture}
\end{center}
\legende{Diagramme circulaire attendu : destinations approximatives des exportations du Cameroun (ordres de grandeur pédagogiques ; source indicative : INS Cameroun, données variables selon les années). La lecture est immédiate : près de la moitié des exportations part vers l'Union européenne.}
\end{popfigure}

\begin{exoresolu}[Calculer et interpréter un angle de secteur]
Énoncé : les exportations camerounaises vers la Chine représentent environ 20 \% du total (ordre de grandeur). Calcule l'angle du secteur correspondant dans un diagramme circulaire, puis formule une phrase d'interprétation.
\tcblower
Solution détaillée :
\begin{align*}
\text{angle} &= \frac{20 \times 360}{100} \\
&= \frac{7200}{100} \\
&= 72 \text{ degrés}
\end{align*}
Interprétation : le secteur « Chine » occupe 72 degrés, soit un cinquième du disque. Cela traduit la montée en puissance de l'Asie comme débouché des matières premières camerounaises (pétrole, bois, coton), aux côtés de l'Europe, partenaire historique mais dont la part relative recule. Piège à éviter : ne jamais reporter le pourcentage directement comme un angle (20 degrés au lieu de 72 serait une erreur grossière).
\end{exoresolu}

\courssub{Interprétation croisée : que révèlent ces représentations ?}

Croisons maintenant les trois documents produits (carte mondiale, croquis du Cameroun, diagramme circulaire). Trois enseignements majeurs se dégagent :

\begin{itemize}
\item \textbf{Une insertion réelle mais dépendante.} Le Cameroun est bien connecté à la mondialisation (deux ports, flux d'IDE, diaspora active), mais il y entre surtout comme \cle{fournisseur de matières premières} non transformées (cacao, pétrole, bois, coton, café, banane) : la valeur ajoutée s'accumule ailleurs.
\item \textbf{Une dépendance aux anciens partenaires, en voie de diversification.} Le diagramme circulaire montre la prédominance européenne (environ 45 \%), héritée de l'histoire, mais l'Asie — la Chine en tête — pèse déjà un quart des échanges : le croquis doit donc faire cohabiter flèches « anciennes » (Europe) et « nouvelles » (Asie).
\item \textbf{Une mondialisation à double sens.} Le pays exporte des matières premières et des hommes (diaspora, dont les transferts financiers sont une ressource vitale), et importe des capitaux (IDE chinois et français), des biens manufacturés et de l'information. Le croquis rend visible cette \textbf{asymétrie} : flèches sortantes de produits bruts, flèches entrantes de produits finis et de capitaux.
\end{itemize}

\begin{aretenir}[La place du Cameroun dans la mondialisation]
Le Cameroun est un acteur \textbf{périphérique mais non marginal} de la mondialisation : intégré par ses ports (Douala, Kribi), ses matières premières et sa diaspora, il reste tributaire des centres de la Triade et de la Chine. L'enjeu national est de monter en gamme : transformer sur place et \cle{labelliser} les produits pour capter davantage de valeur.
\end{aretenir}

\courssub{Agir compétent : atelier « Labelliser les produits camerounais »}

\textbf{Situation.} Tu es membre d'une coopérative de producteurs. Face à l'absence ou à la faible valorisation de certains produits camerounais sur le marché international, ta mission est de proposer un \cle{label} complet pour l'un des produits suivants : cacao, café, banane, poivre de Penja.

\begin{exemplebox}[Le modèle du poivre de Penja]
Le \textbf{poivre de Penja} (département du Moungo, région du Littoral) est le \textbf{premier produit africain à avoir obtenu une Indication Géographique Protégée (IGP)}, reconnue dans l'Union européenne. Ce label garantit l'origine et le savoir-faire local, interdit l'usage du nom hors de la zone délimitée, et permet au poivre de Penja de se vendre à un prix très supérieur au poivre ordinaire sur les marchés européens. C'est la preuve qu'un produit camerounais peut conquérir le marché mondial par la \textbf{qualité} plutôt que par le volume.
\end{exemplebox}

\begin{methode}[Construire un label en quatre livrables]
\begin{enumerate}
\item \textbf{Nom et territoire} : choisir un nom lié à une origine géographique précise (ex. « Cacao des hautes terres de l'Ouest ») et délimiter la zone de production.
\item \textbf{Normes (cahier des charges)} : définir 4 à 6 règles vérifiables — mode de culture sans intrants chimiques excessifs, fermentation et séchage maîtrisés, rémunération minimale des producteurs (dimension \emph{commerce équitable}), traçabilité du lot, contrôle par un organisme indépendant.
\item \textbf{Arguments de commercialisation} : rédiger trois arguments (qualité gustative unique, respect de l'environnement, revenu juste pour le producteur) adaptés au consommateur visé (Europe, Asie, marché régional CEMAC).
\item \textbf{Présentation visuelle} : esquisser le logo et l'emballage mentionnant l'origine, le label et la traçabilité.
\end{enumerate}
\end{methode}

\begin{attention}[Label ne signifie pas slogan]
Un label crédible repose sur des \textbf{normes contrôlables}, pas sur de belles phrases. Dire « cacao de qualité » ne suffit pas : il faut préciser \emph{qui contrôle}, \emph{quoi} (taux d'humidité, absence de pesticides interdits, prix plancher au producteur) et \emph{comment} (audits, certificats). Sans vérification indépendante, le label n'est qu'une publicité et ne résiste pas aux critiques adressées au commerce équitable (coût de certification, faible contrôle réel).
\end{attention}

\courssub{Restitution et évaluation formative}

Chaque groupe présente en cinq minutes ses productions (carte, croquis, diagramme, projet de label). L'évaluation formative s'appuie sur la grille commune suivante :

\begin{center}
\begin{tabularx}{\linewidth}{|l|X|}\hline
\rowcolor{popblueL} \textbf{Critère} & \textbf{Attendus} \\ \hline
Exactitude de la légende & Figurés corrects, rubriques hiérarchisées, 3 à 5 entrées maximum, aucune erreur de localisation. \\ \hline
Proportionnalité & Épaisseur des flèches et taille des secteurs fidèles aux données (angles calculés justes, somme de 360 degrés). \\ \hline
Propreté et lisibilité & Tracés nets, titre présent, orientation indiquée, orthographe des toponymes correcte. \\ \hline
Argumentation & Interprétation croisée pertinente ; label assorti de normes vérifiables et d'arguments de vente convaincants. \\ \hline
\end{tabularx}
\end{center}

\begin{exercice}[Pour t'entraîner]
\begin{enumerate}
\item Sur un fond de planisphère, trace les trois principales routes maritimes de conteneurs et légende-les en une seule rubrique.
\item Les transferts de la diaspora camerounaise représentent plusieurs centaines de millions de dollars par an (ordre de grandeur). Représente ce flux sur ton croquis du TP 2 et justifie ton choix de figuré.
\item Rédige le cahier des charges (5 normes) d'un label « Café des monts Bamboutos » et deux arguments de commercialisation destinés au marché européen.
\end{enumerate}
\end{exercice}

\begin{corrige}[Exercice]
\begin{enumerate}
\item Routes attendues : Asie de l'Est--Europe (par Malacca et Suez), Asie--Amérique du Nord (transpacifique), Europe--Amérique du Nord (transatlantique). Une seule rubrique de légende : « routes maritimes majeures de conteneurs », figuré flèche épaisse.
\item Figuré adapté : flèche entrante (de l'Europe et de l'Amérique du Nord vers le Cameroun), d'épaisseur moyenne, car ce flux financier est important mais inférieur aux flux d'exportation ; on le distingue par une couleur propre aux « flux financiers ».
\item Exemple de normes : café arabica cultivé entre \num{1400} et \num{2000} mètres d'altitude sur les sols volcaniques des Bamboutos ; récolte manuelle sélective ; séchage naturel contrôlé ; prix plancher garanti au producteur ; certification par un organisme indépendant. Arguments : goût d'altitude rare ; achat solidaire qui rémunère justement les producteurs camerounais.
\end{enumerate}
\end{corrige}

Ces travaux pratiques bouclent la leçon : tu sais désormais \textbf{représenter} la mondialisation, \textbf{mesurer} ses flux et \textbf{agir} en acteur compétent pour que le Cameroun y occupe une place plus avantageuse. La leçon suivante te fera découvrir les foyers moteurs de cette économie mondiale : la Triade et ses rivaux.

\newpage

\coursec{Activité d'intégration}

\courssub{Objectifs de l'activité}

À l'issue de cette activité d'intégration, tu seras capable de :

\begin{itemize}
\item \cle{observer}, \cle{localiser} et \cle{décrire} les principaux flux qui relient le Cameroun au reste du monde (marchandises, personnes, capitaux, informations) ;
\item \cle{cartographier} : légender correctement un fond de carte intitulé « Le Cameroun dans la mondialisation » ;
\item \cle{construire un diagramme} en bâtons comparant trois sources de financement extérieur (IDE, aide publique au développement, transferts de la diaspora) ;
\item \cle{interpréter des documents} variés (carte, tableau statistique, photographies, articles de presse) pour établir un diagnostic géographique argumenté ;
\item \cle{agir compétent} : proposer une fiche de labellisation d'un produit camerounais inspirée des normes du commerce équitable, en réponse à une commande professionnelle réaliste.
\end{itemize}

\courssub{Situation}

\textbf{Communiqué de la Chambre de Commerce, d'Industrie, des Mines et de l'Artisanat du Cameroun (CCIMA), Douala.}

Dans le cadre du forum « Invest in Cameroon », la CCIMA souhaite présenter à des investisseurs étrangers un dossier pédagogique sur l'insertion du Cameroun dans la mondialisation. Elle fait appel à ta classe de Terminale : par groupes de quatre, vous devrez produire les quatre pièces du dossier (carte légendée, diagramme, diagnostic rédigé, fiche de labellisation).

\textbf{Pièces documentaires fournies :}

\textbf{Document 1 — Fond de carte à légender.} Croquis du monde très simplifié montrant : le Cameroun, les trois pôles de la Triade (Amérique du Nord, Union européenne, Asie de l'Est), les flux d'exportation de cacao (vers les Pays-Bas et l'Allemagne, environ \SI{200000}{tonnes} par an en ordre de grandeur), de pétrole brut (vers l'Espagne, l'Inde, la Chine), le flux d'importation de produits manufacturés depuis la Chine, et les routes migratoires des jeunes Camerounais vers l'Afrique du Nord puis l'Europe.

\textbf{Document 2 — Tableau statistique : les financements extérieurs du Cameroun (ordres de grandeur annuels récents, en milliards de FCFA).}

\begin{center}
\begin{tabularx}{\linewidth}{|l|X|}\hline
\rowcolor{popblueL} \textbf{Source de financement} & \textbf{Montant annuel (ordre de grandeur)} \\ \hline
Investissements directs étrangers (IDE) & environ 300 milliards de FCFA \\ \hline
Aide publique au développement (APD, fonds bilatéraux et multilatéraux) & environ 400 milliards de FCFA \\ \hline
Transferts d'argent de la diaspora & environ 600 milliards de FCFA \\ \hline
\end{tabularx}
\end{center}

\emph{Source : estimations à partir des données de la Banque mondiale et de la BEAC ; les montants varient fortement d'une année à l'autre, on retient des ordres de grandeur.}

\textbf{Document 3 — Photographies décrites.} (a) Le terminal à conteneurs du port de Kribi, en eau profonde, opéré par un consortium international (Bolloré, CMA-CGM, China Harbour) ; (b) le port de Douala, saturé, qui traite encore l'essentiel du commerce extérieur du pays.

\textbf{Document 4 — Article de presse.} « Le poivre de Penja, première Indication Géographique Protégée (IGP) d'Afrique subsaharienne (2013), se vend jusqu'à dix fois plus cher que le poivre ordinaire sur les marchés européens. Les producteurs, regroupés en association, contrôlent la qualité et négocient directement avec les importateurs. »

\textbf{Document 5 — Extrait de presse.} « Chaque année, des centaines de jeunes Camerounais tentent la traversée du Sahara puis de la Méditerranée. Beaucoup vendent leurs biens pour payer les passeurs ; certains perdent la vie dans le désert ou en mer. »

\courssub{Consignes}

\begin{enumerate}
\item \textbf{Légender la carte.} Sur le fond de carte (document 1), place et nomme : le Cameroun, les trois pôles de la Triade, le port de Kribi, les deux principaux flux d'exportation de marchandises (cacao, pétrole) avec leur destination, le flux d'importation depuis la Chine, et la route migratoire vers l'Europe. Respecte les conventions : flèches orientées pour les flux, épaisseur proportionnelle à l'importance, titre, orientation, légende organisée.
\item \textbf{Construire le diagramme.} À partir du document 2, construis un diagramme en bâtons comparant les trois financements extérieurs. Calcule ensuite la part de chaque source dans le total (en \%) et le rapport entre les transferts de la diaspora et les IDE.
\item \textbf{Rédiger le diagnostic.} En un paragraphe de 10 à 15 lignes, présente les \cle{forces} et les \cle{faiblesses} de l'insertion du Cameroun dans la mondialisation, en exploitant au moins quatre des cinq documents et en citant des chiffres.
\item \textbf{Proposer une fiche de labellisation.} Choisis un produit camerounais (poivre de Penja, cacao, café, miel d'Oku, etc.) et rédige une fiche comportant : le label visé, trois normes à respecter (qualité, prix garanti au producteur, conditions sociales et environnementales), les débouchés visés, et deux critiques ou limites du commerce équitable à anticiper.
\item \textbf{Restitution.} Chaque groupe présente son dossier en 5 minutes ; les autres groupes évaluent avec la grille critériée (exactitude des données, qualité de la légende, correction du diagramme, argumentation du diagnostic, faisabilité de la fiche).
\end{enumerate}

\courssub{Ressources à mobiliser}

\begin{aretenir}[Boîte à outils de la leçon]
\begin{itemize}
\item \textbf{Savoirs :} les flux physiques de marchandises (produits agricoles tropicaux, minerais et hydrocarbures) ; la circulation des personnes (migrations légales et clandestines) ; la circulation de l'information et des services (câbles sous-marins, commerce en ligne) ; les mouvements financiers (fonds bilatéraux et multilatéraux, transferts de la diaspora, IDE) ; les normes et les critiques du commerce équitable ; les problèmes de la mondialisation (immigration clandestine, cybercriminalité, conflits, marginalisation des pays pauvres).
\item \textbf{Savoir-faire cartographique :} titre, orientation (flèche du nord), légende numérotée ou fléchée, flèches de flux d'épaisseur proportionnelle, points nommés pour les villes et ports.
\item \textbf{Calculs :} part en pourcentage : $\text{part} = \dfrac{\text{valeur}}{\text{total}} \times 100$ ; taux de variation : $\dfrac{V_A - V_D}{V_D} \times 100$.
\end{itemize}
\end{aretenir}

\courssub{Démarche et corrigé commenté}

\begin{exoresolu}[Le Cameroun dans les flux mondiaux : corrigé guidé]
Énoncé : produire les quatre pièces du dossier de la CCIMA (carte légendée, diagramme en bâtons, diagnostic, fiche de labellisation).

\tcblower

\textbf{Étape 1 — La carte légendée.} On applique la méthode du croquis : fond simplifié, puis on place d'abord les localisations (points), ensuite les flux (flèches), enfin la légende et le titre.

\begin{popfigure}
\includegraphics[width=\linewidth,height=9.5cm,keepaspectratio]{N12/dev_emerging.png}
\legende{Fond de carte du croquis « Le Cameroun dans la mondialisation » : rose = marchés développés (Amérique du Nord, Europe, Japon, Australie), bleu = marchés émergents (Chine, Inde, Russie, Brésil, Mexique, Afrique du Sud, etc.), gris = autres pays, dont la plus grande partie de l'Afrique. Le Cameroun, en Afrique centrale, est gris. \textbf{Le croquis attendu situe :} 1. les exportations de cacao et de pétrole brut vers l'Union européenne ; 2. les importations de produits manufacturés depuis la Chine ; 3. les routes migratoires vers l'Afrique du Nord puis l'Europe. La carte montre déjà que le Cameroun se trouve à l'écart des pôles développés et émergents. \\ Source : Wikimedia Commons, « Developed and Emerging markets », AlexCovarrubias, domaine public.}
\end{popfigure}

\emph{Commentaire :} la flèche 1 est la plus épaisse car les exportations de matières premières dominent le commerce extérieur camerounais ; la flèche migratoire est en pointillés car une partie de ces migrations est clandestine et mal mesurée.

\textbf{Étape 2 — Le diagramme en bâtons et les calculs.} Total des financements : $300 + 400 + 600 = \num{1300}$ milliards de FCFA. Parts :
\begin{align*}
\text{IDE} &= \frac{300}{1300} \times 100 \approx 23\% &
\text{APD} &= \frac{400}{1300} \times 100 \approx 31\% &
\text{Diaspora} &= \frac{600}{1300} \times 100 \approx 46\%
\end{align*}
Rapport diaspora / IDE : $\dfrac{600}{300} = 2$ : la diaspora envoie environ \textbf{deux fois plus} que les IDE.

\begin{popfigure}
\begin{center}
\begin{tikzpicture}
\begin{axis}[
    ybar, bar width=1.1cm, width=0.85\linewidth, height=6.2cm,
    ymin=0, ymax=700,
    ylabel={Milliards de FCFA (ordre de grandeur)},
    symbolic x coords={IDE, APD, Diaspora},
    xtick=data,
    nodes near coords,
    every node near coord/.append style={font=\small\bfseries, color=popdark},
    ytick={0,100,200,300,400,500,600},
    axis line style={popdark},
    tick label style={color=popdark},
    label style={color=popdark},
]
\addplot[fill=popblue, draw=popdark] coordinates {(IDE,300) (APD,400) (Diaspora,600)};
\end{axis}
\end{tikzpicture}
\end{center}
\legende{Diagramme en bâtons : les trois sources de financement extérieur du Cameroun (ordres de grandeur annuels, milliards de FCFA ; estimations d'après Banque mondiale et BEAC).}
\end{popfigure}

\textbf{Étape 3 — Le diagnostic (paragraphe rédigé).} \emph{Modèle de réponse :} « Le Cameroun est bien inséré dans les flux physiques de la mondialisation : il exporte chaque année environ \SI{200000}{tonnes} de cacao vers l'Europe ainsi que du pétrole brut, et le port en eau profonde de Kribi, opéré par un consortium international, renforce cette ouverture (document 3). Le pays reçoit aussi d'importants capitaux : environ 600 milliards de FCFA de transferts de la diaspora, soit deux fois plus que les IDE (environ 300 milliards) et davantage que l'aide publique (environ 400 milliards). Cependant, cette insertion reste fragile : le Cameroun exporte surtout des matières premières non transformées et importe des produits manufacturés chinois, ce qui dégrade sa balance commerciale ; le port de Douala est saturé ; les jeunes, faute d'emplois, s'engagent dans des migrations clandestines meurtrières (document 5). Enfin, l'expérience du poivre de Penja, première IGP d'Afrique subsaharienne vendu jusqu'à dix fois plus cher, montre qu'une stratégie de qualité et de labellisation peut améliorer la place du pays dans les échanges mondiaux. »

\emph{Commentaire :} le paragraphe cite quatre documents, mobilise des chiffres, et s'achève sur une ouverture qui annonce la quatrième tâche : c'est la structure attendue au Baccalauréat (constats équilibrés, preuves, articulation logique).

\textbf{Étape 4 — La fiche de labellisation.} \emph{Exemple :} produit : cacao fin de la plaine de la Sanaga. Label visé : certification commerce équitable (type Fairtrade/Max Havelaar) adossée à une IGP. Normes : (1) prix plancher garanti au planteur, supérieur au cours mondial, plus une prime de développement versée à la coopérative ; (2) cahier des charges qualité (fermentation et séchage contrôlés, traçabilité du lot) ; (3) interdiction du travail des enfants et limitation des pesticides. Débouchés : chocolatiers européens haut de gamme, vente en ligne, transformation locale (chocolat made in Cameroon). Limites à anticiper : coût de la certification pour les petits planteurs ; marché équitable encore étroit (quelques \% du commerce mondial) ; risque que le prix garanti décourage l'amélioration de la qualité si la demande stagne.
\end{exoresolu}

\begin{attention}[Pièges fréquents dans cette activité]
\begin{itemize}
\item Confondre \cle{IDE} (investissements productifs d'entreprises étrangères) et \cle{transferts de la diaspora} (envois d'argent privés des migrants) : ce sont deux flux financiers de nature différente.
\item Oublier le titre, l'orientation ou la légende du croquis : chaque oubli coûte des points au Baccalauréat.
\item Additionner les pourcentages sans vérifier : $23 + 31 + 46 = 100\%$ (les arrondis doivent toujours être contrôlés).
\item Présenter le commerce équitable comme une solution miracle : le dossier exige aussi ses \cle{critiques} (coût de certification, marché limité).
\end{itemize}
\end{attention}

\courssub{Bilan de l'activité}

Cette activité d'intégration a d'abord montré que la mondialisation n'est pas une abstraction : elle passe par des \cle{flux concrets} et mesurables — conteneurs de cacao à Kribi, virements bancaires de la diaspora, câbles sous-marins de fibre optique, routes migratoires du Sahara. En cartographiant ces flux, tu as vérifié qu'une carte géographique est un \cle{raisonnement} : le choix des flèches, de leur épaisseur et de leur couleur traduit déjà une analyse (dominance des matières premières, asymétrie Nord-Sud).

Le diagramme en bâtons a révélé un fait souvent ignoré : la diaspora camerounaise finance davantage l'économie nationale que les investisseurs étrangers et que l'aide publique réunies séparément. Cette hiérarchie invite à nuancer l'image d'un Cameroun « assisté » : les migrants sont devenus les premiers bailleurs privés du pays, ce qui relie directement la circulation des personnes aux mouvements financiers.

Le diagnostic rédigé a mis en évidence la \cle{contradiction centrale} de l'insertion camerounaise : une ouverture réelle aux flux mondiaux (ports, exportations, capitaux) mais une position subalterne (exportations brutes, importations manufacturées, fuite des jeunes). Enfin, la fiche de labellisation t'a fait passer du constat à l'\cle{action compétente} : à l'image du poivre de Penja, la certification et le commerce équitable peuvent transformer un produit tropical banal en produit de qualité à prix garanti — à condition d'en connaître les normes et d'en anticiper les limites. Tu disposes désormais de tous les outils (observer, cartographier, calculer, argumenter, proposer) pour traiter n'importe quel sujet du Baccalauréat sur la mondialisation.

\newpage

\coursec{Méthodes et exercices résolus}

\coursec{Leçon 12 : Le fonctionnement de la mondialisation}

\courssub{Les flux physiques de marchandises}

\begin{methode}[Analyser une carte de flux commerciaux (marchandises)]
\begin{enumerate}
\item \textbf{Identifier la nature du flux} : produits agricoles (cacao, café, banane, coton), minerais et hydrocarbures (pétrole, bauxite, fer), produits manufacturés. Chaque catégorie a sa logique : les matières premières partent du Sud vers le Nord, les produits manufacturés circulent surtout entre pôles de la Triade.
\item \textbf{Repérer l'origine et la destination} : qui exporte ? qui importe ? Exemple type : le pétrole du golfe de Guinée et du Moyen-Orient vers l'Europe, l'Amérique du Nord et la Chine.
\item \textbf{Lire l'épaisseur des flèches} : elle est proportionnelle au volume ou à la valeur du flux. Comparer les flux majeurs (Asie de l'Est $\rightarrow$ Amérique du Nord) aux flux marginaux (Afrique $\rightarrow$ Afrique).
\item \textbf{Localiser les axes maritimes et les points de passage obligés} : détroits (Malacca, Ormuz, Gibraltar), canaux (Suez, Panama), grands ports (Rotterdam, Shanghai, Singapour, Douala).
\item \textbf{Interpréter} : conclure sur la hiérarchie des espaces (pôles moteurs / interfaces portuaires / arrière-pays marginalisés) et sur la spécialisation des territoires.
\end{enumerate}
\textbf{Réflexes} : toujours citer des chiffres de la carte, toujours nommer les espaces précisément (jamais « les pays riches » mais « l'Amérique du Nord, l'Union européenne, l'Asie de l'Est »).
\end{methode}

\begin{exoresolu}[Les flux mondiaux de pétrole]
\textbf{Énoncé.} Sur une carte des flux pétroliers mondiaux, on lit : le Moyen-Orient exporte environ \num{20} millions de barils par jour (Mb/j), dont les deux tiers vers l'Asie de l'Est ; le golfe de Guinée exporte environ \num{4} Mb/j, principalement vers l'Europe et l'Amérique du Nord. Le détroit d'Ormuz voit transiter environ \num{17} Mb/j.
\begin{enumerate}
\item Décris les grandes orientations des flux pétroliers mondiaux.
\item Calcule la part du détroit d'Ormuz dans les exportations du Moyen-Orient et commente.
\end{enumerate}
\tcblower
\textbf{Solution détaillée.}
\begin{enumerate}
\item \textbf{Description.} Les flux pétroliers mondiaux s'organisent selon une logique Sud $\rightarrow$ Nord et de plus en plus Sud $\rightarrow$ Est : les grandes zones d'extraction (Moyen-Orient, golfe de Guinée, Russie, Amérique latine) alimentent les pôles de consommation de la Triade. Le flux majeur relie le Moyen-Orient à l'Asie de l'Est : $\frac{2}{3}\times 20 \approx \num{13,3}$ Mb/j, soit environ $66\,\%$ des exportations moyen-orientales, ce qui traduit le déplacement du centre de gravité de l'économie mondiale vers la Chine, le Japon et la Corée du Sud. Le golfe de Guinée (Nigeria, Angola, Gabon, Cameroun) exporte environ \num{4} Mb/j vers l'Europe et l'Amérique du Nord : c'est un flux secondaire, cinq fois plus faible que celui du Moyen-Orient, mais stratégique car il sécurise l'approvisionnement atlantique. Ces flux empruntent des routes maritimes jalonnées de points de passage obligés : détroits d'Ormuz et de Malacca, canal de Suez.
\item \textbf{Calcul et commentaire.} Part du détroit d'Ormuz :
\[ \frac{17}{20}\times 100 = 85\,\%. \]
Environ $85\,\%$ des exportations pétrolières du Moyen-Orient transitent par le détroit d'Ormuz. Ce chiffre montre l'extrême vulnérabilité du commerce mondial : un seul point de passage, large de quelques dizaines de kilomètres, conditionne l'approvisionnement énergétique de toute l'Asie de l'Est. Toute tension géopolitique dans cette zone a des répercussions immédiates sur les cours mondiaux.
\end{enumerate}
\textbf{Conclusion.} Les flux pétroliers dessinent une mondialisation hiérarchisée : les pôles de la Triade concentrent la demande, le Sud fournit la matière première, et quelques corridors maritimes étroits structurent l'ensemble du système.
\end{exoresolu}

\courssub{La circulation des personnes}

\begin{methode}[Analyser les migrations internationales]
\begin{enumerate}
\item \textbf{Définir et distinguer} : migrant international (toute personne vivant hors de son pays de naissance), émigration (départ), immigration (arrivée), solde migratoire (immigrés $-$ émigrés), migrations de travail, étudiantes, forcées (réfugiés) et clandestines.
\item \textbf{Chiffrer} : environ \num{280} millions de migrants internationaux dans le monde (ordre de grandeur, ONU), soit environ $3,5\,\%$ de la population mondiale : la majorité des humains ne migre pas.
\item \textbf{Identifier les grands couloirs} : Afrique $\rightarrow$ Europe, Amérique latine $\rightarrow$ Amérique du Nord, Asie du Sud $\rightarrow$ golfe Persique ; mais aussi migrations Sud-Sud (Cameroun $\leftrightarrow$ Tchad, Nigeria, Gabon).
\item \textbf{Expliquer par les facteurs} : facteurs de répulsion (chômage, conflits, pauvreté) et facteurs d'attraction (emplois, sécurité, études).
\item \textbf{Évaluer les effets} : transferts de fonds de la diaspora (ressource majeure pour le Sud), fuite des cerveaux, tensions xénophobes dans les pays d'accueil.
\end{enumerate}
\textbf{Formule à retenir} : solde migratoire $=$ immigrés $-$ émigrés ; taux migratoire $= \dfrac{\text{solde}}{\text{population}}\times 1000$ (en pour mille).
\end{methode}

\begin{exoresolu}[Les transferts de la diaspora camerounaise]
\textbf{Énoncé.} Les transferts d'argent de la diaspora camerounaise sont estimés à environ \num{350} millions de dollars par an (ordre de grandeur, Banque mondiale). Le Cameroun compte environ \num{28} millions d'habitants et un PIB d'environ \num{45} milliards de dollars.
\begin{enumerate}
\item Calcule la part des transferts de la diaspora dans le PIB camerounais.
\item Explique pourquoi ces transferts sont un enjeu de la mondialisation pour le Cameroun.
\end{enumerate}
\tcblower
\textbf{Solution détaillée.}
\begin{enumerate}
\item \textbf{Calcul.} On convertit d'abord dans la même unité : \num{350} millions $= \num{0,35}$ milliard.
\[ \text{Part} = \frac{0,35}{45}\times 100 \approx 0,78\,\%. \]
Les transferts de la diaspora représentent environ $0,8\,\%$ du PIB camerounais. Par habitant, cela représente $\dfrac{350\,000\,000}{28\,000\,000} = \num{12,5}$ dollars par habitant et par an.
\item \textbf{Explication.} Trois arguments sont attendus. (a) Ces transferts sont une \textbf{ressource financière stable} : contrairement aux IDE, ils ne fuient pas en cas de crise et financent directement la consommation, l'éducation, la santé et la construction des familles. (b) Ils dépassent souvent l'aide publique au développement : la diaspora est devenue un acteur majeur du financement du développement, ce qui illustre la circulation des capitaux privés dans la mondialisation. (c) Ils ont une \textbf{contrepartie} : le départ des cadres (médecins, ingénieurs, enseignants) constitue une fuite des cerveaux qui prive le pays de compétences formées à ses frais.
\end{enumerate}
\textbf{Conclusion.} La circulation des personnes n'est pas qu'un déplacement de populations : elle s'accompagne de flux financiers qui intègrent le Cameroun à la mondialisation, avec des effets ambivalents (devises contre fuite des cerveaux).
\end{exoresolu}

\courssub{La circulation de l'information, des services et le commerce en ligne}

\begin{methode}[Analyser les flux immatériels (information, services, e-commerce)]
\begin{enumerate}
\item \textbf{Définir les flux immatériels} : données numériques (câbles sous-marins de fibre optique, satellites), services (finance, tourisme, délocalisations de services ou \emph{offshoring}), commerce en ligne (achats via Internet : Jumia, Alibaba, Amazon).
\item \textbf{Mesurer la fracture numérique} : comparer les taux d'accès à Internet (plus de $90\,\%$ en Amérique du Nord, environ $40\,\%$ en Afrique subsaharienne — ordres de grandeur).
\item \textbf{Localiser les infrastructures} : câbles sous-marins (le Cameroun est relié par des câbles comme SAT-3 et ACE au départ de l'Europe), centres de données, places financières ouvertes 24 h/24 (la succession Tokyo $\rightarrow$ Londres $\rightarrow$ New York).
\item \textbf{Expliquer la logique des services} : délocalisation vers les pays à bas coût de main-d'œuvre qualifiée (Inde pour l'informatique, Maroc pour les centres d'appels francophones).
\item \textbf{Conclure sur la double nature} : les flux immatériels rapprochent les espaces (« village global ») mais creusent les inégalités entre connectés et déconnectés.
\end{enumerate}
\textbf{Réflexe} : ne jamais écrire « l'information circule partout instantanément » sans nuancer par la fracture numérique.
\end{methode}

\begin{exoresolu}[La fracture numérique en chiffres]
\textbf{Énoncé.} En Afrique subsaharienne, environ $40\,\%$ de la population utilise Internet, contre environ $90\,\%$ en Amérique du Nord (ordres de grandeur, UIT). La population de l'Afrique subsaharienne est d'environ \num{1,2} milliard d'habitants.
\begin{enumerate}
\item Calcule le nombre d'habitants d'Afrique subsaharienne non connectés à Internet.
\item À partir de ces chiffres et de tes connaissances, montre que la circulation de l'information renforce la hiérarchie des espaces dans la mondialisation.
\end{enumerate}
\tcblower
\textbf{Solution détaillée.}
\begin{enumerate}
\item \textbf{Calcul.} La part non connectée est $100\,\% - 40\,\% = 60\,\%$.
\[ 1,2 \times \frac{60}{100} = 0,72 \text{ milliard} = 720 \text{ millions d'habitants}. \]
Environ 720 millions d'habitants d'Afrique subsaharienne n'ont pas accès à Internet, soit près de trois fois la population de l'Union européenne.
\item \textbf{Démonstration.} (a) Les flux d'information empruntent des infrastructures coûteuses (câbles sous-marins, centres de données) contrôlées par les pôles de la Triade et leurs firmes (Google, Meta, Microsoft) : la maîtrise des réseaux conforte la domination des pôles moteurs. (b) Les services à haute valeur ajoutée (finance, ingénierie, logiciels) se concentrent dans les métropoles mondiales connectées en permanence, tandis que les espaces mal connectés sont cantonnés à la fourniture de matières premières. (c) Le commerce en ligne illustre cette hiérarchie : un consommateur de Douala peut commander sur Alibaba, mais les plateformes, les paiements et les bénéfices restent aux mains des firmes du Nord et de la Chine. (d) La fracture numérique est donc à la fois une \emph{cause} et une \emph{conséquence} de la marginalisation : sans connexion, pas d'accès aux marchés mondiaux des services ; sans marchés, pas de ressources pour se connecter.
\end{enumerate}
\textbf{Conclusion.} Loin d'effacer les inégalités, la circulation mondiale de l'information les redessine : elle crée un monde à plusieurs vitesses, entre espaces hyperconnectés et espaces marginalisés.
\end{exoresolu}

\courssub{Les mouvements financiers et les IDE}

\begin{methode}[Traiter un sujet sur les flux financiers et les IDE]
\begin{enumerate}
\item \textbf{Définir précisément} : IDE (Investissements Directs à l'Étranger) $=$ investissement durable d'une firme à l'étranger (création ou rachat d'entreprise, contrôle d'au moins $10\,\%$ du capital selon le FMI). À distinguer des placements spéculatifs de portefeuille, des fonds bilatéraux (d'un État à un État) et multilatéraux (FMI, Banque mondiale, BAD).
\item \textbf{Identifier les acteurs} : les firmes transnationales (FTN) comme TotalEnergies, MTN, Bolloré, qui décident de la localisation des IDE.
\item \textbf{Lire la géographie des IDE} : ils circulent surtout \emph{entre} pays développés et vers la Chine ; l'Afrique ne reçoit qu'une faible part des IDE mondiaux (environ $3$ à $5\,\%$ — ordre de grandeur, CNUCED), concentrée sur quelques pays et secteurs extractifs (pétrole, mines).
\item \textbf{Expliquer les facteurs d'attractivité} : coût de la main-d'œuvre, matières premières, stabilité politique, infrastructures, taille du marché.
\item \textbf{Évaluer les effets} : emplois, transferts de technologies et devises, mais aussi rapatriement des bénéfices, pression sur les ressources, dépendance.
\end{enumerate}
\textbf{Réflexe} : un chiffre d'IDE sans destination ni secteur ne vaut rien dans une copie : précise toujours \emph{qui investit, où, dans quoi}.
\end{methode}

\begin{exoresolu}[Les IDE en Afrique : une intégration sélective]
\textbf{Énoncé.} Les IDE mondiaux représentent environ \num{1300} milliards de dollars par an (ordre de grandeur, CNUCED). L'Afrique en reçoit environ \num{45} milliards, dont une grande partie vers cinq pays (Égypte, Afrique du Sud, Nigeria, Maroc, République démocratique du Congo) et vers les secteurs pétrolier et minier.
\begin{enumerate}
\item Calcule la part de l'Afrique dans les IDE mondiaux.
\item Montre que les IDE illustrent une intégration sélective de l'Afrique à la mondialisation.
\end{enumerate}
\tcblower
\textbf{Solution détaillée.}
\begin{enumerate}
\item \textbf{Calcul.}
\[ \frac{45}{1300}\times 100 \approx 3,5\,\%. \]
L'Afrique, qui compte environ $18\,\%$ de la population mondiale, ne reçoit qu'environ $3,5\,\%$ des IDE mondiaux : sa part est cinq fois inférieure à son poids démographique.
\item \textbf{Démonstration en trois points.} (a) \emph{Sélectivité spatiale} : les IDE africains se concentrent sur une poignée de pays dotés de ressources (pétrole du Nigeria, minerais de la RDC) ou de marchés (Égypte, Maroc, Afrique du Sud) ; la majorité des États africains, dont beaucoup de pays enclavés, restent en dehors des circuits. (b) \emph{Sélectivité sectorielle} : les capitaux visent l'extraction (hydrocarbures, mines) et non la transformation locale ; les bénéfices sont largement rapatriés par les FTN, ce qui limite les retombées. (c) \emph{Dépendance} : cette intégration par les matières premières expose les économies africaines aux fluctuations des cours mondiaux et perpétue la division internationale du travail héritée de la colonisation (le Sud fournit le brut, le Nord transforme). Au Cameroun, les IDE ciblent le pétrole, le bois et les grands travaux (port de Kribi), souvent avec des partenaires chinois.
\end{enumerate}
\textbf{Conclusion.} Les mouvements financiers mondiaux intègrent l'Afrique de façon partielle et inégale : le continent est présent dans la mondialisation comme fournisseur de matières premières, mais marginalisé dans les flux de capitaux productifs.
\end{exoresolu}

\courssub{Dossier 3 : Le commerce équitable (activité guidée)}

\begin{methode}[Traiter le dossier sur le commerce équitable]
\begin{enumerate}
\item \textbf{Définir le concept} : commerce équitable $=$ partenariat commercial fondé sur un prix minimum garanti au producteur, une prime de développement, des relations durables et le respect de normes sociales et environnementales (labels : Fairtrade/Max Havelaar).
\item \textbf{Partir du constat du déséquilibre} : les cours des matières premières (cacao, café) sont volatils et fixés sur les bourses du Nord ; le producteur camerounais de cacao ne reçoit qu'une faible part du prix final d'une tablette de chocolat (souvent moins de $10\,\%$).
\item \textbf{Présenter les normes} : prix plancher, prime collective, interdiction du travail des enfants, respect de l'environnement, contrats de long terme.
\item \textbf{Nuancer par les critiques} : coût de la certification inaccessible aux plus pauvres, faibles volumes écoulés (le commerce équitable reste marginal, quelques pourcents du commerce mondial), prix parfois peu avantageux si le cours mondial dépasse le prix plancher, risque de « marketing éthique » des grandes marques.
\item \textbf{Conclure de façon équilibrée} : le commerce équitable corrige localement les injustices sans transformer les règles globales du commerce mondial.
\end{enumerate}
\textbf{Réflexe « agir compétent »} : le programme propose d'agir par la \emph{labellisation des produits} : montre comment un label peut valoriser le cacao ou le café camerounais sur le marché international.
\end{methode}

\begin{exoresolu}[Le cacao camerounais et le commerce équitable]
\textbf{Énoncé.} Un planteur camerounais vend son cacao \num{1000} FCFA le kilogramme. Une coopérative labellisée commerce équitable garantit un prix plancher de \num{1300} FCFA/kg plus une prime de développement de \num{150} FCFA/kg. Le cours mondial chute et le prix conventionnel tombe à \num{900} FCFA/kg.
\begin{enumerate}
\item Calcule le gain relatif du planteur labellisé par rapport au prix conventionnel après la chute du cours.
\item Rédige un paragraphe argumenté (10 à 15 lignes) présentant les atouts et les limites du commerce équitable pour les producteurs camerounais.
\end{enumerate}
\tcblower
\textbf{Solution détaillée.}
\begin{enumerate}
\item \textbf{Calcul.} Revenu du planteur labellisé par kilogramme : $1\,300 + 150 = 1\,450$ FCFA. Gain absolu par rapport au prix conventionnel : $1\,450 - 900 = 550$ FCFA/kg. Gain relatif :
\[ \frac{550}{900}\times 100 \approx 61\,\%. \]
Grâce au prix plancher, le planteur labellisé gagne environ $61\,\%$ de plus que le planteur conventionnel : le commerce équitable joue pleinement son rôle d'amortisseur face à la volatilité des cours. (Si le cours mondial dépassait \num{1450} FCFA/kg, l'avantage s'inverserait : c'est une limite du dispositif.)
\item \textbf{Paragraphe argumenté (corrigé type).} Le commerce équitable apparaît comme une réponse concrète au déséquilibre du commerce mondial, qui défavorise les producteurs de matières premières du Sud. Au Cameroun, où le cacao et le café sont vendus à des cours fixés sur les bourses de Londres ou de New York, le prix plancher garanti protège le planteur des effondrements des cours, comme le montre l'exemple chiffré : il sécurise un revenu et permet d'investir (scolarisation, équipements). La prime de développement, gérée collectivement par la coopérative, finance des infrastructures (puits, écoles, centres de santé) : c'est une démarche de développement local et durable, renforcée par des normes sociales (refus du travail des enfants) et environnementales. Cependant, le dispositif a des limites réelles : le coût de la certification exclut souvent les producteurs les plus pauvres ; les volumes labellisés restent faibles, car la demande mondiale de produits équitables est limitée à des consommateurs sensibilisés du Nord ; enfin, certaines grandes marques utilisent le label comme argument marketing sans changer leurs pratiques globales. Le commerce équitable corrige donc localement les injustices du libre-échange, mais il ne remet pas en cause les règles du commerce mondial, dominé par les pôles de la Triade et leurs firmes.
\end{enumerate}
\textbf{Conclusion.} Le commerce équitable est un outil efficace de stabilisation des revenus et de développement local, mais sa portée reste marginale à l'échelle du commerce mondial.
\end{exoresolu}

\courssub{Dossier 4 : Les problèmes liés à la mondialisation (activité guidée)}

\begin{methode}[Construire une étude sur les problèmes de la mondialisation]
\begin{enumerate}
\item \textbf{Classer les problèmes en catégories} : sociaux (immigration clandestine, trafics), numériques (cybercriminalité : arnaques, piratage, \emph{scamming}), sécuritaires (conflits armés alimentés par le commerce des matières premières et des armes), économiques (marginalisation des pays pauvres, dumping, évasion fiscale des FTN).
\item \textbf{Montrer le lien avec la mondialisation} : c'est la \emph{même} libre circulation (personnes, capitaux, données) qui produit les richesses et les dérives ; les réseaux criminels utilisent les mêmes infrastructures (ports, Internet, système bancaire) que les acteurs licites.
\item \textbf{Illustrer par des exemples précis} : routes migratoires méditerranéennes, cyberarnaques en Afrique de l'Ouest, rôle des ressources (diamants, coltan) dans les conflits africains, part faible des PMA dans le commerce mondial (environ $1\,\%$).
\item \textbf{Évoquer les réponses} : coopération policière internationale (Interpol), régulation financière, aide au développement, politiques migratoires.
\item \textbf{Conclure} : la mondialisation appelle une gouvernance mondiale que les institutions actuelles (ONU, OMC, FMI) ne parviennent qu'imparfaitement à assurer.
\end{enumerate}
\textbf{Réflexe} : évite le ton moralisateur ; le correcteur attend une \emph{analyse géographique} (acteurs, espaces, flux, échelles), pas un jugement de valeur.
\end{methode}

\begin{exoresolu}[Commentaire composé : les faces cachées de la mondialisation]
\textbf{Énoncé.} « La mondialisation profite à tous les espaces et à tous les acteurs. » À partir d'exemples précis, discute cette affirmation en rédigeant un développement construit (introduction, deux parties, conclusion).
\tcblower
\textbf{Solution détaillée (corrigé type).}

\emph{Introduction.} La mondialisation, processus d'intensification des échanges de marchandises, de personnes, de capitaux et d'informations à l'échelle planétaire, a dopé la croissance mondiale et enrichi les pôles moteurs. L'affirmation selon laquelle elle « profite à tous » mérite cependant d'être discutée : si elle crée des richesses considérables, elle génère aussi des problèmes spécifiques et marginalise des espaces entiers. On montrera d'abord que la mondialisation produit des dérives qui frappent tous les espaces, puis qu'elle creuse les inégalités en marginalisant les pays pauvres.

\emph{I. Une mondialisation qui génère des problèmes transnationaux.} La libre circulation a son revers. (a) \textbf{L'immigration clandestine} : les écarts de richesse poussent chaque année des milliers d'Africains vers l'Europe par les routes du Sahara et de la Méditerranée, au péril de leur vie ; les réseaux de passeurs exploitent les mêmes corridors que les flux licites. (b) \textbf{La cybercriminalité} : la circulation instantanée des données permet arnaques, piratages bancaires et vols d'identité à l'échelle planétaire ; aucun État ne peut y répondre seul, d'où la coopération via Interpol. (c) \textbf{Les conflits armés} se mondialisent : le commerce des matières premières (diamants, coltan) et le trafic d'armes financent des rebellions, notamment en Afrique centrale, tandis que les groupes armés s'insèrent dans les réseaux internationaux.

\emph{II. Une mondialisation inégalitaire qui marginalise les pays pauvres.} (a) Les pays les moins avancés ne pèsent qu'environ $1\,\%$ du commerce mondial alors qu'ils regroupent près d'un milliard d'habitants : ils restent fournisseurs de matières premières brutes aux cours instables. (b) L'Afrique ne reçoit qu'environ $3,5\,\%$ des IDE mondiaux, concentrés sur quelques pays et secteurs extractifs. (c) La fracture numérique (environ $40\,\%$ d'internautes en Afrique subsaharienne contre $90\,\%$ en Amérique du Nord) exclut les pays pauvres des activités les plus rentables (services, finance, commerce en ligne). (d) L'évasion fiscale des firmes transnationales prive enfin les États du Sud de recettes essentielles.

\emph{Conclusion.} L'affirmation est donc largement discutable : la mondialisation profite d'abord aux pôles de la Triade et à leurs firmes, tandis qu'elle engendre des dérives transnationales et marginalise les pays pauvres. Sa régulation exige une gouvernance mondiale renforcée et des initiatives correctrices comme le commerce équitable.
\end{exoresolu}

\courssub{TP : Cartographier la mondialisation et agir compétent}

\begin{methode}[Réaliser un croquis de synthèse sur les flux mondiaux]
\begin{enumerate}
\item \textbf{Tracer le fond de carte} : planisphère très simplifié (blocs de continents) ou espace régional ; ajoute le titre en haut, l'orientation et l'échelle indicative.
\item \textbf{Choisir peu d'informations hiérarchisées} (maximum 4 ou 5 rubriques) : pôles de la Triade (aplats colorés), flux majeurs (flèches d'épaisseur proportionnelle), points de passage (symboles), espaces marginalisés (hachures).
\item \textbf{Respecter la sémiologie} : flèches pour les flux, figurés ponctuels pour les ports et places financières, surfaces pour les pôles ; jamais de couleur sans signification.
\item \textbf{Construire la légende} : ordonnée (du plus important au moins important), chaque symbole reproduit à l'identique, titre de légende explicite.
\item \textbf{Rédiger une courte analyse} du croquis (5 à 8 lignes) : ce que montre le document, ce qu'il prouve sur la hiérarchie des espaces.
\end{enumerate}
\textbf{Réflexe} : un croquis sans titre, sans nord ou avec une légende désordonnée perd systématiquement des points, même s'il est exact.
\end{methode}

\begin{exoresolu}[Croquis : les pôles et les flux de la mondialisation]
\textbf{Énoncé.} Construis un croquis du monde montrant : les trois pôles de la Triade, les principaux flux maritimes de marchandises, deux points de passage stratégiques et les espaces marginalisés. Rédige ensuite un commentaire de 6 lignes.
\tcblower
\textbf{Solution détaillée.}

\emph{Étape 1 — le fond de carte.} On trace un planisphère simplifié en blocs : Amérique du Nord à gauche en haut, Europe-Afrique au centre, Asie à droite, Amérique du Sud en bas à gauche.

\emph{Étape 2 — les informations.} (a) Aplats colorés sur les trois pôles : Amérique du Nord, Union européenne, Asie de l'Est (Japon, Chine, Corée du Sud). (b) Flèches épaisses : Asie de l'Est $\rightarrow$ Amérique du Nord (Pacifique nord), Asie de l'Est $\rightarrow$ Europe (océan Indien, canal de Suez), Europe $\leftrightarrow$ Amérique du Nord (Atlantique nord) ; flèche fine : golfe de Guinée $\rightarrow$ Europe (pétrole). (c) Figurés ponctuels : canal de Suez et détroit de Malacca. (d) Hachures sur l'intérieur de l'Afrique subsaharienne, le Sahel et une partie de l'Asie centrale : espaces marginalisés.

\emph{Étape 3 — la légende ordonnée} : 1. Pôles moteurs (Triade) ; 2. Flux majeurs de marchandises ; 3. Flux secondaires ; 4. Points de passage stratégiques ; 5. Espaces marginalisés.

\emph{Commentaire type (6 lignes).} Ce croquis montre un monde organisé autour de trois pôles moteurs — Amérique du Nord, Union européenne et Asie de l'Est — qui concentrent les richesses et s'échangent l'essentiel des marchandises le long de trois axes maritimes majeurs (Atlantique nord, Pacifique nord, route de Suez). Ces flux empruntent des points de passage stratégiques (Suez, Malacca) dont le contrôle est un enjeu géopolitique. À l'écart de ces corridors, de vastes espaces, notamment en Afrique subsaharienne, restent marginalisés : la mondialisation dessine ainsi une géographie profondément inégale, faite de pôles, de flux et de marges.

\begin{popfigure}
\includegraphics[width=\linewidth,height=9.5cm,keepaspectratio]{N12/dev_emerging.png}
\legende{Fond de carte pour le croquis de synthèse : rose = marchés développés (Amérique du Nord, Europe, Japon, Australie), bleu = marchés émergents (Chine, Inde, Russie, Brésil, Mexique, Afrique du Sud, etc.), gris = autres pays, dont la plus grande partie de l'Afrique. Le croquis attendu y situe : la Triade (aplats foncés sur l'Amérique du Nord, l'Union européenne et l'Asie de l'Est), les flux majeurs de marchandises (flèches épaisses), le flux secondaire de pétrole du golfe de Guinée vers l'Europe (flèche pointillée), les points de passage stratégiques (Suez, Malacca) et les espaces marginalisés (hachures sur l'Afrique subsaharienne, le Sahel et une partie de l'Asie centrale). \\ Source : Wikimedia Commons, « Developed and Emerging markets », AlexCovarrubias, domaine public.}
\end{popfigure}
\end{exoresolu}

\courssub{Synthèse et vigilance}

\begin{attention}[Les 5 erreurs qui coûtent des points au Bac]
\begin{enumerate}
\item \textbf{Confondre les catégories de flux.} Ne mélange pas IDE (investissement durable d'une firme), aide publique (fonds bilatéraux et multilatéraux) et transferts de la diaspora (envois privés des migrants) : le correcteur sanctionne toute confusion entre ces trois flux financiers.
\item \textbf{Affirmer que « tout le monde migre » ou que « l'Afrique est inondée d'IDE ».} Les migrants internationaux ne représentent qu'environ $3,5\,\%$ de l'humanité et l'Afrique ne reçoit qu'environ $3,5\,\%$ des IDE mondiaux : des chiffres faux ou exagérés décrédibilisent toute la copie.
\item \textbf{Oublier la proportionnalité et les unités dans les calculs.} Dans un calcul de part, convertis d'abord dans la même unité (millions/milliards) avant de diviser, et termine toujours par $\times 100$ pour obtenir un pourcentage ; annonce l'unité du résultat ($\%$, FCFA/kg, Mb/j).
\item \textbf{Présenter le commerce équitable comme une solution miracle.} Une copie qui vante le commerce équitable sans citer ses critiques (coût de la certification, volumes marginaux, marketing éthique) reste en dessous de la moyenne : le programme exige explicitement les \emph{normes} ET les \emph{critiques}.
\item \textbf{Rendre un croquis sans titre, sans nord ou avec une légende désordonnée.} En TP cartographique, la forme compte autant que le fond : titre, orientation, légende ordonnée et symboles cohérents sont des critères de notation à part entière.
\end{enumerate}
\end{attention}

\newpage

\coursec{Exercices}

\courssub{Je vérifie mes connaissances}

\begin{exercice}[QCM : les notions clés de la mondialisation]
Pour chaque question, choisis la bonne réponse et justifie ton choix en une phrase.
\begin{enumerate}
\item Un \cle{IDE} (Investissement Direct à l'Étranger) est :
\begin{enumerate}
\item un don accordé par un État riche à un État pauvre ;
\item un investissement réalisé par une entreprise hors de son pays d'origine, avec prise de contrôle durable d'une activité ;
\item un prêt du FMI à un État endetté ;
\item un transfert d'argent effectué par un migrant.
\end{enumerate}
\item Le \cle{commerce équitable} repose principalement sur :
\begin{enumerate}
\item la suppression de toutes les barrières douanières ;
\item un prix garanti au producteur, une prime de développement et des relations commerciales durables ;
\item la vente exclusive des produits sur Internet ;
\item l'interdiction des importations agricoles.
\end{enumerate}
\item Les \cle{fonds multilatéraux} sont des financements provenant :
\begin{enumerate}
\item d'un seul État vers un autre État ;
\item d'organisations regroupant plusieurs États (Banque mondiale, FMI, Union européenne) ;
\item des entreprises privées uniquement ;
\item des migrants uniquement.
\end{enumerate}
\item La \cle{fracture numérique} désigne :
\begin{enumerate}
\item la coupure des câbles sous-marins ;
\item les inégalités d'accès aux technologies de l'information et de la communication entre pays et entre populations ;
\item le coût élevé des ordinateurs ;
\item la concurrence entre opérateurs téléphoniques.
\end{enumerate}
\end{enumerate}
\end{exercice}

\begin{exercice}[Vrai ou faux ? Justifie ta réponse]
Indique si chacune des affirmations suivantes est vraie ou fausse, puis justifie en deux ou trois lignes.
\begin{enumerate}
\item Les flux maritimes concentrent l'essentiel du commerce mondial de marchandises en volume.
\item Les migrations internationales concernent uniquement les déplacements du Sud vers le Nord.
\item Les transferts d'argent de la diaspora vers le Cameroun sont négligeables par rapport à l'aide publique au développement.
\item La cybercriminalité est un problème qui ne concerne que les pays développés.
\item Les IDE en Afrique se dirigent majoritairement vers les secteurs pétrolier et minier.
\end{enumerate}
\end{exercice}

\begin{exercice}[Définitions express]
Définis en deux ou trois lignes chacun des termes suivants, en donnant à chaque fois un exemple concret (camerounais si possible) :
\begin{enumerate}
\item flux de marchandises ;
\item migration internationale ;
\item commerce en ligne (e-commerce) ;
\item aide publique au développement (APD) ;
\item immigration clandestine ;
\item marginalisation (d'un pays dans la mondialisation).
\end{enumerate}
\end{exercice}

\begin{exercice}[Calculs directs sur les flux financiers]
On donne les montants approximatifs des flux financiers reçus par l'Afrique subsaharienne au cours d'une année récente (ordres de grandeur, en milliards de dollars) : transferts des migrants : 48 ; IDE : 40 ; aide publique au développement : 30.
\begin{enumerate}
\item Calcule le montant total de ces trois flux.
\item Calcule la part (en \%) de chaque flux dans ce total. Arrondis à une décimale.
\item Les transferts des migrants représentent-ils plus que l'APD ? Exprime l'écart en milliards de dollars et en pourcentage relatif par rapport à l'APD.
\item Si les IDE augmentent de 10 \% l'année suivante, quel serait leur nouveau montant ?
\end{enumerate}
\end{exercice}

\courssub{Je m'entraîne}

\begin{exercice}[Classer les flux de la mondialisation]
Voici une liste de phénomènes observés dans l'espace mondial :
\begin{enumerate}
\item un camionneur transporte du cacao de Kumba vers le port de Douala pour l'exportation ;
\item une étudiante camerounaise s'inscrit dans une université au Maroc ;
\item une entreprise française installe une filiale de téléphonie à Yaoundé ;
\item un habitant de Bafoussam commande un téléphone sur un site chinois ;
\item un migrant camerounais installé en France envoie de l'argent à sa famille par transfert électronique ;
\item le pétrole tchadien transite par l'oléoduc Tchad-Cameroun jusqu'à Kribi ;
\item une banque nigériane ouvre une agence à Douala ;
\item un jeune de Bamenda suit une formation en ligne délivrée par une université américaine.
\end{enumerate}
\textbf{Travail à faire :} classe chaque phénomène dans l'une des quatre catégories suivantes : flux physiques de marchandises ; circulation des personnes ; circulation de l'information et des services ; mouvements financiers et IDE. Présente ton résultat sous forme de tableau à deux colonnes (catégorie / numéros des phénomènes), puis rédige une phrase de conclusion sur la diversité des flux de la mondialisation.
\end{exercice}

\begin{exercice}[Lire un tableau statistique : les exportations du Cameroun]
On donne la structure approximative des exportations camerounaises en valeur (ordres de grandeur récents) :
\begin{center}
\begin{tabularx}{\linewidth}{|l|X|}\hline
\rowcolor{popblueL} \textbf{Produit exporté} & \textbf{Part approximative (\%)} \\ \hline
Pétrole brut & 40 \\ \hline
Cacao & 12 \\ \hline
Bois et produits du bois & 10 \\ \hline
Coton & 6 \\ \hline
Café & 3 \\ \hline
Bananes & 3 \\ \hline
Autres produits & 26 \\ \hline
\end{tabularx}
\end{center}
\textbf{Travail à faire :}
\begin{enumerate}
\item Calcule la part cumulée des produits agricoles (cacao, bois, coton, café, bananes).
\item Construis un diagramme circulaire ou semi-circulaire représentant cette structure (indique les angles calculés pour chaque secteur).
\item Dégage en cinq lignes deux caractéristiques du commerce extérieur camerounais révélées par ce tableau.
\end{enumerate}
\end{exercice}

\begin{exercice}[Construire un diagramme en bâtons : les migrants dans le monde]
On donne le nombre approximatif de migrants internationaux dans le monde (ordres de grandeur, en millions) : 1990 : 153 ; 2000 : 173 ; 2010 : 221 ; 2020 : 281.
\begin{enumerate}
\item Construis un diagramme en bâtons représentant l'évolution du nombre de migrants internationaux (échelle conseillée : 1 cm pour 25 millions ; 1 cm pour 10 ans).
\item Calcule le taux de variation du nombre de migrants entre 1990 et 2020 (formule : $\dfrac{\text{valeur finale} - \text{valeur initiale}}{\text{valeur initiale}} \times 100$). Arrondis à l'unité.
\item Rédige un court commentaire (5 à 8 lignes) de ce diagramme en soulignant ce qu'il révèle de l'intensification de la circulation des personnes dans la mondialisation.
\end{enumerate}
\end{exercice}

\begin{exercice}[Schématiser une filière mondialisée : le cacao camerounais]
À partir de tes connaissances et des informations suivantes — le cacao est produit dans les régions du Centre, du Sud et du Sud-Ouest ; il est acheminé vers le port de Douala ; il est exporté principalement vers l'Union européenne (Pays-Bas, notamment) où il est transformé ; les produits finis (chocolat) sont revendus dans le monde entier, y compris au Cameroun — réalise un schéma légendé de la filière cacao. Ton schéma comportera : les lieux (zones de production, port, destinations), les flux (flèches orientées), et une légende organisée. Termine par une phrase qui montre en quoi cette filière illustre l'insertion du Cameroun dans la mondialisation.
\end{exercice}

\courssub{Je me prépare au Bac}

\begin{exercice}[Type Bac — Étude de documents : les flux de marchandises et la place du Cameroun dans le commerce mondial]
\textbf{Document 1 (carte décrite).} La carte des grands flux maritimes mondiaux montre trois axes dominants : l'axe Atlantique Nord entre l'Amérique du Nord et l'Europe ; l'axe Asie-Europe par le canal de Suez ; l'axe transpacifique entre l'Asie de l'Est et l'Amérique du Nord. Les principaux détroits stratégiques (Ormuz, Malacca, Gibraltar, Panama) concentrent le trafic. Les côtes africaines, hors Afrique du Sud et détroits du Nord, sont traversées par des flux secondaires : le trafic des ports d'Afrique subsaharienne représente moins de 5 \% du trafic portuaire mondial (ordre de grandeur).

\textbf{Document 2 (tableau).} Principales destinations des exportations camerounaises (parts approximatives en valeur, ordres de grandeur récents) :
\begin{center}
\begin{tabularx}{\linewidth}{|l|X|}\hline
\rowcolor{popblueL} \textbf{Destination} & \textbf{Part approximative (\%)} \\ \hline
Union européenne (dont Pays-Bas, Espagne, Italie, France) & 45 \\ \hline
Chine & 15 \\ \hline
Inde & 6 \\ \hline
Pays de la CEMAC et Afrique & 12 \\ \hline
Reste du monde & 22 \\ \hline
\end{tabularx}
\end{center}

\textbf{Document 3.} Le Cameroun exporte surtout des produits primaires (pétrole brut, cacao, bois, coton, café, bananes) et importe des produits manufacturés (machines, véhicules, produits pharmaceutiques), des produits alimentaires transformés et des carburants raffinés. Sa part dans le commerce mondial est inférieure à 0,05 \% (ordre de grandeur).

\textbf{Questions :}
\begin{enumerate}
\item À partir du document 1, décris en quelques lignes la répartition mondiale des grands flux maritimes de marchandises. (4 points)
\item À partir du document 2, présente les principales caractéristiques des destinations des exportations camerounaises. (4 points)
\item En t'appuyant sur les trois documents, dégage la place du Cameroun dans le commerce mondial (nature des échanges, poids, dépendances). (8 points)
\item Propose deux solutions permettant au Cameroun de mieux s'intégrer dans le commerce mondial. (4 points)
\end{enumerate}
\emph{Ton devoir, rédigé, ne doit pas dépasser deux pages.}
\end{exercice}

\begin{exercice}[Type Bac — Croquis : « Les principaux flux de la mondialisation »]
Sur un fond de carte muette du monde (fourni ou reproduit à main levée), construis un croquis légendé intitulé \emph{« Les principaux flux de la mondialisation »}.

\textbf{Consignes :}
\begin{enumerate}
\item Représente par des flèches de natures différentes au moins quatre types de flux : flux de marchandises (grands axes maritimes), migrations internationales (au moins trois grands courants migratoires), IDE (des pôles de la Triade vers les pays émergents et l'Afrique), flux d'informations (câbles sous-marins ou axes numériques majeurs).
\item Localise et nomme les trois pôles de la Triade (Amérique du Nord, Union européenne, Asie de l'Est) ainsi que le Cameroun.
\item Soigne la légende : elle doit être organisée, avec des symboles distincts pour chaque type de flux, et placée hors de la carte.
\item Rédige une courte analyse (10 à 15 lignes) de ton croquis : montre que la mondialisation est hiérarchisée (pôles moteurs, espaces intégrés, espaces marginalisés) et précise la situation de l'Afrique et du Cameroun.
\end{enumerate}
\emph{Critères d'évaluation : exactitude des localisations (5 points), diversité et lisibilité des flux (5 points), qualité de la légende et du titre (5 points), analyse (5 points).}
\end{exercice}

\begin{exercice}[Type Bac — Sujet de synthèse : mondialisation, commerce équitable et pays pauvres]
\textbf{Sujet.} « La mondialisation : opportunités et problèmes pour les pays pauvres. Le commerce équitable peut-il corriger les déséquilibres économiques mondiaux ? Illustrez votre développement à partir du cas camerounais. »

\textbf{Documents d'appui :}

\textbf{Document A.} Le commerce équitable garantit au producteur un prix plancher et une prime de développement versée à la coopérative, en échange du respect de normes sociales (refus du travail des enfants) et environnementales. Au Cameroun, des coopératives de producteurs de cacao et de café de l'Ouest et du Centre sont certifiées et exportent vers l'Europe.

\textbf{Document B.} Le commerce équitable concerne moins de 2 \% du commerce mondial (ordre de grandeur). Les coûts de certification sont élevés pour les petits producteurs. Par ailleurs, les pays pauvres restent confrontés à d'autres problèmes liés à la mondialisation : immigration clandestine, cybercriminalité, conflits armés favorisés par le trafic de ressources, marginalisation dans les flux commerciaux et financiers.

\textbf{Travail à faire :} rédige un développement construit et argumenté comportant :
\begin{enumerate}
\item une introduction avec une définition de la mondialisation, la présentation du problème posé et l'annonce du plan ;
\item une première partie montrant les opportunités offertes par la mondialisation aux pays pauvres (accès aux marchés, IDE, transferts de la diaspora, diffusion de l'information), illustrée par des exemples camerounais ;
\item une deuxième partie présentant les problèmes (déséquilibres des échanges, marginalisation, immigration clandestine, cybercriminalité) ;
\item une troisième partie évaluant la capacité du commerce équitable à corriger ces déséquilibres (normes et apports, puis limites et critiques) ;
\item une conclusion nuancée avec ouverture.
\end{enumerate}
\emph{Barème indicatif : introduction (3 points), développement en trois parties (12 points), exactitude des exemples (3 points), conclusion (2 points).}
\end{exercice}

\newpage

\coursec{Corrigés des exercices}

\begin{corrige}[Exercice 1 — QCM : les notions clés de la mondialisation]
\begin{enumerate}
\item \textbf{Réponse b.} Un \cle{IDE} (Investissement Direct à l'Étranger) est un investissement réalisé par une entreprise hors de son pays d'origine, avec prise de contrôle durable d'une activité (création de filiale, rachat d'entreprise). Exemple : une entreprise française qui installe une cimenterie à Douala. Les réponses a et c relèvent de l'aide publique, la réponse d des transferts de la diaspora.
\item \textbf{Réponse b.} Le \cle{commerce équitable} repose sur un prix garanti (prix plancher) au producteur, une prime de développement versée à la coopérative et des relations commerciales durables et directes. Il ne supprime pas les barrières douanières et ne passe pas nécessairement par Internet.
\item \textbf{Réponse b.} Les \cle{fonds multilatéraux} proviennent d'organisations regroupant plusieurs États (Banque mondiale, FMI, Union européenne, BAD). Un financement d'un seul État vers un autre est un fonds \emph{bilatéral}.
\item \textbf{Réponse b.} La \cle{fracture numérique} désigne les inégalités d'accès aux technologies de l'information et de la communication (Internet, téléphonie, équipements) entre pays développés et pays pauvres, mais aussi au sein d'un même pays (ville/campagne, riches/pauvres).
\end{enumerate}
\textbf{Points de vigilance :} ne pas confondre IDE et APD ; ne pas réduire la fracture numérique à une question de prix des ordinateurs.
\end{corrige}

\begin{corrige}[Exercice 2 — Vrai ou faux ?]
\begin{enumerate}
\item \textbf{Vrai.} Environ 80~\% du commerce mondial de marchandises en volume (ordre de grandeur) transite par voie maritime, car le navire est le mode de transport le moins coûteux pour les grandes masses (pétrole, minerais, céréales, conteneurs).
\item \textbf{Faux.} Les migrations Sud-Nord sont les plus médiatisées, mais les migrations Sud-Sud sont tout aussi importantes en nombre (ex. migrations intra-africaines : Nigérians et Tchadiens vers le Cameroun). Il existe aussi des migrations Nord-Nord et même Nord-Sud (retraités, cadres expatriés).
\item \textbf{Faux.} Les transferts de la diaspora vers le Cameroun représentent plusieurs centaines de millions de dollars par an (ordre de grandeur : plus de 300 millions USD selon la Banque mondiale) ; à l'échelle de l'Afrique subsaharienne, ils dépassent l'aide publique au développement (environ 48 milliards contre 30 milliards USD, ordres de grandeur 2019).
\item \textbf{Faux.} La cybercriminalité touche tous les pays connectés. Le Cameroun et l'Afrique centrale connaissent arnaques en ligne, fraudes bancaires et usurpations d'identité ; la faiblesse des législations y rend parfois les victimes plus vulnérables.
\item \textbf{Vrai.} En Afrique, les IDE se concentrent majoritairement dans les industries extractives (pétrole, gaz, minerais) : pétrole au Nigeria et en Angola, mines en RDC et en Zambie, oléoduc Tchad-Cameroun. Les IDE vers l'agriculture et les services restent minoritaires.
\end{enumerate}
\textbf{Points de vigilance :} une justification de deux ou trois lignes est exigée : une simple mention « vrai » ou « faux » ne rapporte pas tous les points.
\end{corrige}

\begin{corrige}[Exercice 3 — Définitions express]
\begin{enumerate}
\item \textbf{Flux de marchandises :} déplacement physique de biens (matières premières, produits agricoles, manufacturés) entre des lieux de production et des lieux de consommation, par voie maritime, terrestre ou aérienne. \emph{Exemple :} l'exportation du cacao de Kumba vers le port de Douala puis vers les Pays-Bas.
\item \textbf{Migration internationale :} déplacement durable de personnes franchissant une frontière d'État, pour des raisons économiques, d'études ou politiques. \emph{Exemple :} une étudiante camerounaise qui s'installe au Maroc pour ses études.
\item \textbf{Commerce en ligne (e-commerce) :} achat et vente de biens ou de services par voie électronique, via Internet, sans contact physique entre vendeur et acheteur. \emph{Exemple :} un habitant de Bafoussam qui commande un téléphone sur un site chinois et paie par mobile money.
\item \textbf{Aide publique au développement (APD) :} dons et prêts à conditions avantageuses accordés par les États riches (aide bilatérale) ou les organisations internationales (aide multilatérale) aux pays pauvres. \emph{Exemple :} un financement de la Banque mondiale pour un programme d'électrification rurale au Cameroun.
\item \textbf{Immigration clandestine :} entrée ou séjour d'une personne sur le territoire d'un État sans autorisation légale (sans visa ni titre de séjour). \emph{Exemple :} des jeunes d'Afrique centrale qui tentent la traversée de la Méditerranée vers l'Europe sans documents.
\item \textbf{Marginalisation :} situation d'un pays ou d'une région faiblement intégré aux flux de la mondialisation (commerce, IDE, information), relégué à la périphérie du système mondial. \emph{Exemple :} le Cameroun, dont la part dans le commerce mondial est inférieure à 0,05~\%.
\end{enumerate}
\end{corrige}

\begin{corrige}[Exercice 4 — Calculs directs sur les flux financiers]
\begin{enumerate}
\item \textbf{Montant total :}
\[ 48 + 40 + 30 = 118 \text{ milliards de dollars.} \]
\item \textbf{Part de chaque flux :}
\begin{align*}
\text{Transferts des migrants : } & \dfrac{48}{118}\times 100 \approx 40,7\,\% \\
\text{IDE : } & \dfrac{40}{118}\times 100 \approx 33,9\,\% \\
\text{APD : } & \dfrac{30}{118}\times 100 \approx 25,4\,\%
\end{align*}
Vérification : $40,7 + 33,9 + 25,4 = 100,0\,\%$. Le total est correct.
\item \textbf{Comparaison migrants / APD :} oui, les transferts des migrants (48) dépassent l'APD (30). Écart absolu :
\[ 48 - 30 = 18 \text{ milliards de dollars.} \]
Écart relatif par rapport à l'APD :
\[ \dfrac{48-30}{30}\times 100 = \dfrac{18}{30}\times 100 = 60\,\%. \]
Les transferts des migrants sont donc supérieurs de 60~\% à l'APD.
\item \textbf{IDE après augmentation de 10~\% :}
\[ 40 \times \left(1 + \dfrac{10}{100}\right) = 40 \times 1,1 = 44 \text{ milliards de dollars.} \]
\end{enumerate}
\textbf{Points de vigilance :} pour le pourcentage relatif, le diviseur est la valeur de référence (l'APD, soit 30), pas le total ; pour une hausse de 10~\%, multiplier par 1,1 et non ajouter 10.
\end{corrige}

\begin{corrige}[Exercice 5 — Classer les flux de la mondialisation]
\begin{center}
\begin{tabularx}{\linewidth}{|l|X|}\hline
\rowcolor{popblueL} \textbf{Catégorie de flux} & \textbf{Numéros des phénomènes} \\ \hline
Flux physiques de marchandises & 1 (cacao vers Douala), 6 (pétrole par l'oléoduc Tchad-Cameroun) \\ \hline
Circulation des personnes & 2 (étudiante camerounaise au Maroc) \\ \hline
Circulation de l'information et des services & 4 (achat en ligne sur un site chinois), 8 (formation en ligne d'une université américaine) \\ \hline
Mouvements financiers et IDE & 3 (filiale française de téléphonie à Yaoundé = IDE), 5 (transfert d'argent du migrant), 7 (agence bancaire nigériane à Douala = IDE) \\ \hline
\end{tabularx}
\end{center}
\textbf{Conclusion.} La mondialisation ne se réduit pas au commerce de marchandises : elle combine des flux matériels (cacao, pétrole), humains (étudiants, migrants), immatériels (informations, services en ligne) et financiers (IDE, transferts de la diaspora), qui relient étroitement le Cameroun au reste du monde.

\textbf{Points de vigilance :} le cas 5 est un \emph{mouvement financier} (transfert de la diaspora), pas une circulation de personnes : c'est l'argent qui circule, pas le migrant. Le cas 4 est un flux d'information/service (commerce en ligne), même si le téléphone finit par être livré : le phénomène décrit est l'acte d'achat électronique.
\end{corrige}

\begin{corrige}[Exercice 6 — Lire un tableau statistique : les exportations du Cameroun]
\begin{enumerate}
\item \textbf{Part cumulée des produits agricoles et forestiers :}
\[ 12 + 10 + 6 + 3 + 3 = 34\,\%. \]
(Cacao 12 + bois 10 + coton 6 + café 3 + bananes 3.)
\item \textbf{Calcul des angles du diagramme circulaire} (angle = part $\times$ 3,6) :
\begin{align*}
\text{Pétrole : } & 40 \times 3,6 = 144^\circ &
\text{Cacao : } & 12 \times 3,6 = 43,2^\circ \\
\text{Bois : } & 10 \times 3,6 = 36^\circ &
\text{Coton : } & 6 \times 3,6 = 21,6^\circ \\
\text{Café : } & 3 \times 3,6 = 10,8^\circ &
\text{Bananes : } & 3 \times 3,6 = 10,8^\circ \\
\text{Autres : } & 26 \times 3,6 = 93,6^\circ
\end{align*}
Vérification : $144 + 43,2 + 36 + 21,6 + 10,8 + 10,8 + 93,6 = 360^\circ$.
\begin{popfigure}
\begin{center}
\begin{tikzpicture}[scale=1.6]
\fill[popgold] (0,0) -- (90:1) arc (90:-54:1) -- cycle;
\fill[poporange] (0,0) -- (-54:1) arc (-54:-97.2:1) -- cycle;
\fill[popgreen] (0,0) -- (-97.2:1) arc (-97.2:-133.2:1) -- cycle;
\fill[popblue] (0,0) -- (-133.2:1) arc (-133.2:-154.8:1) -- cycle;
\fill[poppurple] (0,0) -- (-154.8:1) arc (-154.8:-165.6:1) -- cycle;
\fill[poppink] (0,0) -- (-165.6:1) arc (-165.6:-176.4:1) -- cycle;
\fill[popmuted] (0,0) -- (-176.4:1) arc (-176.4:-270:1) -- cycle;
\draw (0,0) circle (1);
\node[font=\scriptsize] at (18:0.62) {\textbf{40~\%}};
\node[font=\scriptsize, rotate=-75] at (-75:0.68) {12~\%};
\node[font=\scriptsize, rotate=-115] at (-115:0.68) {10~\%};
\node[font=\tiny, rotate=-144] at (-144:0.72) {6~\%};
\node[font=\tiny, rotate=-160] at (-160:0.74) {3~\%};
\node[font=\tiny, rotate=-171] at (-171:0.74) {3~\%};
\node[font=\scriptsize, rotate=-223] at (-223:0.68) {26~\%};
\end{tikzpicture}
\end{center}
\legende{Diagramme circulaire de la structure des exportations camerounaises en valeur (ordres de grandeur) : pétrole brut 40~\% (or), cacao 12~\% (orange), bois 10~\% (vert), coton 6~\% (bleu), café 3~\% (violet), bananes 3~\% (rose), autres 26~\% (gris).}
\end{popfigure}
\item \textbf{Deux caractéristiques du commerce extérieur camerounais.} Premièrement, les exportations sont dominées par les produits primaires : le pétrole brut à lui seul représente environ 40~\% des recettes, et avec les produits agricoles et forestiers (34~\%), les produits bruts ou peu transformés forment près des trois quarts des exportations. Deuxièmement, cette structure révèle une forte dépendance à quelques produits de base dont les prix fluctuent sur les marchés mondiaux, ce qui fragilise les recettes extérieures du pays.
\end{enumerate}
\textbf{Points de vigilance :} pour un diagramme circulaire, multiplier chaque part par 3,6 (et non par 360) ; toujours vérifier que la somme des angles fait 360° ; le bois est ici compté avec les produits agricoles/forestiers, comme le demande l'énoncé.
\end{corrige}

\begin{corrige}[Exercice 7 — Construire un diagramme en bâtons : les migrants dans le monde]
\begin{enumerate}
\item \textbf{Diagramme en bâtons} (échelle : 1~cm pour 25 millions de migrants ; hauteurs : 1990 : 6,1~cm ; 2000 : 6,9~cm ; 2010 : 8,8~cm ; 2020 : 11,2~cm) :
\begin{popfigure}
\begin{center}
\begin{tikzpicture}
\begin{axis}[
  ybar, bar width=0.9cm,
  width=0.85\linewidth, height=7cm,
  ymin=0, ymax=320,
  ylabel={Migrants internationaux (millions)},
  symbolic x coords={1990,2000,2010,2020},
  xtick=data,
  nodes near coords,
  every node near coord/.append style={font=\small},
  ymajorgrids=true, grid style={popline!40},
  axis lines=left,
]
\addplot[fill=popblue, draw=popdark] coordinates {(1990,153) (2000,173) (2010,221) (2020,281)};
\end{axis}
\end{tikzpicture}
\end{center}
\legende{Évolution du nombre de migrants internationaux dans le monde, 1990-2020 (en millions, ordres de grandeur ; source : ONU).}
\end{popfigure}
\item \textbf{Taux de variation entre 1990 et 2020 :}
\[ \dfrac{281 - 153}{153}\times 100 = \dfrac{128}{153}\times 100 \approx 83,7\,\% \approx 84\,\%. \]
Le nombre de migrants internationaux a augmenté d'environ 84~\% en trente ans.
\item \textbf{Commentaire.} Le diagramme montre une augmentation continue et forte du nombre de migrants internationaux : de 153 millions en 1990, il passe à 281 millions en 2020, soit une hausse d'environ 84~\% en trente ans. La croissance s'accélère après 2000 (+48 millions entre 2000 et 2010, puis +60 millions entre 2010 et 2020). Cette évolution traduit l'intensification de la circulation des personnes dans la mondialisation : recherche d'emplois, études à l'étranger, mais aussi déplacements forcés liés aux conflits. Les migrations relient désormais toutes les régions du monde, du Sud vers le Nord mais aussi entre pays du Sud.
\end{enumerate}
\textbf{Points de vigilance :} respecter l'échelle annoncée et la reporter sur la copie ; dans le taux de variation, diviser par la valeur \emph{initiale} (153) ; un commentaire doit décrire (chiffres à l'appui) puis interpréter.
\end{corrige}

\begin{corrige}[Exercice 8 — Schématiser une filière mondialisée : le cacao camerounais]
\begin{popfigure}
\begin{center}
\begin{tikzpicture}[font=\small, >=Stealth]
% Zones de production
\node[draw, fill=popgreenL, rounded corners, minimum width=4.2cm, minimum height=1.6cm, align=center] (prod) at (0,0) {\textbf{Zones de production}\\ Centre, Sud, Sud-Ouest\\ (plantations villageoises)};
% Port de Douala
\node[draw, fill=popblueL, rounded corners, minimum width=3.2cm, minimum height=1.2cm, align=center] (port) at (6.5,0) {\textbf{Port de Douala}\\ (embarquement)};
% Europe
\node[draw, fill=poporangeL, rounded corners, minimum width=4cm, minimum height=1.6cm, align=center] (ue) at (12.5,1.8) {\textbf{Union européenne}\\ Pays-Bas (transformation :\\ chocolat)};
% Monde / réimportation
\node[draw, fill=poppinkL, rounded corners, minimum width=3.6cm, minimum height=1.4cm, align=center] (monde) at (12.5,-1.8) {\textbf{Marché mondial}\\ (produits finis)};
% Flèches
\draw[->, very thick, popgreen!70!black] (prod) -- (port) node[midway, above, font=\scriptsize]{piste, rail, route};
\draw[->, very thick, popblue] (port) -- (ue) node[midway, above, sloped, font=\scriptsize]{fèves brutes (voie maritime)};
\draw[->, very thick, poporange] (ue) -- (monde) node[midway, right, font=\scriptsize]{chocolat};
\draw[->, thick, poppink, dashed] (monde.south) .. controls (6.5,-3.2) .. (prod.south) node[midway, below, font=\scriptsize]{réimportation de chocolat au Cameroun};
\end{tikzpicture}
\end{center}
\legende{Schéma de la filière cacao camerounaise : 1. production (Centre, Sud, Sud-Ouest) ; 2. acheminement vers le port de Douala ; 3. exportation des fèves brutes vers l'Union européenne (Pays-Bas) ; 4. transformation en chocolat et vente mondiale ; 5. réimportation des produits finis au Cameroun.}
\end{popfigure}
\textbf{Phrase de conclusion.} Cette filière illustre l'insertion du Cameroun dans la mondialisation comme fournisseur de matière première brute : le pays produit et exporte les fèves, mais la transformation (et l'essentiel de la valeur ajoutée) se réalise en Europe, ce qui montre une intégration réelle mais subordonnée aux flux commerciaux mondiaux.

\textbf{Points de vigilance :} un schéma de filière exige des flèches \emph{orientées} (sens du flux), des lieux nommés et une légende organisée ; ne pas oublier la boucle de réimportation, qui montre le paradoxe d'un pays producteur important du produit fini.
\end{corrige}

\begin{corrige}[Exercice 9 — Type Bac : étude de documents (les flux de marchandises et la place du Cameroun)]
\textbf{1. Répartition mondiale des grands flux maritimes (4 points ; 1 point par idée).} Les flux maritimes mondiaux sont concentrés dans l'hémisphère Nord, autour de trois axes dominants reliant les pôles de la Triade : l'axe Atlantique Nord (Amérique du Nord--Europe), l'axe Asie--Europe par Suez et l'axe transpacifique (Asie de l'Est--Amérique du Nord). Le trafic est canalisé par des détroits et canaux stratégiques (Ormuz, Malacca, Gibraltar, Panama, Suez). À l'inverse, les côtes d'Afrique subsaharienne ne connaissent que des flux secondaires : moins de 5~\% du trafic portuaire mondial.

\textbf{2. Destinations des exportations camerounaises (4 points).} Les exportations camerounaises sont fortement concentrées : l'Union européenne absorbe environ 45~\% des ventes (Pays-Bas, Espagne, Italie, France), ce qui traduit la persistance des liens historiques et coloniaux. La Chine (15~\%) et l'Inde (6~\%) confirment la montée en puissance des partenaires asiatiques. Les échanges régionaux (CEMAC et Afrique : 12~\%) restent faibles, signe d'une intégration régionale limitée.

\textbf{3. La place du Cameroun dans le commerce mondial (8 points).} Le Cameroun occupe une place marginale et subordonnée. Sa part dans le commerce mondial est inférieure à 0,05~\% (doc. 3) et ses ports ne drainent que des flux secondaires (doc. 1). Ses échanges sont de type \emph{colonial} ou \emph{néocolonial} : il exporte des produits primaires bruts (pétrole, cacao, bois, coton) dont les prix sont fixés sur les marchés mondiaux, et importe des produits manufacturés à forte valeur ajoutée (machines, véhicules, médicaments) ainsi que des carburants raffinés — alors même qu'il exporte du pétrole brut. Cette structure engendre une dépendance commerciale (concentration sur l'UE), une dépendance aux cours mondiaux des matières premières et une dépendance logistique (un seul grand port, Douala). Le Cameroun est donc intégré à la mondialisation, mais en périphérie du système.

\textbf{4. Deux solutions (4 points ; 2 points par solution argumentée).}
\begin{itemize}
\item \textbf{Transformer localement les matières premières} (usines de chocolat, scieries, raffinage, filature de coton) pour exporter des produits à plus forte valeur ajoutée et réduire la dépendance aux cours des produits bruts.
\item \textbf{Diversifier les partenaires et renforcer l'intégration régionale} : développer les échanges avec la CEMAC et la Zone de libre-échange continentale africaine (ZLECAf), moderniser les infrastructures (port de Kribi, routes, rail) pour abaisser les coûts d'exportation.
\end{itemize}
\textbf{Points de vigilance :} citer les documents (« d'après le document 2... ») ; ne pas recopier les documents mais les exploiter ; chaque solution doit être explicitée, pas simplement nommée.
\end{corrige}

\begin{corrige}[Exercice 10 — Type Bac : croquis « Les principaux flux de la mondialisation »]
\textbf{Éléments attendus sur le croquis} (fond de planisphère simplifié) :
\begin{itemize}
\item \textbf{Triade localisée et nommée :} Amérique du Nord, Union européenne, Asie de l'Est (Chine, Japon, Corée du Sud) ; le Cameroun localisé en Afrique centrale.
\item \textbf{Flux de marchandises :} flèches pleines épaisses sur les trois axes maritimes majeurs (Atlantique Nord, Asie--Europe par Suez, transpacifique) ; détroits stratégiques nommés (Suez, Malacca, Panama, Ormuz).
\item \textbf{Migrations :} flèches d'un second symbole : Amérique latine $\to$ Amérique du Nord ; Afrique du Nord et subsaharienne $\to$ Europe ; Asie du Sud $\to$ golfe Persique.
\item \textbf{IDE :} flèches d'un troisième symbole partant des pôles de la Triade vers la Chine, l'Asie du Sud-Est, l'Amérique latine et (plus fines) vers l'Afrique.
\item \textbf{Flux d'information :} traits fins ou pointillés figurant les câbles sous-marins reliant les continents.
\item \textbf{Légende organisée} hors de la carte, titre en haut, flèche du nord.
\end{itemize}
\begin{popfigure}
\includegraphics[width=\linewidth,height=9.5cm,keepaspectratio]{N12/shipping_routes.jpg}
\legende{Fond de carte pour le croquis « Les principaux flux de la mondialisation » : la densité du trafic maritime mondial (plus le bleu est foncé, plus le trafic est dense ; carte sans noms de lieux). Le croquis attendu y ajoute : la Triade localisée, les flux de marchandises sur les trois axes maritimes majeurs (Atlantique Nord, Asie--Europe par Suez, transpacifique), les détroits stratégiques, les migrations, les IDE et les câbles d'information, le Cameroun étant localisé en Afrique centrale. On lit sur le fond que les trois axes concentrent l'essentiel du trafic. \\ Source : Wikimedia Commons, « Shipping routes » (densité du trafic maritime commercial), T. Hengl, CC BY-SA 3.0.}
\end{popfigure}
\textbf{Analyse rédigée (attendue : 10 à 15 lignes).} Le croquis montre que la mondialisation est un système hiérarchisé. Au sommet, les trois pôles de la Triade (Amérique du Nord, Union européenne, Asie de l'Est) concentrent et commandent les flux : les grands axes maritimes de marchandises relient ces pôles entre eux, les IDE en partent vers les pays émergents, et les câbles d'information y aboutissent. Autour d'eux, des espaces intégrés (pays émergents d'Asie, d'Amérique latine) captent une part croissante des flux. En contrepartie, les migrations internationales s'écoulent des espaces pauvres vers les pôles riches. Enfin, l'Afrique subsaharienne, dont le Cameroun, apparaît marginalisée : traversée par des flux faibles (IDE tournés vers le pétrole et les mines, trafic portuaire secondaire), elle fournit des matières premières sans maîtriser les circuits. Le Cameroun, localisé en Afrique centrale, illustre cette position périphérique : relié au monde par le port de Douala, il reste à l'écart des grands axes de la mondialisation.

\textbf{Barème indicatif :} exactitude des localisations (5 pts), diversité et lisibilité des flux (5 pts), titre et légende organisée (5 pts), analyse (5 pts).
\textbf{Points de vigilance :} la légende doit être placée \emph{hors} de la carte ; chaque type de flux doit avoir un symbole distinct et constant ; ne pas oublier le titre ni l'orientation ; l'analyse doit mobiliser le vocabulaire du cours (pôles, hiérarchie, marginalisation).
\end{corrige}

\begin{corrige}[Exercice 11 — Type Bac : sujet de synthèse (mondialisation, commerce équitable et pays pauvres)]
\textbf{Introduction (3 points).} La mondialisation est le processus d'intensification des échanges de marchandises, de personnes, de capitaux et d'informations à l'échelle de la planète, qui relie les territoires en un système unique mais hiérarchisé. Pour les pays pauvres comme le Cameroun, elle suscite un débat : est-elle une chance de développement ou une machine à creuser les inégalités ? Le commerce équitable, qui garantit un prix plancher et une prime aux producteurs, prétend corriger ces déséquilibres. Nous montrerons d'abord que la mondialisation offre de réelles opportunités aux pays pauvres, puis qu'elle engendre de graves problèmes, avant d'évaluer la capacité — réelle mais limitée — du commerce équitable à corriger les déséquilibres mondiaux.

\textbf{I. Les opportunités de la mondialisation pour les pays pauvres (4 points).} La mondialisation ouvre aux pays pauvres l'accès aux marchés mondiaux : le Cameroun exporte cacao, café, coton et pétrole vers l'Europe et l'Asie, ce qui procure des devises. Elle attire les IDE : l'oléoduc Tchad-Cameroun, le port en eau profonde de Kribi ou les filiales de téléphonie mobile créent emplois et infrastructures. Les transferts de la diaspora camerounaise (plusieurs centaines de millions de dollars par an) soutiennent la consommation, l'éducation et l'investissement des familles. Enfin, la diffusion de l'information et des services (Internet, mobile money, formations en ligne, commerce en ligne) permet aux jeunes d'accéder au savoir et à de nouveaux débouchés économiques.

\textbf{II. Les problèmes liés à la mondialisation (4 points).} Mais ces opportunités s'accompagnent de déséquilibres. Les échanges restent inégaux : le Cameroun vend des produits primaires bruts bon marché et achète des produits manufacturés coûteux, ce qui dégrade les termes de l'échange. Le pays reste marginalisé : moins de 0,05~\% du commerce mondial, des flux d'IDE concentrés sur le pétrole et les mines. La mondialisation favorise aussi des fléaux : l'immigration clandestine (jeunes qui risquent leur vie en Méditerranée), la cybercriminalité (arnaques et fraudes en ligne), et les conflits armés alimentés par le trafic de ressources (diamants, bois, pétrole). La fracture numérique exclut en outre les campagnes et les plus pauvres des bénéfices de l'information.

\textbf{III. Le commerce équitable peut-il corriger ces déséquilibres ? (4 points).} Le commerce équitable apporte des corrections concrètes : il garantit un prix plancher aux producteurs, verse une prime de développement à la coopérative, impose des normes sociales (refus du travail des enfants) et environnementales, et instaure des relations commerciales durables. Au Cameroun, des coopératives de cacao et de café de l'Ouest et du Centre, certifiées, exportent vers l'Europe et améliorent le revenu de leurs membres (document A). Mais ses limites sont fortes : il concerne moins de 2~\% du commerce mondial, les coûts de certification sont élevés pour les petits producteurs, et il ne règle ni l'immigration clandestine, ni la cybercriminalité, ni la marginalisation financière (document B). Il corrige donc localement les inégalités sans transformer les règles du système mondial.

\textbf{Conclusion (2 points).} La mondialisation est ambivalente pour les pays pauvres : source d'opportunités (marchés, IDE, transferts, information), elle aggrave aussi les déséquilibres et engendre de nouveaux problèmes. Le commerce équitable est un correctif utile mais insuffisant, car marginal. Une véritable réduction des inégalités suppose aussi la transformation locale des matières premières, l'intégration régionale africaine (ZLECAf) et une régulation internationale plus juste des échanges.

\textbf{Barème indicatif :} introduction (3 pts), développement en trois parties (12 pts), exactitude des exemples camerounais (3 pts), conclusion (2 pts).
\textbf{Points de vigilance :} annoncer clairement le plan en introduction ; mobiliser les deux documents d'appui et les citer ; illustrer chaque partie par au moins un exemple camerounais précis ; éviter le catalogue de notions sans articulation ; la conclusion doit être nuancée et comporter une ouverture.
\end{corrige}

\newpage

\coursec{Fiche bilan}

\courssub{Fiche bilan : Le fonctionnement de la mondialisation}

\begin{popfigure}
\centering
\begin{tikzpicture}[mindmap, grow cyclic,
  every node/.style={concept, text=white, font=\small\bfseries},
  root concept/.append style={concept color=popink, font=\bfseries, minimum size=2.6cm},
  level 1 concept/.append style={level distance=3.6cm, sibling angle=60, font=\scriptsize\bfseries},
  level 2 concept/.append style={level distance=2.2cm, font=\tiny},
  scale=0.92, transform shape]
\node[root concept]{Le fonctionnement de la mondialisation}
  child[concept color=popblue]{ node{Flux physiques de marchandises}
    child[concept color=popblue!70]{ node{produits agricoles, minerais, conteneurs} }}
  child[concept color=popgreen]{ node{Circulation des personnes}
    child[concept color=popgreen!70]{ node{migrations, tourisme, diaspora} }}
  child[concept color=poporange]{ node{Information et services}
    child[concept color=poporange!70]{ node{Internet, câbles, commerce en ligne} }}
  child[concept color=poppurple]{ node{Mouvements financiers}
    child[concept color=poppurple!70]{ node{IDE, transferts diaspora, aides} }}
  child[concept color=popgold]{ node{Commerce équitable (Dossier 3)}
    child[concept color=popgold!70]{ node{normes, prix juste, critiques} }}
  child[concept color=poppink]{ node{Problèmes (Dossier 4)}
    child[concept color=poppink!70]{ node{immigration clandestine, cybercriminalité, conflits, marginalisation} }};
\end{tikzpicture}
\legende{Carte mentale de la leçon : les six branches à maîtriser (4 flux + 2 dossiers).}
\end{popfigure}

\begin{bilan}
\textbf{Définitions clés à connaître par cœur}
\begin{itemize}
  \item \cle{Mondialisation} : processus d'intensification des échanges (biens, personnes, capitaux, informations) qui met les territoires du monde en interdépendance.
  \item \cle{Flux} : mouvement de marchandises, de personnes, de capitaux ou d'informations entre deux ou plusieurs lieux.
  \item \cle{IDE} (Investissements Directs à l'Étranger) : investissement d'une entreprise (souvent une firme transnationale) dans un pays étranger pour y créer ou contrôler une activité productive.
  \item \cle{Firme transnationale (FTN)} : entreprise implantée dans plusieurs pays, qui organise sa production à l'échelle mondiale.
  \item \cle{Transferts de la diaspora} : envois d'argent des migrants vers leur pays d'origine (pour le Cameroun : plusieurs centaines de milliards de FCFA par an, ordre de grandeur).
  \item \cle{Commerce équitable} : échange fondé sur un prix rémunérateur garanti au producteur, des relations durables et le respect de normes sociales et environnementales.
  \item \cle{Commerce en ligne (e-commerce)} : achat et vente de biens et services via Internet.
  \item \cle{Cybercriminalité} : ensemble des délits communs via les réseaux numériques (arnaques, piratage, vol de données).
\end{itemize}

\textbf{Repères chiffrés essentiels (ordres de grandeur)}

\begin{center}
\begin{tabularx}{\linewidth}{|l|X|}
\hline
\rowcolor{popblueL} \textbf{Indicateur} & \textbf{Ordre de grandeur} \\
\hline
Part de la Triade (Amérique du Nord, UE, Asie de l'Est) dans les échanges mondiaux & environ les deux tiers (dominance nette) \\
\hline
Part du transport maritime dans le commerce mondial (volume) & environ 80 à 90 \% \\
\hline
Migrants internationaux dans le monde & environ 280 millions (environ 3,5 \% de l'humanité) \\
\hline
Transferts de la diaspora vers les pays en développement & plusieurs centaines de milliards de dollars par an, supérieurs à l'aide publique au développement \\
\hline
Position du Cameroun dans la mondialisation & périphérie : exporte des produits bruts (pétrole, cacao, bois, coton), importe des produits manufacturés \\
\hline
\end{tabularx}
\end{center}

\textbf{Méthodes express}
\begin{itemize}
  \item \emph{Croquis de flux} : fond de carte simplifié $\rightarrow$ flèches d'épaisseur proportionnelle à l'importance du flux $\rightarrow$ légende numérotée + titre + orientation.
  \item \emph{Commentaire de carte} : observer (localiser) $\rightarrow$ décrire (directions, hiérarchie des flux) $\rightarrow$ interpréter (acteurs, inégalités Nord/Sud).
  \item \emph{Étude de cas} : toujours illustrer par un exemple camerounais (port de Douala, cacao, transferts de la diaspora, Jumia, MTN Mobile Money).
\end{itemize}
\end{bilan}

\begin{aretenir}[Checklist avant le Bac]
\begin{itemize}
  \item $\square$ Je sais définir la mondialisation, un flux, un IDE, une FTN et le commerce équitable.
  \item $\square$ Je connais les trois grands produits des flux physiques : agricoles, minerais (et hydrocarbures), manufacturés.
  \item $\square$ Je sais que le transport maritime (conteneurisation) domine le commerce mondial.
  \item $\square$ Je distingue les migrations volontaires, forcées et le tourisme international.
  \item $\square$ Je sais expliquer le rôle d'Internet et des câbles sous-marins dans la circulation de l'information et le commerce en ligne.
  \item $\square$ Je distingue IDE, aides bilatérales/multilatérales et transferts de la diaspora.
  \item $\square$ Je connais les normes du commerce équitable (prix garanti, prime de développement, conditions de travail) et ses critiques.
  \item $\square$ Je peux citer quatre problèmes de la mondialisation : immigration clandestine, cybercriminalité, conflits armés, marginalisation des pays pauvres.
  \item $\square$ Je sais construire et légender un croquis de flux mondiaux avec flèches proportionnelles.
  \item $\square$ Je sais illustrer chaque notion par un exemple camerounais précis (Douala, cacao, diaspora, mobile money).
  \item $\square$ Je rédige un commentaire de document en trois étapes : observer, décrire, interpréter.
  \item $\square$ Je vérifie mes chiffres : j'annonce les ordres de grandeur quand la valeur exacte est incertaine.
\end{itemize}
\end{aretenir}

\findecours

\end{document}
